\documentclass[11pt,a4paper]{article}

\usepackage[utf8]{inputenc}
\usepackage[T1]{fontenc}
\usepackage{lmodern}
\usepackage[french]{babel}
\usepackage[margin=2.1cm]{geometry}
\usepackage{microtype}
\usepackage{amsmath,amssymb}
\usepackage{eurosym}
\usepackage{booktabs}
\usepackage{array}
\usepackage{tabularx}
\usepackage[table]{xcolor}
\usepackage{graphicx}
\usepackage{caption}
\usepackage{enumitem}
\usepackage[most]{tcolorbox}
\usepackage{tikz}
\usetikzlibrary{arrows.meta,positioning,calc,fit,backgrounds,decorations.pathreplacing,shapes.geometric}
\usepackage[hidelinks]{hyperref}

\definecolor{navy}{HTML}{1F3864}
\definecolor{bluemid}{HTML}{2E5496}
\definecolor{lightblue}{HTML}{D6E0F0}
\definecolor{accent}{HTML}{C0491F}
\definecolor{vert}{HTML}{3D8361}

\captionsetup{font=small,labelfont={bf,color=navy},labelsep=period}
\graphicspath{{../figures/}}

\newtcolorbox{phrasebox}{colback=vert!7,colframe=vert!70!black,boxrule=0.7pt,
  arc=2pt,left=7pt,right=7pt,top=5pt,bottom=5pt,before skip=6pt,after skip=8pt}
\newtcolorbox{cle}[1][]{colback=lightblue!40,colframe=navy,boxrule=0.6pt,
  arc=2pt,left=7pt,right=7pt,top=5pt,bottom=5pt,#1}

\providecommand{\Prob}{\mathbb{P}}
\providecommand{\E}{\mathbb{E}}
\providecommand{\R}{\mathbb{R}}

\newcommand{\lequoi}[1]{\smallskip\noindent\textcolor{navy}{\textbf{Le problème.}}\ #1\par}
\newcommand{\lidee}[1]{\smallskip\noindent\textcolor{navy}{\textbf{L'idée.}}\ #1\par}
\newcommand{\laformule}[1]{\smallskip\noindent\textcolor{navy}{\textbf{La formule.}}\ #1\par}
\newcommand{\cheznous}[1]{\smallskip\noindent\textcolor{navy}{\textbf{Chez nous.}}\ #1\par}
\newcommand{\lepiege}[1]{\smallskip\noindent\textcolor{accent}{\textbf{Le piège.}}\ #1\par}
\newcommand{\laphrase}[1]{\begin{phrasebox}\textbf{À dire à voix haute.} \emph{#1}\end{phrasebox}}

% Questions du jury : la question, puis la reponse a tenir.
\newtcolorbox{question}[1]{colback=accent!5,colframe=accent!75!black,boxrule=0.6pt,
  arc=2pt,left=7pt,right=7pt,top=5pt,bottom=5pt,before skip=7pt,after skip=4pt,
  fonttitle=\bfseries\color{white},title={#1}}
\newcommand{\rep}[1]{\smallskip\noindent\textcolor{navy}{\textbf{La réponse.}}\ #1\par}
\newcommand{\appui}[1]{\smallskip\noindent\small\textcolor{navy}{\textbf{Appui~:}} #1\par\normalsize}

% Partie X, chapitre par chapitre : les reperes de pages, le chapitre en une phrase, et la
% fiche \og par coeur\fg{} (tableau a deux colonnes dans un cadre, titre optionnel).
\newcommand{\reperes}[1]{\noindent{\small\color{navy!75}#1}\par\smallskip}
\newcommand{\enbref}[1]{\noindent\textcolor{navy}{\textbf{En une phrase.}}\ #1\par}
\newtcolorbox{parcoeur}[1][À savoir par c\oe ur]{colback=navy!3,colframe=navy!85!black,
  boxrule=0.6pt,arc=2pt,left=6pt,right=6pt,top=4pt,bottom=4pt,before skip=7pt,after skip=5pt,
  fonttitle=\bfseries\small\color{white},title={#1}}
\newenvironment{fichetab}{\small\renewcommand{\arraystretch}{1.1}%
  \tabularx{\linewidth}{@{}>{\raggedright\arraybackslash\bfseries}p{3.2cm}
  >{\raggedright\arraybackslash}X@{}}}{\endtabularx}

% Styles des schemas
\tikzset{
  brique/.style={draw=navy,fill=lightblue!45,rounded corners=2pt,inner sep=5pt,
                 align=center,font=\footnotesize,minimum height=9mm},
  pilier/.style={draw=navy,fill=white,circle,inner sep=1.5pt,minimum size=8mm,
                 font=\footnotesize\bfseries,text=navy},
  fleche/.style={-{Stealth[length=2mm]},draw=navy,line width=0.5pt},
  flechef/.style={-{Stealth[length=2.6mm]},draw=accent,line width=1.1pt},
  flechep/.style={-{Stealth[length=1.8mm]},draw=navy!45,line width=0.4pt},
  etiq/.style={font=\scriptsize\itshape,text=navy!75},
}

\setcounter{secnumdepth}{2}
\emergencystretch=1.2em   % noms propres longs (Kolmogorov-Smirnov, Pickands-Balkema)

\title{\vspace{-1.3cm}\color{navy}\bfseries Le mémoire, concept par concept\\[3pt]
  \large Ce que chaque outil résout, comment il marche,
  et la phrase qui l'explique en trente secondes}
\author{Kélian Kaddouri}
\date{30 septembre 2026}

\begin{document}
\maketitle
\vspace{-0.8cm}

\begin{cle}
\textbf{À quoi sert ce document.} Le mémoire empile une vingtaine d'outils, chacun choisi
pour une raison précise, et il est facile de les manipuler correctement sans savoir les
\emph{défendre}. Chaque section suit le même gabarit~: le problème que l'outil résout,
l'idée en langage courant, la formule, ce que ça donne dans notre cas, la phrase à dire à
voix haute quand on vous interroge, et le piège à éviter.

\textbf{La phrase encadrée en vert est le livrable de chaque section.} Si vous ne retenez
qu'elle, vous tenez l'oral.

\textbf{La partie~X reprend le mémoire chapitre par chapitre.} Pour chacun : le chapitre en une
phrase, un encadré de ce qu'il faut savoir \emph{par c\oe ur}, puis les questions que le jury peut
poser, avec la réponse à tenir et la page où s'appuyer. C'est la partie à relire la veille de la
soutenance.
\end{cle}

\smallskip
\noindent\textcolor{accent}{\textbf{Mise à jour du 30 septembre 2026, jour du dépôt.}}
L'ancienne partie~X, qui rangeait les questions par thème, est remplacée par une partie
\textbf{chapitre par chapitre}, recalée sur le PDF déposé ($199$~pages). Les anciennes questions y
sont reprises et corrigées, et leur nombre passe d'une quarantaine à près de cent cinquante. Une
relecture intégrale du mémoire contre ce cours a corrigé des erreurs, dont
\textcolor{accent}{trois qui auraient coûté cher à l'oral} : le capital est une
$\mathrm{VaR}_{99{,}5\,\%}$ \emph{non centrée}, et non \og VaR moins perte attendue\fg{} ; la phrase de
clôture annonçait que l'ordre des piliers ne bouge jamais, alors que seul le \emph{groupe de tête},
P1 et P4, résiste à l'ignorance de la direction ; et les $66\,\%$ de valeur d'information étaient
présentés comme un acquis, alors que c'est la valeur d'une information que le mémoire ne revendique
que sur deux paires. Également corrigés : le seuil des excès dans la section qui le présente
($20{,}03$~M\euro, plus $4{,}176$), la surdispersion ($1{,}19$ mesuré, $9{,}20$ posé), la fréquence
d'entité ($0{,}092$) et le facteur $3\,719$, le processus de Hawkes (testé sur données datées, puis
écarté par équivalence observationnelle), les liens faibles du corpus (les quatre sorties de P2),
le facteur du bootstrap ($2{,}5$), les tests d'adéquation, les trois étages d'incertitude lus sur
une même échelle, les postures de VaR, la co-occurrence (qui ne voit $S$ qu'au premier ordre),
l'élicitation (diffusée et restée sans réponse), et la table des échelles. Les nombres absents du
mémoire et du matériel complémentaire ont été retirés.

\bigskip
\noindent\textbf{Le schéma d'ensemble.} Cinq étages, de la donnée au capital. Chaque partie de ce
document en traite un.

\begin{center}
\begin{tikzpicture}[node distance=4.5mm]
\node[brique] (d) {\textbf{Données}\\OpRisk, PRC\\Hackmageddon};
\node[brique,right=of d] (f) {\textbf{Fréquence}\\binomiale\\négative};
\node[brique,right=of f] (s) {\textbf{Sévérité}\\corps\,+\,queue GPD\\$\xi=0{,}5954$};
\node[brique,right=of s] (c) {\textbf{Cascade}\\dirigée $W$\\$\rho(W)<1$};
\node[brique,right=of c] (m) {\textbf{Conformité}\\multi-états\\(Markov)};
\node[brique,right=of m,fill=vert!15,draw=vert!70!black] (k)
  {\textbf{Capital}\\$\mathrm{VaR}_{99{,}5\%}(L)$};
\draw[fleche] (d)--(f); \draw[fleche] (f)--(s); \draw[fleche] (s)--(c);
\draw[fleche] (c)--(m); \draw[fleche] (m)--(k);
\draw[flechef] (d.south) to[out=-60,in=-120] node[etiq,below,pos=0.5]
  {la direction de $W$ n'est PAS identifiable ici} (c.south);
\end{tikzpicture}
\end{center}

\noindent La flèche rouge est le sujet du mémoire~: tout le reste se calibre sur la donnée,
\emph{sauf} la direction de la contagion. C'est de là que vient l'identification partielle,
et c'est la question que le jury posera en premier.

\tableofcontents
\newpage

% =====================================================================
\part*{I. D'où viennent les nombres}
\addcontentsline{toc}{part}{I. D'où viennent les nombres}

\section{Les trois sources, et ce que chacune peut dire}

\lequoi{Pour chiffrer un capital, il faut trois choses~: combien d'incidents par an, combien
coûte un incident, et comment les incidents s'enchaînent. Aucune source publique ne porte
les trois.}

\lidee{On ne cherche pas \emph{la} bonne base, on répartit les questions entre les sources
selon ce que chacune peut porter, et on déclare ce qui manque.}

\cheznous{Trois bases pour les marges, un corpus pour l'ordre, deux sources de contrôle.
\begin{itemize}[itemsep=1pt,topsep=2pt]
\item \textbf{SAS OpRisk Global Data}~: pertes opérationnelles en dollars, catégories de Bâle ;
$582$ pertes cyber $\times$ finance sur 2000--2026 pour la sévérité, $4\,329$ transitions
annuelles pour le test de direction. Elle porte le \textbf{niveau} de sévérité et la fréquence
d'entrée. Son pas est annuel, ce qui efface l'ordre des événements.
\item \textbf{Privacy Rights Clearinghouse (PRC)}~: $15\,053$ incidents retenus, exprimés en
\emph{nombre d'enregistrements compromis}, pas en euros. Elle sert d'\textbf{axe de sensibilité}
après conversion, et de clé de répartition ($\lambda_{ref}=341$).
\item \textbf{Hackmageddon}~: $1\,041$ incidents, $840$ à vecteur identifié. Elle donne les
\textbf{parts par vecteur d'attaque}, jamais un niveau ; c'est une citation, non recalculable.
\item \textbf{Dix rapports post-incident} ($22$ transitions)~: peu de volume, mais une enquête y
reconstitue la \emph{séquence} des défaillances. C'est la seule source qui porte
l'\textbf{ordre}.
\item \textbf{VERIS} ($10\,591$ incidents) montre que le bon schéma existe sans être rempli, et
\textbf{quatre SFCR} donnent l'échelle d'entité.
\end{itemize}}

\laphrase{Aucune base ne dit tout. OpRisk donne le niveau mais pas l'ordre, PRC donne du volume
mais pas d'euros, Hackmageddon la structure par vecteur, et les rapports d'enquête donnent l'ordre
mais pas le volume. Je les utilise pour ce que chacune peut porter, et je dis explicitement ce qui
reste absent.}

\lepiege{Ne jamais laisser croire qu'on a \og une base de données DORA\fg{}. Elle n'existe pas.
Le règlement met en place le reporting qui la produira, et c'est justement l'objet de la
recommandation finale du mémoire.}

\section{La conversion Jacobs}

\lequoi{PRC compte des enregistrements compromis, pas des euros. Il faut passer de l'un à
l'autre.}

\lidee{Une régression publiée établit, sur un échantillon d'incidents dont on connaît à la
fois le nombre d'enregistrements et le coût, une relation log-log entre les deux.}

\laformule{$\ln L = 7{,}68 + 0{,}76\,\ln X + \varepsilon$, où $X$ est le nombre
d'enregistrements et $L$ la perte en dollars. Le résidu a un écart-type de $0{,}523$ en logarithme :
à volume donné, le coût réel varie d'un facteur $2{,}4$ autour de la droite, à $90\,\%$.}

\cheznous{L'exposant $0{,}76 < 1$ dit que le coût croît \emph{moins vite} que le nombre
d'enregistrements~: un incident dix fois plus gros ne coûte pas dix fois plus cher, mais
environ six fois. Il y a un coût fixe important dans tout incident.}

\laphrase{Jacobs n'est pas un lien fréquence-sévérité, c'est une conversion de sévérité
seule~: elle transforme un nombre d'enregistrements en montant. Et mon capital principal
n'en dépend pas, il est calculé sur OpRisk, qui est déjà en dollars.}

\lepiege{Le modèle traitait d'abord cette conversion comme \emph{exacte}, en ignorant le résidu.
En le réinjectant enregistrement par enregistrement, le niveau du surcoût PRC bouge de $7\,\%$ mais
l'intervalle ne s'élargit presque pas, le bruit se moyennant dans l'ajustement GPD ($\xi$ passe de
$1{,}03$ à $1{,}07$, matériel complémentaire). Et ne jamais lire $\xi_{\text{PRC}}=1{,}03$ comme un
risque plus lourd : une puissance multiplie l'indice de queue par son exposant, donc cet indice
décrit la queue des \emph{enregistrements} transportée en euros.}

% =====================================================================
\section{Le seuil de collecte, et les trois fréquences}

\lequoi{Trois nombres du mémoire portent le nom de \og{}fréquence\fg{} et vont de $0{,}092$ à
$341$ par an. Confondre deux d'entre eux fait dire n'importe quoi.}

\lidee{Ils ne comptent pas la même chose, parce qu'ils ne sont pas lus au même endroit. Une
base ne retient que les pertes au-dessus d'un \textbf{seuil de collecte} : le taux qu'on y lit
est celui des incidents \emph{matériels}, générateurs d'une perte enregistrée, et non celui de
tous les incidents au sens du reporting réglementaire.}

\cheznous{\;\\[-6pt]
\begin{center}\small
\begin{tabular}{@{}l l l@{}}
\toprule
$\lambda_{ref}=341$/an & notifications PRC, organisations financières & \emph{clé de répartition} par vecteur \\
$\lambda=21{,}6$/an & incidents TIC matériels, secteur & \textbf{l'entrée du moteur} \\
$\lambda_{\text{entité}}=0{,}092$/an & les mêmes, lus à la taille de l'entité & l'échelle d'une entité \\
\bottomrule
\end{tabular}
\end{center}
Les seaux du panel valent $0{,}10$ (toutes firmes) et $0{,}21$ (firmes à au moins dix événements),
mais ils ne décrivent aucune entité : la lecture juste estime $\lambda$ en fonction de la taille et
la lit à celle de l'entité visée. La base PRC entière compte $2\,150$ notifications par an ; le
$341$ ne porte que sur les organisations financières.}

\laphrase{Le facteur $3\,719$ entre ma plus grande et ma plus petite fréquence n'est pas une
incohérence, c'est trois niveaux d'observation. Un seul alimente le moteur, les deux autres servent
à répartir et à interpréter. Le seuil de collecte est ce qui les sépare.}

\lepiege{Le taux $\lambda$ est une propriété d'\textbf{entité}, lisible sur un panel
firme-année ; la charge systémique commune, elle, ne se lit qu'au niveau \textbf{agrégé} et
après retrait de tendance. Les confondre revient à imputer au risque systémique ce qui n'est
qu'un effet de taille des firmes.}

\section{Mesurer un biais au lieu de le déclarer}

\lequoi{Tout le monde déclare que sa base est biaisée. Presque personne ne dit de combien. Or
déclarer ne coûte rien, mesurer engage.}

\lidee{Deux biais classiques d'une base de pertes se mesurent avec la base elle-même. Le
\textbf{biais de taille}, parce que la base porte les actifs de la firme sinistrée. La
\textbf{troncature à droite}, parce qu'elle porte la date de saisie en plus de l'année de
survenance.}

\laformule{Pour la taille, une régression log-log $\ln L = a + b\ln(\text{actifs})$ donne
l'élasticité $b$. Mais elle est biaisée par le seuil de collecte : une petite firme n'entre
dans la base que si sa perte est anormalement grande, ce qui \emph{aplatit} la pente. On
corrige par un maximum de vraisemblance sur une lognormale \textbf{tronquée à gauche} en $u_0$,
où la contribution de chaque observation est divisée par $1-\Phi\bigl((\ln u_0-\mu)/s\bigr)$.
Pour la troncature à droite, on construit le triangle survenance $\times$ délai de remontée et
l'on applique un \textbf{chain-ladder} aux comptes.}

\cheznous{Élasticité brute $0{,}042$ avec un $R^2$ de $0{,}005$ (matériel complémentaire), ce qui
\emph{réfuterait} l'explication usuelle. Corrigée de la troncature, elle double à
$\mathbf{0{,}087}$ $[0{,}026\,;0{,}148]$ et devient significative, mais reste faible : la sévérité
transposée à l'entité notionnelle n'est multipliée que par $0{,}85$ environ ($0{,}854$ au
chapitre~10). Côté remontée, le délai médian est de trois ans (matériel complémentaire) et les
$97$ événements observés sur 2023--2025 deviennent \textbf{$255$} à l'ultime, soit $+163\,\%$.}

\laphrase{Mes deux biais les plus souvent opposés, je ne me contente pas de les signaler : je
les chiffre. Le biais de taille coûte quinze pour cent environ, pas un facteur dix. Et la baisse
apparente des incidents après 2022 n'est pas une baisse du risque, c'est une remontée
incomplète que le chain-ladder rattrape.}

\lepiege{Les deux corrections vont en sens opposés et se compensent presque exactement sur la
\emph{perte moyenne} ($0{,}861\times1{,}147=0{,}987$). Elles ne se compensent \textbf{pas} sur
le capital : à queue lourde, le quantile suit la sévérité et non la fréquence. Deux biais qui
s'annulent sur l'espérance peuvent ne pas s'annuler sur le quantile, et c'est le quantile qui
fait le SCR.}

% =====================================================================
\part*{II. Combien d'incidents par an}
\addcontentsline{toc}{part}{II. Combien d'incidents par an}

\section{Le processus de Poisson, et pourquoi il ne suffit pas}

\lequoi{Modéliser un nombre d'événements par an, sans savoir quand ils tombent.}

\lidee{Si les événements arrivent indépendamment, à un rythme constant, et jamais deux
exactement en même temps, alors leur nombre annuel suit une loi de Poisson. C'est le modèle
de comptage par défaut, et il n'a qu'un paramètre.}

\laformule{$\mathbb{P}(N=k)=e^{-\lambda}\lambda^k/k!$, avec la propriété caractéristique
\[
  \mathbb{E}[N]=\mathrm{Var}(N)=\lambda .
\]
Moyenne et variance sont \emph{égales}. C'est à la fois sa force et sa faiblesse.}

\cheznous{Cette égalité est testable, et elle est rejetée massivement. Sur $14\,944$
cellules entité-année, un test du rapport de vraisemblance donne $p\approx10^{-30}$.
Autrement dit~: la probabilité d'observer nos données si Poisson était vrai est de l'ordre
d'un sur $10^{30}$.}

\laphrase{Poisson impose que la variance égale la moyenne. Sur les cellules firme-année, la
variance dépasse la moyenne d'un facteur $1{,}19$ : cela paraît peu, mais sur quinze mille
cellules c'est un rejet à $10^{-30}$. Donc j'ai changé de loi.}

\lepiege{Ne pas dire \og Poisson sous-estime le risque\fg{}. Poisson donne la \emph{bonne}
moyenne, il se trompe sur la \emph{dispersion}, donc sur la queue. Le SCR est un quantile
extrême~: c'est la dispersion qui le fait, pas la moyenne.}

\section{La surdispersion et la négative binomiale}

\lequoi{Les données montrent beaucoup plus de variabilité que Poisson n'en autorise. Pourquoi,
et comment en tenir compte ?}

\lidee{Parce que le rythme lui-même varie. Certaines années sont plus exposées que
d'autres, certaines entités plus que d'autres. On rend donc $\lambda$ \emph{aléatoire} au
lieu de le fixer. Un mélange de Poisson par une loi Gamma donne exactement la négative
binomiale.}

\laformule{Si $N\mid\Lambda\sim\mathcal{P}(\Lambda)$ et $\Lambda\sim\mathrm{Gamma}$, alors
$N$ est négative binomiale et
\[
  \mathbb{E}[N]=\lambda, \qquad \mathrm{Var}(N)=\varphi\,\lambda, \qquad \varphi>1 .
\]
Le facteur $\varphi$ est la \textbf{surdispersion}~: le rapport variance sur moyenne.}

\cheznous{L'indice de dispersion d'un mélange de Poisson vaut $1+\lambda/r$ : il n'est pas
invariant d'échelle. Il vaut $1{,}19$ sur les cellules firme-année, $2{,}24$ sur la série annuelle
détendancée, et le moteur retient $\varphi=9{,}20$, choix de prudence posé et non mesure. Ce
$9{,}20$ et la charge systémique $a=0{,}60$ du modèle multi-états décrivent la même surdispersion
agrégée à $5\,\%$ près ($a=0{,}57$). À variance égale, la \emph{forme} du mélange (Gamma ou
lognormale) déplace encore la queue du nombre de $11\,\%$ ($74$ contre $82$).}

\laphrase{La négative binomiale, c'est Poisson dont le rythme est lui-même incertain.
Concrètement, certaines années sont plus exposées que d'autres, et ce sont ces années-là
qui font le capital.}

\lepiege{Le changement de loi de comptage \textbf{ne déplace pas la moyenne}, il déplace la
VaR, et de $+9\,\%$ seulement : à $\xi\approx0{,}6$, le quantile annuel est porté par la perte
unique dominante. Sur la série détendancée, la surdispersion mesurée ($2{,}24$) est quatre fois
plus faible que celle retenue ($9{,}20$) : le calage est conservateur, pas calibré, et il faut le
dire soi-même.}

\section{Pourquoi pas un processus de Hawkes}

\lequoi{Un processus auto-excité, où chaque événement augmente temporairement la
probabilité du suivant, serait un candidat naturel pour de la contagion.}

\lidee{Hawkes modélise la contagion \emph{dans le temps}~: un incident aujourd'hui rend un
autre incident plus probable demain. Notre cascade modélise la contagion \emph{entre
domaines}~: un incident sur un pilier rend un incident plus probable sur un autre pilier.
Ce ne sont pas les mêmes axes.}

\cheznous{Le test a été fait, sur la chronologie PRC datée, contre les noyaux de la
littérature (Bessy-Roland, Boumezoued et Hillairet ; Boumezoued, Cherkaoui et Hillairet).
L'excitation mesurée est intra-journalière, avec un pic à $3{,}3$~heures, et le ratio de
branchement de $0{,}551$ tombe vers zéro dès que les co-occurrences du même jour sont traitées
comme exogènes. Surtout, une cascade \emph{sans aucune} auto-excitation, simulée puis ajustée comme
une chronologie réelle, donne déjà $0{,}480$ au jour : à la résolution disponible, les deux
modèles ne sont pas distinguables.}

\laphrase{Hawkes décrit la contagion dans le temps, ma cascade décrit la contagion entre domaines.
Je l'ai testé contre les noyaux de la littérature, et à la résolution des données il n'est pas
distinguable d'une cascade sans auto-excitation. Je retiens la cascade par parcimonie, parce
qu'elle porte en plus la structure entre piliers.}

\lepiege{Ne pas dire que Hawkes est rejeté ou réfuté : c'est une équivalence observationnelle
suivie d'un choix de parcimonie. Attention au sens : dégrader la résolution fait \emph{monter} le
ratio de branchement estimé. Et ne pas présenter ce choix comme une critique de cette littérature,
qui travaille sur des données horodatées plus finement.}

% =====================================================================
\part*{III. Combien coûte un incident}
\addcontentsline{toc}{part}{III. Combien coûte un incident}

\section{Pourquoi la théorie des valeurs extrêmes}

\lequoi{Le SCR est un quantile à $99{,}5\,\%$. On veut donc décrire précisément la région
où l'on a le moins d'observations. Ajuster une loi sur l'ensemble des données optimise
l'ajustement là où il y a le plus de points, c'est-à-dire au \emph{centre}, exactement là
où l'on s'en moque.}

\lidee{On abandonne l'idée d'une loi unique pour tout. On modélise \emph{seulement} la
queue, avec une famille dont on peut démontrer qu'elle est la seule limite possible.}

\laphrase{Une loi ajustée sur toutes les données est calée sur le centre, or je n'ai besoin
que de la queue. La théorie des valeurs extrêmes dit quelle famille de lois est la seule
possible tout au bout, et je n'ajuste que celle-là.}

\section{Les deux théorèmes fondateurs}

\lidee{Deux résultats, structurés exactement comme le théorème central limite~: ils disent
qu'\emph{une seule} famille peut apparaître à la limite, quelle que soit la loi de départ.}

\laformule{\textbf{Fisher-Tippett-Gnedenko.} Si le maximum renormalisé de $n$ variables
indépendantes converge en loi, la limite ne peut être que la loi des valeurs extrêmes
généralisée~:
\[
  G_\xi(x)=\exp\!\left\{-\left(1+\xi x\right)^{-1/\xi}\right\}, \qquad 1+\xi x>0 .
\]
\textbf{Pickands-Balkema-de Haan.} De façon équivalente, la loi des \emph{excès} au-dessus
d'un seuil $u$ converge, quand $u$ monte vers le point terminal, vers la loi de Pareto
généralisée~:
\[
  \mathbb{P}(X-u\le y \mid X>u)\;\longrightarrow\;
  1-\left(1+\frac{\xi y}{\sigma}\right)^{-1/\xi}.
\]}

\cheznous{C'est le second qu'on utilise, parce qu'il exploite \emph{tous} les dépassements
de seuil et pas seulement le maximum de chaque bloc. On perd beaucoup moins d'information.}

\laphrase{C'est le même genre de résultat que le théorème central limite. Là où le TCL dit
que les moyennes tendent vers une gaussienne quelle que soit la loi de départ, ces deux
théorèmes disent que les extrêmes tendent vers une seule famille. Je n'ai donc pas
\emph{choisi} la loi de queue, elle est imposée.}

\lepiege{Ces théorèmes sont \emph{asymptotiques}. Ils ne garantissent pas que la GPD est
correcte à un seuil fini~: ils garantissent qu'aucune autre famille ne peut l'être à la
limite. C'est plus faible, il faut le dire, et c'est pourquoi on teste l'adéquation.}

\section{La méthode des excès et la loi de Pareto généralisée}

\lidee{On fixe un seuil $u$. En dessous, on garde la distribution empirique telle quelle. Au
dessus, on ajuste une GPD. La loi obtenue est dite \emph{recollée} (spliced).}

\laformule{La fonction de survie s'écrit, pour $x>u$~:
\[
  \mathbb{P}(X>x)\;=\;p_u\left(1+\frac{\xi (x-u)}{\sigma}\right)^{-1/\xi},
  \qquad p_u=\mathbb{P}(X>u).
\]}

\cheznous{$\xi=0{,}5954$ sur OpRisk, $\xi=1{,}03$ sur PRC. Le paramètre $\xi$ est
l'\textbf{indice de queue}, et il commande tout~:
\begin{itemize}[itemsep=1pt,topsep=2pt]
\item $\xi<1/2$~: variance finie ;
\item $1/2\le\xi<1$~: variance \emph{infinie}, moyenne finie : c'est le cas de la calibration
publiée ;
\item $\xi\ge1$~: moyenne elle-même infinie.
\end{itemize}
Le $\xi=1{,}03$ de la PRC dépasse $1$, mais il ne mesure pas un risque plus lourd : la sévérité
PRC est une puissance du nombre d'enregistrements, et une puissance multiplie l'indice de queue par
son exposant. C'est pourtant ce régime qui fait tomber la TVaR comme mesure de couverture : sous
lui, elle n'existe pas.}

\laphrase{L'indice de queue mesure à quel point les catastrophes dominent. Sur ma calibration il
vaut $0{,}6$ : la variance des pertes est infinie, leur moyenne finie, et un seul sinistre peut
peser plus que tous les autres réunis. C'est ce qu'on observe en cyber : sur OpRisk, le décile
supérieur porte plus de $92\,\%$ de la perte totale.}

\lepiege{Avec $\xi$ proche de $1$, la moyenne empirique n'est pas un estimateur stable et
les intervalles de confiance sont larges par nature. Ne jamais annoncer un SCR ponctuel
sans sa bande dans ce régime.}

\section{Le choix du seuil}

\lequoi{Le seuil $u$ arbitre entre biais et variance~: trop bas, on inclut le corps de la
distribution et l'approximation GPD est fausse ; trop haut, il reste trop peu de points et
l'estimation devient bruitée.}

\lidee{On ne le choisit pas au hasard : on le lit sur des diagnostics, puis on vérifie qu'une
règle écrite le retrouve.}

\laformule{\textbf{Vie résiduelle moyenne.} Pour une GPD, la fonction
$e(u)=\mathbb{E}[X-u\mid X>u]$ est \emph{linéaire} en $u$ : on choisit $u$ au début de la zone
linéaire. \textbf{Diagnostic GPD.} On trace $\hat\xi$ en fonction de $u$ : le seuil doit se
situer là où $\xi$ passe sous $1$ (espérance finie) en gardant au moins trente excès.}

\cheznous{Seuil retenu $u=20{,}03$~M\euro{} sur le périmètre cyber $\times$ finance d'OpRisk, soit
$91$ excès, et $p_u=0{,}1509$, valeur gelée (le recomptage donnerait $91/582=0{,}1564$, écart
chiffré au chapitre~11). Une règle de stabilité écrite, appliquée à dix percentiles candidats,
retrouve $19{,}87$~M\euro{} avec le même quantile ; la règle la plus permissive descendrait à
$12{,}88$ et donnerait un quantile supérieur de $16\,\%$. Attention : $\hat\xi$ n'est stable qu'au
\emph{voisinage} du seuil ; de $0{,}60$ il descend à $0{,}38$ quand on remonte au percentile $95$,
dans l'intervalle déjà publié.}

\laphrase{Le seuil se lit sur la vie résiduelle moyenne, qui devient droite au-delà de vingt
millions, et une règle de stabilité écrite le retrouve à un pour cent près. Et il ne m'arrange pas :
la règle la plus permissive donnerait un capital plus élevé.}

\lepiege{Ne pas dire que le graphique de Hill confirme le seuil : Hill vaut $1{,}32$ au seuil
retenu, plus du double du maximum de vraisemblance, et le mémoire en fait une corroboration d'un
autre genre, puisque c'est exactement ce que Hill produit en échantillon fini sous $\xi=0{,}595$.
Et ne pas confondre avec les $4{,}176$~M\euro{} d'une table descriptive : c'est le percentile $85$
de la sévérité PRC dérivée, jamais un seuil.}

% =====================================================================
\part*{IV. Comment les pertes s'additionnent}
\addcontentsline{toc}{part}{IV. Comment les pertes s'additionnent}

\section{Le modèle collectif}

\lidee{On ne modélise pas la perte annuelle directement. On modélise séparément
\emph{combien} d'incidents et \emph{combien coûte} chacun, puis on somme. C'est le cadre
standard en actuariat non-vie.}

\laformule{
\[
  L \;=\; \sum_{i=1}^{N} X_i,
\]
avec $N$ le nombre d'incidents (fréquence) et $X_i$ leurs coûts (sévérité), les $X_i$
indépendants de $N$ et entre eux.}

\cheznous{Des dizaines de milliers d'années simulées par Monte-Carlo, avec le bruit de simulation
publié à côté de chaque grandeur. La loi de $L$ n'a pas de forme fermée simple avec une queue
lourde et une cascade, mais elle se calcule sans simuler par inversion de Fourier de sa fonction
génératrice : c'est le contrôle du moteur, qui retrouve le quantile à moins d'une erreur type
($5\,945$ contre $6\,049$ au conforme).}

\laphrase{Je sépare deux questions qui n'ont rien à voir~: combien d'incidents par an, et
combien coûte un incident. Chacune se calibre sur ses propres données, et la perte annuelle
est leur somme.}

\lepiege{L'indépendance entre $N$ et les $X_i$ est une hypothèse, et elle a été testée hors
échantillon : elle tient sur les excès. La pente négative qu'on voit sur toute la distribution
disparaît au-delà du seuil, c'est un artefact de profondeur de collecte. Et la cascade ne rend pas
$N$ et $X$ dépendants : elle fait dépendre le coût d'un sinistre du nombre de piliers qu'il
touche.}

\section{VaR, TVaR, et pourquoi le SCR est une VaR}

\lidee{La VaR à $99{,}5\,\%$ est le montant qu'on ne dépasse que dans une année sur deux
cents. La TVaR est la moyenne des pertes \emph{au-delà} de ce seuil.}

\laformule{
\[
  \mathrm{VaR}_\alpha(L)=\inf\{x:\ \mathbb{P}(L\le x)\ge\alpha\},
  \qquad
  \mathrm{TVaR}_\alpha(L)=\mathbb{E}[L\mid L>\mathrm{VaR}_\alpha(L)].
\]}

\cheznous{Solvabilité~II impose la VaR à $99{,}5\,\%$ à horizon un an, et le mémoire la prend
\emph{non centrée} : la perte attendue d'une non-conformité n'étant couverte par aucune provision,
elle n'est pas déduite. La TVaR sert de diagnostic, jamais de couverture : noyau de l'allocation
d'Euler, et indicateur de forme, $\mathrm{TVaR}/\mathrm{VaR}\approx2{,}1$ contre l'asymptote
$1/(1-\xi)=2{,}5$.}

\laphrase{Le capital, c'est la perte annuelle qu'on ne dépasse qu'une année sur deux cents. Ce
n'est pas mon choix de mesure, c'est celui de Solvabilité~II. La TVaR, je m'en sers pour
diagnostiquer la queue, pas pour couvrir.}

\lepiege{La VaR n'est pas sous-additive, la TVaR est cohérente : c'est vrai, et le mémoire ne dit
pas le contraire. Mais l'argument décisif contre une couverture en TVaR n'est pas de degré, il est
d'existence : sous la calibration PRC, $\xi=1{,}03>1$, l'espérance est infinie et la TVaR n'existe
pas. Une mesure de couverture qui cesse d'être définie selon la source ne peut pas porter un
capital.}

% =====================================================================
\part*{V. Comment un problème en entraîne un autre}
\addcontentsline{toc}{part}{V. Comment un problème en entraîne un autre}

\section{La convention d'indices, à ne jamais improviser}

\lequoi{Une matrice dirigée a deux lectures possibles, et les deux sont défendables. Si le
document en utilise l'une ici et l'autre là, tous ses énoncés directionnels deviennent
ambigus.}

\lidee{On fixe la convention une fois, on l'écrit, et on la vérifie sur un fait connu.}

\laformule{\[
  W_{jk}\;=\;\text{propension de } j \text{ (\textbf{source}) à entraîner la défaillance de }
  k \text{ (\textbf{cible}).}
\]
La \textbf{ligne} $j$ décrit ce que $j$ \emph{émet}, la \textbf{colonne} $k$ ce que $k$
\emph{reçoit}. La somme de ligne $s_j$ mesure donc une \textbf{émission}.}

\cheznous{Vérification sur la matrice d'expert : P1 a la plus grande somme de ligne ($2{,}6$)
et une somme de colonne trois fois inférieure à celle de P2. C'est exactement la revendication
centrale \og{}P1 émet fort et reçoit peu\fg{}, et elle serait \emph{fausse} sous la convention
inverse. La forme réduite du modèle latent s'écrit alors $(I-W^{\top})^{-1}$, le rayon spectral
étant de toute façon invariant par transposition.}

\laphrase{La ligne émet, la colonne reçoit. Je le vérifie sur ma matrice : le pilier que je dis
émetteur a bien la plus grosse somme de ligne. Une convention n'est pas un détail de notation,
c'est ce qui rend un énoncé directionnel falsifiable.}

\lepiege{Ce défaut a réellement existé dans le mémoire : l'équation structurelle sommait sur
les lignes alors que la définition de la cascade et tout le code sommaient sur les colonnes. Le
code et le classeur ont tranché. Vérifier une convention sur un fait connu coûte trente
secondes et évite une matrice transposée.}

\section{La matrice de contagion et sa normalisation}

\lequoi{On veut dire \og un problème sur le pilier $j$ aggrave le pilier $k$\fg{}, avec une
intensité. Mais si on met des coefficients au hasard, la cascade peut exploser.}

\begin{center}
\begin{tikzpicture}[scale=1.0]
\node[pilier] (p1) at (90:2.0)  {P1};
\node[pilier] (p2) at (18:2.0)  {P2};
\node[pilier] (p3) at (-54:2.0) {P3};
\node[pilier] (p4) at (234:2.0) {P4};
\node[pilier] (p5) at (162:2.0) {P5};
\node[etiq,above=1pt of p1] {gouvernance TIC};
\node[etiq,above right=0pt and 1pt of p2] {incidents};
\node[etiq,below=1pt of p3] {tests};
\node[etiq,below=1pt of p4] {tiers};
\node[etiq,left=1pt of p5]  {partage};
% P1 emet fort, recoit peu
\draw[flechef] (p1) to[bend left=12] (p2);
\draw[flechef] (p1) to[bend left=12] (p3);
\draw[flechef] (p1) to[bend right=12] (p4);
\draw[fleche]  (p4) to[bend right=12] (p2);
\draw[fleche]  (p4) to[bend left=12]  (p3);
\draw[flechep] (p5) to[bend left=12]  (p1);
\draw[flechep] (p3) to[bend left=12]  (p2);
\node[align=left,font=\scriptsize,text=navy!75,anchor=west] at (4.3,0.9)
  {\textbf{Lire la matrice~:}\\$W_{jk}$ = propension de $j$\\à entraîner $k$.\\
   \emph{La ligne émet, la colonne reçoit.}};
\node[align=left,font=\scriptsize,text=accent,anchor=west] at (4.3,-0.9)
  {P1 émet fort et reçoit peu.\\C'est pourquoi il sort\\premier en vue Shapley.};
\end{tikzpicture}
\end{center}

\noindent La cascade est un \emph{processus de branchement}~: un incident amorcé sur un pilier
tire, pour chaque arc, une pièce de Bernoulli de probabilité $W_{jk}$, et recommence depuis les
piliers atteints. Sans garde-fou, ce processus peut ne jamais s'éteindre.

\lidee{On part du jugement dirigé (la matrice \textsc{trans} du classeur), et on la
renormalise pour garantir mathématiquement que la propagation s'éteint.}

\laformule{
\[
  W \;=\; g\,\frac{\mathrm{TRANS}}{\max_j \sum_k \mathrm{TRANS}_{jk}},
  \qquad g=0{,}90 ,
\]
ce qui garantit $\rho(W)\le g<1$, où $\rho$ est le rayon spectral (la plus grande valeur
propre en module).}

\cheznous{Le facteur $g$ est un curseur de propagation~: $g=0$ signifie aucune contagion,
$g\to1$ signifie une cascade au bord de l'explosion. On le fixe à $0{,}90$, ce qui est
volontairement élevé.}

\laphrase{La normalisation garantit que la cascade s'éteint. Sans elle, un choc pourrait
s'amplifier indéfiniment en tournant en boucle entre les piliers, et le capital serait
infini.}

\lepiege{Les valeurs de $g$, $0{,}45$ au conforme, $0{,}68$ au partiel et $0{,}90$ au non
conforme, sont des \emph{choix}, pas des mesures. À dire ainsi : \og elles sont posées, mais la
thèse n'en dépend pas : le capital est croissant en $g$, donc toute correspondance qui respecte
l'ordre reproduit l'écart, et seule son amplitude est un scénario\fg{}.}

\section{L'inverse de Leontief et la sous-criticité}

\lidee{Le choc initial se propage, ce qu'il déclenche se propage à son tour, et ainsi de
suite. La somme de tous les tours est une série géométrique matricielle.}

\laformule{
\[
  (I-W)^{-1} \;=\; I + W + W^2 + W^3 + \cdots
\]
Cette série converge \emph{si et seulement si} $\rho(W)<1$. Le terme $I$ est le choc
initial, $W$ ce qu'il déclenche directement, $W^2$ ce que cela déclenche à son tour.}

\cheznous{L'objet vient de l'économie (Leontief, tableaux entrées-sorties)~: la même algèbre
décrit comment une commande dans un secteur en génère dans les autres. Sur la matrice du mémoire,
$\rho(W)=0{,}562\,g$, soit $0{,}506$ à $g=0{,}9$ ; le seuil critique $\rho(W)=1$ serait atteint à
$g=1{,}78$, hors du domaine admissible.}

\laphrase{C'est la même formule que les tableaux entrées-sorties en économie. Le premier
terme est le choc direct, le deuxième ce qu'il déclenche, le troisième ce que cela déclenche
à son tour. La condition de convergence est exactement la condition pour que la cascade
s'éteigne.}

\lepiege{Le passage à $W^2$ est celui où la \emph{direction} commence à compter, alors
qu'elle ne compte pas dans $W$ seul si on ne regarde que les co-occurrences. C'est le point
technique le plus important du mémoire, voir la section~\ref{sec:trace}.}

\section{Le seuil de déclenchement, et le rôle de la loi normale}

\lequoi{À partir de quand un pilier bascule ? Il faut un mécanisme, pas seulement une
probabilité.}

\lidee{On suppose une variable latente $X_j$ qui mesure la fragilité du pilier $j$, et le
pilier bascule quand elle franchit un seuil. C'est le mécanisme de Merton, repris par
Vasicek pour le risque de crédit.}

\laformule{Le pilier $j$ bascule quand
\[
  X_j \;\ge\; K_j \;=\; F^{-1}(1-p_j),
\]
où $p_j$ est sa probabilité marginale de basculement et $F$ la fonction de répartition de la
latente.}

\cheznous{$K_4=1{,}03$ pour le pilier des prestataires, dont la probabilité de déclenchement
vaut $p_4=0{,}151$.}

\laphrase{Il y a deux questions derrière le mot seuil, et il ne faut pas les mélanger. Pour
savoir \emph{quand un pilier bascule tout seul}, la loi normale est inoffensive~: elle ne
sert que d'unité de mesure, la vraie queue lourde est dans la sévérité, pas ici. Pour savoir
\emph{combien de piliers basculent ensemble}, c'est autre chose, et là le gaussien pose un
vrai problème.}

\lepiege{Le choix de la loi de la latente est une \emph{normalisation} pour la marge~:
n'importe quelle loi continue donnerait la même probabilité marginale, seul le seuil change
de valeur. Ce n'est pas vrai pour la dépendance. Confondre les deux est l'erreur classique.}

\section{Les copules et la dépendance de queue}

\lequoi{Comment plusieurs piliers basculent-ils \emph{ensemble} ? La probabilité marginale
de chacun ne suffit pas à le dire.}

\lidee{Une copule sépare les marges de la structure de dépendance. On peut garder les mêmes
marges et changer complètement la façon dont les événements extrêmes se coordonnent.}

\laformule{Le coefficient de \textbf{dépendance de queue supérieure} mesure la
co-occurrence extrême~:
\[
  \lambda_U \;=\; \lim_{q\to1^-}\ \mathbb{P}\bigl(U_2>q \mid U_1>q\bigr).
\]
Pour la copule gaussienne, $\lambda_U = 0$ dès que la corrélation est $<1$. Pour la copule
de Gumbel de paramètre $\theta$,
\[
  \lambda_U \;=\; 2-2^{1/\theta}.
\]}

\cheznous{Avec $\theta=1{,}8$, on a $\lambda_U = 0{,}53$. À $\tau$ de Kendall \emph{égal},
la gaussienne donne $0$ et Gumbel $0{,}53$. La probabilité que les cinq piliers basculent
ensemble passe de $0{,}17$ à $0{,}39$ selon la copule retenue. On encadre au lieu de
trancher~: plancher gaussien, plafond de Gumbel. Le co-déclenchement d'un autre pilier avec P4
tient alors dans $[0{,}48\,;0{,}58]$, sous la borne $\gamma=0{,}68$ du choc commun, qui reçoit
ainsi son origine structurelle.}

\laphrase{Une dépendance gaussienne a une propriété gênante~: dans les extrêmes, elle prédit
que les choses se \emph{découplent}. C'est le contraire de ce qu'on observe en crise. Donc
je n'ai pas tranché entre les deux, j'ai encadré, et l'écart entre le plancher et le plafond
est un résultat en soi.}

\lepiege{Deux copules peuvent avoir le \emph{même} $\tau$ de Kendall, donc la même
dépendance \og moyenne\fg{}, et des comportements extrêmes opposés. Comparer deux copules par
leur corrélation ne dit rien sur la queue.}

% =====================================================================
\part*{VI. Ce qu'on ne peut pas savoir, et pourquoi}
\addcontentsline{toc}{part}{VI. Ce qu'on ne peut pas savoir, et pourquoi}

\section{La décomposition $W = S + A$}

\lidee{Toute matrice se coupe de manière unique en une partie symétrique et une partie
antisymétrique. La première dit \emph{qui est lié à qui}, la seconde dit \emph{dans quel
sens}.}

\laformule{
\[
  S=\tfrac12\bigl(W+W^{\top}\bigr), \qquad A=\tfrac12\bigl(W-W^{\top}\bigr),
  \qquad W=S+A .
\]
Les deux parties sont \textbf{orthogonales}~: $\langle S,A\rangle=0$, d'où un Pythagore
$\|W\|^2=\|S\|^2+\|A\|^2$.}

\cheznous{Sur la matrice ordonnée estimée à partir du corpus de rapports, la part de direction
$\|A\|/\|W\|$ vaut $60\,\%$, et elle ne bouge pas quand on corrige le dénominateur, alors que le
niveau des $p_{jk}$ tombe de $0{,}27$ à $0{,}07$. La direction n'est pas un détail de la matrice.}

\laphrase{On coupe la matrice en deux~: une partie qui dit quels piliers sont liés, et une partie
qui dit dans quel sens. Les deux sont orthogonales, et la co-occurrence que voit la donnée ne
dépend, au premier ordre, que de la première.}

\lepiege{La décomposition propre porte sur le \emph{flux stationnaire} $\pi_j P_{jk}$, pas
sur la matrice de transition brute. Notre $S=(W+W^{\top})/2$ en est le cas où $\pi$ est
uniforme, et l'écart mesuré est de $0{,}069$, donc faible. Le dire soi-même met à l'abri.}

\medskip
\noindent\textbf{Le schéma qui fait comprendre la thèse en dix secondes.} Deux matrices
opposées, une seule statistique observable, deux décisions contraires.

\begin{center}
\begin{tikzpicture}
% cas W
\node[pilier] (a1) at (0,1.1) {P1};
\node[pilier] (b1) at (0,-0.5) {P2};
\draw[flechef] (a1) -- (b1);
\node[font=\scriptsize\bfseries,text=navy] at (0,1.9) {matrice $W$};
\node[etiq,align=center] at (0,-1.3) {P1 entraîne P2\\\emph{remédier P1}};
% cas W transpose
\node[pilier] (a2) at (3.4,1.1) {P1};
\node[pilier] (b2) at (3.4,-0.5) {P2};
\draw[flechef] (b2) -- (a2);
\node[font=\scriptsize\bfseries,text=navy] at (3.4,1.9) {matrice $W^{\top}$};
\node[etiq,align=center] at (3.4,-1.3) {P2 entraîne P1\\\emph{remédier P2}};
% ce que la donnee voit
\node[draw=vert!70!black,fill=vert!8,rounded corners=2pt,align=center,
      font=\footnotesize,inner sep=6pt] (obs) at (8.4,0.3)
  {\textbf{Ce que la donnée agrégée voit}\\
   \og P1 et P2 tombent ensemble\fg{}\\[2pt]
   au premier ordre, $\Prob(j \text{ et } k)$\\
   ne dépend que de $S$, et $S(W)=S(W^{\top})$};
\draw[fleche] (1.0,0.3) -- (obs.west |- 0,0.3);
\draw[fleche] (4.4,0.3) -- (obs.west |- 0,0.3);
\node[align=center,font=\scriptsize,text=accent] at (4.25,-2.3)
  {\textbf{Indiscernables pour la donnée. Matériellement différentes pour la décision.}};
\end{tikzpicture}
\end{center}

\noindent C'est l'énoncé de la frontière d'identification, et il ne repose sur \emph{aucun} dire
d'expert : au premier ordre et sous amorce uniforme, la co-occurrence est une fonction de $S$ seule,
et $S$ est identique pour $W$ et pour sa transposée. \textcolor{accent}{Ce n'est pas une
impossibilité de principe} : au-delà du premier ordre, la loi complète des piliers touchés par
incident distingue $W$ de $W^\top$, une fourche d'un collisionneur ($62{,}31\,\%$ de sinistres
multi-piliers sous $W$, $50{,}77\,\%$ sous $W^\top$). La frontière est celle des statistiques que
les sources fournissent, et c'est pourquoi le registre recommandé consigne tous les piliers
touchés.

\section{Le renversement du temps et la fluctuation martingale}

\lequoi{Pourquoi n'arrive-t-on pas à mesurer $A$ ? Et est-ce une limite de méthode ou une
limite de principe ?}

\lidee{Passer le film à l'envers. La partie symétrique ne change pas, la partie
antisymétrique s'inverse. Et la fluctuation d'un processus, une fois la tendance retranchée,
ne dépend que de la première.}

\laformule{La chaîne renversée a pour flux $F^{*}=F^{\top}=S-A$. Et pour toute fonction test
$f$, l'intensité de la fluctuation vaut
\[
  \mathbb{E}_\pi\bigl[\Gamma(f,f)\bigr]
  \;=\;\sum_{j,k} S_{jk}\bigl(f(k)-f(j)\bigr)^2 ,
\]
la partie $A$ y contribuant exactement zéro. La démonstration tient en une ligne~:
$(f(k)-f(j))^2$ est symétrique en $(j,k)$, $A$ est antisymétrique, les termes $(j,k)$ et
$(k,j)$ se compensent.}

\cheznous{Vérifié sur trois cents fonctions test tirées au hasard : l'énergie de fluctuation est
la même pour la chaîne posée et pour sa symétrisée, les points tombent sur la diagonale (figure du
chapitre~8, \S~8.3).}

\laphrase{Renverser le temps laisse la partie symétrique intacte et retourne le signe de la partie
directionnelle. Or la fluctuation d'un processus, sa partie martingale, ne dépend que de la
symétrique. Donc une statistique invariante par renversement du temps ne peut pas voir la
direction. La loi complète des piliers touchés par incident, elle, le pourrait : c'est une donnée
qui manque, pas une impossibilité.}

\lepiege{Ne pas dire que $W$ \emph{est} une chaîne de Markov. C'est un noyau de contagion~: on
emprunte l'algèbre de la réversibilité comme cadre de lecture, on n'ajoute pas une hypothèse. Et le
mot martingale s'entend au sens des chaînes de Markov, sous la mesure physique : aucune mesure
risque-neutre, aucun test de martingalité.}

\section{Le test placebo}
\label{sec:placebo}

\lequoi{Même des données sans aucune direction produisent un déséquilibre, par simple
fluctuation d'échantillonnage. Observer un déséquilibre ne prouve donc rien tant qu'on ne
sait pas ce que le hasard produit tout seul.}

\lidee{On fabrique une version des données dont la direction a été \textbf{détruite
volontairement}, tout le reste étant conservé, et on compare. C'est l'essai clinique~: si le
médicament ne fait pas mieux que la pilule vide, il ne fait rien.}

\laformule{On mesure l'asymétrie de la matrice de transitions comptées~:
\[
  \mathcal{A}(M)=\frac{\|M-M^{\top}\|_1}{2}
  =\frac12\sum_{j,k}\bigl|M_{jk}-M_{kj}\bigr| ,
\]
c'est-à-dire le surplus de transitions d'un sens sur l'autre, sommé sur toutes les paires.
Puis on \textbf{permute les années à l'intérieur de chaque entité}~: l'ensemble des
catégories touchées chaque année est conservé, seul l'ordre est détruit. On recompte, on
recalcule, et on compare~:
\[
  z=\frac{\mathcal{A}(M_{\text{obs}})-\overline{\mathcal{A}^{\,\text{null}}}}
        {\widehat{s}(\mathcal{A}^{\,\text{null}})} .
\]}

\cheznous{
\begin{center}\small
\begin{tabular}{@{}l r r r l@{}}
\toprule
\textbf{Source} & \textbf{observée} & \textbf{sous placebo} & \textbf{$z$} & \textbf{Verdict}\\
\midrule
OpRisk, pas annuel & $119$ & $128\pm27$ & $-0{,}33$ & aucun signal \\
7 post-mortems & \multicolumn{2}{c}{\emph{même protocole}} & $+3{,}93$ & signal net \\
\bottomrule
\end{tabular}
\end{center}
Sur OpRisk, l'observé est \emph{sous} la moyenne du bruit. Un test binomial par paire le
confirme~: aucune paire ne survit à la correction de Bonferroni.}

\laphrase{Je voulais savoir si la direction que je crois voir est réelle. Le problème, c'est
que même des données au hasard produisent un déséquilibre. Alors j'ai pris mes données, j'ai
détruit l'ordre en gardant tout le reste, et j'ai comparé. Mes vraies données font
$119$, le hasard fait $128$~: elles font moins bien que le hasard. Puis j'ai refait
exactement le même test sur les sept rapports d'enquête, et là le signal apparaît nettement.
Même test, autre source, résultat opposé~: le problème venait de la donnée, pas de la
méthode.}

\lepiege{Trois choses à ne pas dire. \textbf{(i)} Le test du \og temps inversé\fg{} est
\emph{dégénéré} sur un comptage~: renverser la chronologie donne $M^{\top}$ par
construction, donc une corrélation de $-1$ quel que soit le contenu. Il ne prouve rien.
\textbf{(ii)} Un test non rejeté n'est pas une preuve d'absence~: la donnée est
\emph{compatible} avec l'absence de direction. \textbf{(iii)} Sur les post-mortems, le test
établit la \emph{cohérence} du codage, pas son \emph{exogénéité}~: le biais de narration des
rapports d'enquête n'est pas levé par le calcul.}

\section{La production d'entropie}

\lidee{Un seul nombre qui dit à quel point le processus est dirigé, et qui a une
interprétation exacte~: à quel point le film à l'endroit est distinguable du film à
l'envers.}

\laformule{
\[
  \sigma=\sum_{j<k} 2A_{jk}\,\ln\!\frac{S_{jk}+A_{jk}}{S_{jk}-A_{jk}} \;\ge\; 0,
\]
nulle si et seulement si $A=0$. Elle est aussi le taux d'entropie relative entre la loi de
la trajectoire lue à l'endroit et celle lue à l'envers.}

\cheznous{$\sigma$ est strictement positive pour la chaîne posée et s'annule à la
symétrisation (figure du chapitre~8, \S~8.3). \textbf{Et la conséquence qui compte}, qui se lit
sur la formule : chaque terme de la somme est \emph{pair} en $A_{jk}$, donc $\sigma$ ne dépend que
des \emph{amplitudes} de la direction, pas de ses signes. Les $1\,024$ sommets de l'ensemble
admissible, qui ne diffèrent que par les signes, ont donc la même production d'entropie, alors
que le capital y varie d'un facteur $1{,}27$.}

\laphrase{La production d'entropie mesure à quel point on peut distinguer le film à l'endroit du
film à l'envers. Si elle est nulle, personne ne peut dire dans quel sens le temps coule. Mais
attention~: elle ne voit que la \emph{quantité} de direction, jamais son \emph{sens}. Les mille
vingt-quatre configurations ont la même production d'entropie, et pourtant elles étalent le
capital de plus d'un quart.}

\lepiege{Conséquence à tirer~: ce n'est pas $\sigma$ qu'il faut chercher à estimer, ce sont
les \emph{signes}, un par un. C'est ce qui justifie que la valeur d'information se
hiérarchise par paire et non globalement.}

\section{$z$ scalaire ou $z$ par paire}

\lequoi{L'écriture $P=S+zA$ est ambiguë tant qu'on n'a pas dit si $z$ est un nombre ou dix.}

\laformule{\textbf{Lecture globale}~: $z\in[-1,1]$ règle l'\emph{intensité} d'un motif de
direction fixé. Un seul degré de liberté. Elle ne peut pas changer quel pilier précède
quel autre.
\textbf{Lecture par paire}, celle du mémoire~:
\[
  A_{jk}=z_{jk}\,t\,S_{jk}, \qquad z_{jk}\in[-1,1], \qquad j<k,
\]
soit \textbf{dix degrés de liberté}. Le paramètre $t$ est l'amplitude maximale d'asymétrie,
les $z_{jk}$ portent chacun un sens indépendant.}

\cheznous{Les $2^{10}=1024$ sommets sont exactement les configurations $z_{jk}=\pm1$. La
contrainte de positivité du flux impose $|A_{jk}|\le S_{jk}$, donc $t\le1$~; $t=1$ est
l'ignorance totale.}

\laphrase{Un $z$ global ne fait que renforcer ou affaiblir un ordre déjà choisi. Ce qu'il me
faut, c'est un curseur par paire, parce que c'est le \emph{sens} de chaque lien que je ne
connais pas, pas son intensité globale.}

\section{Le troisième échec : le schéma existe, il n'est pas rempli}

\lequoi{Deux bases ont échoué à identifier la direction, l'une faute de taxonomie, l'autre
faute de résolution temporelle. On pouvait encore croire que personne n'avait \emph{prévu} de
collecter la bonne donnée.}

\lidee{On teste l'hypothèse la plus favorable : un standard de description d'incidents qui,
lui, prévoit exactement les bons champs. Le schéma \textsc{veris}, référence de la profession
et support du \emph{Data Breach Investigations Report}, porte littéralement les deux objets
manquants : un champ \texttt{control\_failure} qui nomme le domaine de contrôle défaillant, et
une chronologie en quatre étapes qui porte l'antériorité.}

\cheznous{Sur les $10\,591$ incidents de la base publique associée : \texttt{control\_failure}
renseigné dans \textbf{$0{,}12\,\%$} des cas ($13$ incidents, dont un seul en secteur
financier), chronologie datée sur au moins deux étapes dans \textbf{$1{,}19\,\%$}. Et
l'\textbf{intersection} des deux, qui est ce qu'identifier $W$ exigerait simultanément, ne
compte qu'\textbf{un seul incident}, et \textbf{aucun} en secteur financier.}

\laphrase{Le champ qui donnerait la réponse existe dans le standard de la profession. Il est
rempli treize fois sur dix mille incidents. Le problème n'est donc pas de savoir quoi collecter, il
est d'en
rendre le remplissage obligatoire, et c'est précisément ce qu'un régime de notification
réglementaire peut imposer là où une base déclarative volontaire ne le peut pas.}

\lepiege{C'est le plus utile des trois échecs, parce qu'il transforme la recommandation. On ne
demande plus d'inventer une taxonomie, ce qui serait long et contestable, on demande de rendre
obligatoires deux champs \textbf{déjà normalisés}. Coût de conception nul.}

\section{La matrice ordonnée $p_{jk}$}

\lequoi{Le codage documentaire, tel qu'on l'exploite d'abord, compare $n(j\to k)$ à
$n(k\to j)$ \emph{à l'intérieur} d'une paire. Ce rapport élimine le dénominateur, donc il
donne un \emph{sens} et ne dit rien de l'\emph{amplitude}.}

\lidee{On ajoute le dénominateur. Chaque incident déclare non seulement ses transitions, mais
l'ensemble des \textbf{piliers défaillants}, y compris ceux qui ont cédé sans rien entraîner.
On peut alors estimer directement chaque case de la matrice ordonnée.}

\laformule{Avec $N_j$ le nombre d'incidents où $j$ a cédé et $m_{jk}$ ceux où le rapport
établit $j\to k$ :
\[
  p_{jk}=\Prob(k \text{ cède}\mid j \text{ a cédé}),
  \qquad
  p_{jk}\mid\text{données}\ \sim\ \mathrm{Beta}\bigl(1+m_{jk},\ 1+N_j-m_{jk}\bigr).
\]
C'est la matrice ordonnée \emph{complète}, donc $S$ \textbf{et} $A$.}

\cheznous{Sur dix rapports et $22$ transitions : \textbf{douze} liens sur vingt ont un crédible
à $90\,\%$ entièrement \emph{sous} $0{,}5$, et \textbf{aucun} au-dessus. Les plus élevés,
$P_3\to P_2$ ($0{,}77$, $[0{,}48\,;0{,}95]$) et $P_1\to P_4$ ($0{,}68$, $[0{,}40\,;0{,}89]$),
recouvrent encore $0{,}5$. Les douze liens faibles comprennent les \emph{quatre} sorties de
$P_2$, soit toutes, et trois des quatre entrées de $P_1$, l'arc $P_5\to P_1$ n'étant pas établi
comme faible : le puits est établi comme tel, la source pure seulement en partie.}

\laphrase{Mon corpus établit très bien où la contagion \emph{ne va pas}, et pas du tout où elle
va fort. Le puits et la source pure sont donc établis en tant que tels, mais aucun coefficient
n'est établi comme grand. C'est un ordre, ce n'est pas une magnitude.}

\lepiege{Une enquête raconte la \emph{chaîne} qui a produit le sinistre : elle documente bien
les contrôles qui ont cédé \emph{et} propagé, mal ceux qui ont cédé sans propager, pas du tout
ceux qui ont tenu. Le dénominateur $N_j$ est donc sous-estimé et $p_{jk}$ est une \textbf{borne
supérieure}. En ajoutant des défaillances non propagées, le \emph{niveau} s'effondre de
$0{,}27$ à $0{,}07$ tandis que la part de direction $\lVert A\rVert/\lVert W\rVert$ ne bouge
pas ($60\,\%$). Le corpus établit un ordre, pas une magnitude, et cette phrase-là est robuste.}

\section{L'identification partielle}

\lidee{Quand un paramètre n'est pas identifiable, il reste deux attitudes~: le poser
arbitrairement et faire semblant, ou énumérer tout ce qu'il pourrait être et rapporter
l'intervalle qui en résulte. La seconde est l'identification partielle, au sens de Manski.}

\laformule{On définit l'ensemble admissible
\[
  \mathcal{W}=\bigl\{\,W=S+A \ :\ |A_{jk}|\le t\,S_{jk}\ \bigr\},
\]
et l'on rapporte $\bigl[\inf_{\mathcal{W}}\mathrm{SCR},\ \sup_{\mathcal{W}}\mathrm{SCR}\bigr]$
au lieu d'un point.}

\cheznous{Socle sans contagion $5\,275$~M\euro. Avec contagion et ignorance totale, capital
borné dans $[6\,858\,;\,8\,697]$~M\euro, un facteur $1{,}27$ : la contagion ajoute entre $+30$ et
$+65\,\%$ au socle. L'énumération porte sur $1\,024$ sommets, et $3\,000$ points intérieurs
tirés au hasard ne sortent pas de l'enveloppe des sommets, ni en VaR ni en perte moyenne (bornée
dans $[849\,;1\,100]$). Le point d'expert, $8\,110$, est un point de l'ensemble, jamais sa
mesure.}

\laphrase{Je ne prétends pas connaître la direction, donc je ne donne pas un chiffre, je
donne une fourchette. Ce n'est pas de la prudence, c'est la seule réponse possible~: une
grandeur qui compte pour le résultat et qui est invisible dans mes données.}

\lepiege{L'atteinte des extrêmes aux sommets est une \emph{propriété vérifiée
numériquement}, pas un théorème démontré. C'est déclaré comme telle dans le mémoire.}

\medskip
\noindent\textbf{Le schéma à projeter si l'on vous demande \og quel est votre SCR\fg{}.}

\begin{center}
\begin{tikzpicture}[xscale=1.0]
\def\y{0}
% axe
\draw[->,draw=navy!60] (-0.4,\y-1.25) -- (11.2,\y-1.25) node[right,font=\scriptsize,text=navy]{M\euro};
\foreach \x/\v in {0/5275, 3.6/6858, 7.4/8697} {
  \draw[draw=navy!60] (\x,\y-1.15) -- (\x,\y-1.35);
  \node[font=\scriptsize,text=navy] at (\x,\y-1.62) {$\v$};
}
% socle
\fill[vert!25,draw=vert!70!black] (0,\y-0.35) rectangle (3.6,\y+0.35);
\node[font=\scriptsize,align=center,text=vert!35!black] at (1.8,\y) {\textbf{socle}\\sans contagion};
% bande d'ignorance
\fill[accent!18,draw=accent!80!black] (3.6,\y-0.35) rectangle (7.4,\y+0.35);
\node[font=\scriptsize,align=center,text=accent!70!black] at (5.5,\y)
  {\textbf{bande d'ignorance}\\largeur $1\,839$};
% resserrement
\draw[decorate,decoration={brace,amplitude=4pt},draw=navy] (7.4,\y+0.5) -- (5.1,\y+0.5)
  node[midway,above=4pt,font=\scriptsize,text=navy,align=center]{un registre spécifié\\en retire $66\,\%$};
\draw[densely dashed,draw=navy!70] (5.1,\y+0.36) -- (5.1,\y+0.5);
\node[align=left,font=\scriptsize,text=navy,anchor=west] at (7.9,\y+0.15)
  {$1\,024$ sommets\\énumérés};
\end{tikzpicture}
\end{center}

\noindent La largeur de la bande n'est pas une incertitude subie~: c'est une grandeur
\textbf{calculée à l'avance}, et le mémoire dit combien elle se réduirait si telle donnée
existait. C'est ce qui transforme une limite en livrable.

\section{Pourquoi $A$ compte quand même}
\label{sec:trace}

\lequoi{Si $A$ est invisible dans la donnée, pourquoi s'en soucier ?}

\laformule{Parce que le capital passe par $(I-W)^{-1}=I+W+W^2+\cdots$, et dès le terme
quadratique $W^2=S^2+SA+AS+A^2$ fait apparaître $A$. En prenant la trace~:
\[
  \operatorname{tr}(W^2)\;=\;\|S\|^2-\|A\|^2 .
\]
Le signe moins vient de l'antisymétrie~: $A^{\top}=-A$ implique $A^2=-AA^{\top}$.}

\laphrase{La direction est invisible au premier ordre mais elle apparaît dès qu'on met la
matrice au carré, et le calcul du capital passe par là. Une grandeur qui compte pour la
réponse et qui est absente de la donnée~: c'est la définition exacte d'une identification
partielle.}

% =====================================================================
\part*{VII. À quel point on y croit}
\addcontentsline{toc}{part}{VII. À quel point on y croit}

\section{Le bootstrap}

\lidee{On ne connaît pas la loi de l'estimateur. Alors on la fabrique~: on rééchantillonne
les données avec remise, on réestime à chaque fois, et la dispersion des estimations
obtenues approche la dispersion vraie.}

\cheznous{L'intervalle de confiance bootstrap du quantile de sévérité à $99{,}5\,\%$,
$[411\,;1\,037]$~M\euro, vaut un \textbf{facteur $2{,}5$} entre ses deux bouts. Le bootstrap est
aussi utilisé en version
\emph{paramétrique} pour les tests d'adéquation~: on simule sous la loi ajustée pour obtenir
la loi de la statistique de test.}

\laphrase{Le bootstrap fabrique la loi de mon estimateur en rééchantillonnant mes propres
données. Et le résultat est brutal~: rien que cette incertitude-là vaut un facteur deux et
demi sur le quantile d'un sinistre, et près de neuf sur la charge annuelle.}

\lepiege{Le bootstrap capture l'incertitude \emph{de paramètre} à modèle fixé. Il ne dit rien du
risque que le modèle lui-même soit faux. Et le facteur $2{,}5$ porte sur le quantile d'\emph{un}
sinistre : porté sur la charge annuelle, le même intervalle de $\xi$ donne un facteur $8{,}93$. Ne
jamais comparer tel quel ce $2{,}5$ à un facteur lu sur la charge annuelle.}

\section{Les tests d'adéquation}

\lidee{Une loi ajustée peut être fausse. On teste la distance entre la loi ajustée et les
données, et on regarde si cette distance est anormalement grande.}

\laformule{Kolmogorov-Smirnov mesure l'écart maximal entre les fonctions de répartition.
Anderson-Darling pondère ce même écart pour donner davantage de poids aux \emph{queues}, ce
qui est exactement la région qui nous intéresse.}

\cheznous{Par bootstrap paramétrique : Anderson-Darling $p=0{,}96$ et Kolmogorov-Smirnov
$p=0{,}89$ pour la famille GPD ; avec les paramètres publiés imposés, $0{,}99$ et $0{,}89$, ce
qui atteste le chiffre et pas seulement le choix de loi. La couverture réelle de l'intervalle à
$90\,\%$ de $\xi$ vaut $86\,\%$ : l'intervalle est \emph{trop étroit}, et il faudrait élargir sa
demi-largeur d'environ $10\,\%$. La binomiale négative bat Poisson avec $p\approx10^{-30}$.}

\laphrase{Les tests ne rejettent pas, avec des valeurs $p$ de $0{,}89$ à $0{,}99$. Attention à ce
que cela veut dire~: les \emph{formes} sont validées, mais c'est le \emph{niveau} qui reste une
bande.}

\lepiege{Les valeurs critiques tabulées de KS et AD sont fausses quand les paramètres sont
\emph{estimés} sur les mêmes données. Il faut passer par un bootstrap paramétrique, ce que
nous faisons, et il faut le dire, parce que beaucoup de mémoires ne le font pas.}

\section{Les trois étages d'incertitude}

\lidee{L'incertitude ne se réduit pas à l'intervalle de confiance. Elle a trois sources qui
se cumulent et dont les tailles n'ont rien à voir.}

\cheznous{
\begin{center}\small
\begin{tabular}{@{}l l r r@{}}
\toprule
\textbf{Étage, sur la charge annuelle} & \textbf{Ce qui varie} & \textbf{Bande (M\euro)} &
\textbf{Facteur} \\
\midrule
Paramètre & $\xi$ sur son intervalle à $90\,\%$ & $[3\,042\,;27\,152]$ & $8{,}93$ \\
Identification & la direction, non identifiée & $[6\,858\,;8\,697]$ & $1{,}27$ \\
Modèle, dépendance & de l'indépendance au comonotone & $[5\,322\,;9\,806]$ & $1{,}84$ \\
Modèle, famille de queue & de la lognormale à la queue lourde & $[4\,330\,;39\,901]$ & $9{,}21$ \\
\bottomrule
\end{tabular}
\end{center}}

\laphrase{Lues sur la même grandeur, la charge annuelle, l'incertitude de paramètre et celle de
modèle sont du même ordre, un facteur de l'ordre de neuf chacune, et elles écrasent celle de
l'identification. La bande que je reporte cumule le paramètre et l'identification ; le choix de
famille, qui ne se moyenne pas, est publié comme axe de sensibilité. Et le résultat qui tient est
ailleurs : \textbf{le niveau bouge dans la bande, l'ordre d'amorce des piliers ne bouge pas d'une
famille à l'autre}.}

\lepiege{Deux fautes à éviter. Comparer le facteur $2{,}5$ du quantile d'un sinistre au facteur
$5{,}5$ de la famille de queue (lognormale $4\,330$ contre GPD à $\xi=0{,}90$, $24\,016$) : ce sont
deux échelles. Et laisser croire que rien ne tient : toujours enchaîner avec les tests d'adéquation
et l'invariance de l'ordre. Les largeurs de paramètre font varier $\xi$ seul, ce sont des
majorants.}

\section{VaR prédictive et VaR robuste}

\lequoi{La remarque de Caroline Hillairet~: prendre un quantile d'un objet incertain est
mauvais. Comment y répondre sans abandonner la VaR, qui est imposée ?}

\lidee{Deux réponses différentes, à ne pas confondre.
\textbf{La VaR prédictive} intègre l'incertitude de paramètre dans la loi elle-même~: on
mélange les lois sur la loi a posteriori des paramètres, puis on prend le quantile du
mélange.
\textbf{La VaR robuste} prend le pire cas sur un ensemble d'ambiguïté~: on ne mélange pas,
on borne.}

\laformule{
\[
  \mathrm{VaR}^{\text{préd}}_\alpha = \mathrm{VaR}_\alpha\!
  \left(\int f(\cdot\mid\theta)\,\pi(\theta)\,d\theta\right),
  \qquad
  \mathrm{VaR}^{\text{rob}}_\alpha = \sup_{\theta\in\Theta}\ \mathrm{VaR}_\alpha
  \bigl(f(\cdot\mid\theta)\bigr).
\]}

\cheznous{Six postures, en M\euro{} : point $657$ ; prédictive $645\pm7$ (le $648$ d'un autre
calcul est le même nombre à la précision du bruit) ; robuste à $\beta=90\,\%$ $939\pm13$, à
$95\,\%$ $1\,037\pm9$, à $99\,\%$ $1\,247\pm24$ ; pire cas de modèle, $\xi$ poussé à $0{,}90$ :
$1\,331$. \textbf{La prédictive coïncide presque avec le point} à $99{,}5\,\%$, à deux fois le
bruit ; elle ne s'en écarte qu'en profondeur, $+4\,\%$ à $99{,}9\,\%$. Et la robuste à $95\,\%$
est exactement la borne haute de l'intervalle déjà publié, $1\,037$ : la même quantité sous deux
noms. Le mémoire reporte donc le point avec sa bande.}

\laphrase{Caroline a raison, et c'est mesurable. Ma réponse n'est pas d'abandonner la VaR,
qui est imposée, mais de rendre l'incertitude visible sur trois étages. Et j'ai un résultat
un peu contre-intuitif~: la VaR prédictive coïncide presque avec la VaR ponctuelle. Ce n'est
pas l'intégration qui manquait, c'est la largeur de bande qui est le vrai sujet.}

\lepiege{Ne pas présenter la prédictive comme \og la bonne réponse\fg{} qui remplacerait le
point. Elle règle un problème (l'incertitude de paramètre) et pas l'autre (l'ambiguïté de modèle).
C'est la robuste qui règle le second, et les deux ne se substituent pas. Et la posture la plus
prudente est la moins précise : $\pm24$ à $\beta=99\,\%$, quand les deux extrémités de l'échelle
n'ont aucun bruit.}

% =====================================================================
\part*{VIII. Les outils annexes}
\addcontentsline{toc}{part}{VIII. Les outils annexes}

\section{Le kappa de Cohen}

\lequoi{Deux codeurs lisent les mêmes récits. Comment mesurer leur accord, sachant que même
deux personnes répondant au hasard seraient parfois d'accord par chance ?}

\laformule{
\[
  \kappa = \frac{p_o - p_e}{1 - p_e},
\]
où $p_o$ est l'accord observé et $p_e$ l'accord attendu par pur hasard. $\kappa=1$ est un
accord parfait, $\kappa=0$ un accord au niveau du hasard.}

\cheznous{Le dispositif de second codage en aveugle est prêt : dix récits réduits aux faits,
anonymisés, et une grille de $100$ cellules dont $22$ sont directionnelles dans le codage de
référence, ce qui fixe la base d'accord (matériel complémentaire). Sans second jeu, le script ne
fabrique aucun coefficient. Par convention, un $\kappa$ au-delà de $0{,}6$ se lit comme un accord
substantiel ; en dessous, cela voudrait dire que la séquence des défaillances ne se lit pas de
façon reproductible dans les rapports d'enquête, ce qui serait un résultat en soi.}

\laphrase{Le kappa corrige l'accord du hasard : deux personnes qui répondraient au hasard seraient
parfois d'accord, et le kappa retranche cette part. Le second codage n'est pas encore reçu ; le
dispositif est prêt, et il ne produit rien sans lui.}

\section{La valeur d'information}

\lidee{Une recommandation de reporting ne vaut rien tant qu'elle n'est pas chiffrée. On la
chiffre en mesurant de combien la bande de capital rétrécit quand on acquiert
l'information.}

\laformule{Avec $L(\mathcal{V})=\sup_{\mathcal{V}}\mathrm{SCR}-\inf_{\mathcal{V}}\mathrm{SCR}$,
\[
  V(\mathcal{I}) = 1-\frac{L(\mathcal{W}_{\mathcal{I}})}{L(\mathcal{W})} \in[0,1].
\]
Cette quantité est \emph{monotone} (une information plus riche ne vaut pas moins) et
\emph{sans dimension} (donc comparable d'une source de sévérité à l'autre).}

\cheznous{Si l'on retenait le sens que les sept premiers rapports documentent sur six
dépendances, la largeur de la bande tomberait de $66\,\%$. Ce n'est pas un acquis : le mémoire
n'en revendique que deux paires, et ce chiffre est la valeur d'une information de cette étendue.
La valeur se \textbf{hiérarchise}~: documenter la seule dépendance gouvernance-incidents lèverait
$28{,}9\,\%$ de l'ambiguïté, quand le lien tests-prestataires n'en lèverait \emph{rien}.}

\laphrase{J'ai chiffré ce que vaut, en capital, une donnée qu'on n'a pas. Une information de sens
sur six dépendances resserrerait la bande des deux tiers. Et surtout, un superviseur qui ne
pourrait exiger qu'une seule chose sait laquelle exiger, parce que l'écart entre les dépendances
est énorme.}

\lepiege{Une bande subsisterait même les dix sens connus, imputable à la seule
\emph{amplitude}. L'identification partielle n'est pas un état provisoire que la donnée
finirait par lever, elle est structurelle.}

\section{Le bayésien conjugué}

\lidee{Plutôt que d'estimer un signe de direction par oui ou non, on lui attribue une
probabilité, qu'on met à jour avec les observations documentaires.}

\laformule{Modèle Beta-binomial~: si $\theta\sim\mathrm{Beta}(\alpha_0,\beta_0)$ est la
probabilité que la direction aille dans un sens, et qu'on observe $n_{lo}$ transitions dans
ce sens et $n_{hi}$ dans l'autre, alors
\[
  \theta \mid \text{données} \ \sim\ \mathrm{Beta}(\alpha_0+n_{lo},\ \beta_0+n_{hi}),
\]
et l'on pose $A = S\,(2\theta-1)$, de sorte que $\theta=1/2$ donne $A=0$.}

\cheznous{Le crédible à $90\,\%$ du capital est $81\,\%$ plus resserré que les bornes de pire
cas. Sous un prior sceptique, la médiane ne bouge presque pas : la conclusion vient des comptes
documentés. Mais seulement $2$~paires documentées sur $6$ ont un intervalle qui exclut $0{,}50$, et
les quatre paires que le corpus ne documente pas gardent un $\theta$ uniforme : le posterieur
reproduit de lui-même l'identification partielle.}

\laphrase{Le bayésien me permet de dire \og ce lien va probablement dans ce sens\fg{} au lieu de
trancher par oui ou non. Deux résultats~: le prior n'a presque aucune influence, ce qui est
rassurant, et seulement deux paires sur six sont réellement tranchées, ce qui est honnête. Je
publie les deux objets : les bornes, opposables car sans prior, et le crédible, informatif.}

\lepiege{Le bayésien ne crée pas d'information. Il exprime autrement la même incertitude, de
façon plus lisible et plus fine, mais il ne remplace pas la donnée manquante.}

% =====================================================================
\part*{IX. Du capital à la décision}
\addcontentsline{toc}{part}{IX. Du capital à la décision}

\section{À quelle échelle se lit un SCR}

\lequoi{Un capital de plusieurs milliards paraît absurde pour une compagnie d'assurance. C'est
la première objection qu'on vous fera, et une clause de style n'y répond pas.}

\lidee{Le niveau ne dépend pas de la cascade, il dépend de la \textbf{fréquence d'amorce}. On le
teste : on redescend $\lambda$ au niveau d'une entité, à sévérité, matrice $W$ et moteur
\textbf{inchangés}, et l'on regarde ce qui sort.}

\cheznous{\;\\[-6pt]
\begin{center}\small
\begin{tabular}{@{}l r r r r@{}}
\toprule
\textbf{Lecture} & $\lambda$/an & \textbf{Sévérité} & \textbf{SCR} & \textbf{Moyenne} \\
\midrule
Secteur financier mondial & $21{,}56$ & $\times1$ & $8\,123$~M\euro & $983$~M\euro \\
Seau \og{}au moins dix événements\fg{} & $0{,}210$ & $\times1$ & $488$~M\euro & $9{,}7$~M\euro \\
Seau \og{}toutes firmes\fg{} & $0{,}100$ & $\times1$ & $229$~M\euro & $4{,}6$~M\euro \\
$\lambda$ lu à la taille de l'entité & $0{,}092$ & $\times1$ & $198$~M\euro & $4{,}3$~M\euro \\
\textbf{Les deux canaux composés} & $0{,}092$ & $\times0{,}854$ & $\mathbf{169}$~M\euro & $3{,}6$~M\euro \\
\bottomrule
\end{tabular}
\end{center}
La dernière ligne est la lecture cohérente : fréquence \emph{et} sévérité prises à la taille de
l'entité. Sa bande sur les $1\,024$ sommets va de $131{,}5$ à $183{,}3$~M\euro.}

\laphrase{Mes montants sont un besoin de capital du \emph{secteur financier mondial}, pas d'une
entreprise. Redescendu à l'échelle d'une entité, fréquence et sévérité lues à sa taille, le même
modèle donne $169$~millions, que le mémoire publie comme un majorant d'ordre de grandeur, l'entité
étant hors du domaine de validité. Attention au piège que j'ai moi-même tendu : au seau des firmes à
au moins dix événements, je trouvais $488$, tout près du repère forfaitaire de $450$, et je croyais
l'écart refermé ; la coïncidence ne survit pas à la lecture à la taille. Ce qui parle, ce n'est pas
le niveau, c'est que le forfait ne bouge pas d'un euro entre un profil exemplaire et un profil
défaillant.}

\lepiege{La non-linéarité est le vrai résultat. Diviser $\lambda$ par $235$ divise la moyenne par
$270$, donc à proportion, mais ne divise le capital que par $48$ : c'est le \textbf{principe du
grand saut unique}, le quantile d'une entité décrivant un sinistre isolé et non une accumulation.
Diviser le capital du secteur par le nombre d'entités le sous-estimerait d'un facteur de l'ordre de
cinq. Et le repère de $3\,\%$ des provisions n'est qu'un ordre de grandeur : il s'écarte de
$-8\,\%$ et $+56\,\%$ des modules publiés dans les SFCR.}

\section{Détenir ou transférer}

\lequoi{Le mémoire compare son SCR à la Formule Standard, aux pertes réelles et au marché
cyber. Trois comparaisons de \emph{niveau}. Il en manque une de \emph{prix}.}

\lidee{Un SCR n'est pas une perte, c'est du capital \textbf{immobilisé}, et le capital coûte.
Solvabilité~II chiffre ce coût explicitement. En face, transférer coûte une prime. Les deux
grandeurs sont homogènes, donc comparables, et c'est l'arbitrage que fait un directeur des
risques.}

\laformule{Le coût annuel de détention vaut $\mathrm{CoC}\times\mathrm{SCR}$, où
$\mathrm{CoC}$ est le taux de coût du capital de la marge de risque : $6\,\%$ historiquement,
ramené à $\mathbf{4{,}75\,\%}$ par la directive (UE)~2025/2. Pour un traité en excès de perte
annuel de portée $L$, on minimise
\[
  \mathrm{co\hat ut}(L)=\underbrace{\E[\min(X,L)]/0{,}485}_{\text{prime}}
  \;+\;\mathrm{CoC}\times\underbrace{\mathrm{VaR}_{99{,}5\%}\bigl(X-\min(X,L)\bigr)}_{\text{capital résiduel}} .
\]}

\cheznous{Multiple de capital $\mathrm{SCR}/\E[X]=\mathbf{47}$ à l'échelle entité contre $8{,}3$
au secteur. Coût de détention $8{,}0$~M\euro/an au taux révisé ($10{,}1$ à $6\,\%$) contre
$3{,}6$~M\euro{} de sinistralité attendue, soit $2{,}2$ à $2{,}8$ fois. Le capital résiduel vaut
$\max(\mathrm{SCR}-L,0)$, donc l'optimum tombe exactement en $L^\star=\mathrm{SCR}=169$~M\euro, et
le coût total passe à $\mathbf{2{,}9}$~M\euro, soit $\mathbf{-64\,\%}$. Une entité devenue conforme
voit en outre ce coût baisser de $81\,\%$ : la conformité déplace le point d'équilibre entre
rétention et transfert.}

\laphrase{Immobiliser le capital coûte deux fois et demie ce que le risque coûte en moyenne.
Près de l'optimum, un euro de portée coûte un peu plus d'un pour cent de prime et libère
quatre virgule sept cinq pour cent de coût du capital. C'est pour ça que le transfert domine, et
c'est un calcul, pas une opinion.}

\lepiege{Ce gain est une \textbf{borne supérieure}, pour trois raisons : la capacité du marché
plafonne près de $500$~M\euro{} sur un seul risque ; les polices excluent la guerre et les
infrastructures essentielles ; et sous Solvabilité~II céder ne supprime pas l'exigence, cela
\emph{substitue} au risque cyber un risque de \textbf{défaut de contrepartie}. En résumant les
trois par un coefficient d'inefficacité $\kappa$, le seuil critique vaut
$\kappa^\star=1-0{,}0103/0{,}0475=\mathbf{78\,\%}$ : il faudrait que plus des trois quarts de la
couverture soient inefficaces pour renverser la conclusion. La conclusion est robuste, le
chiffre ne l'est pas.}

% =====================================================================
\newpage
\part*{X. Chapitre par chapitre : ce qu'il faut savoir par c\oe ur}
\addcontentsline{toc}{part}{X. Chapitre par chapitre : ce qu'il faut savoir par c\oe ur}

\begin{cle}
\textbf{Comment lire cette partie.} Une section par chapitre du mémoire, dans l'ordre du
document. Chacune donne le chapitre en une phrase, puis un encadré \emph{à savoir par c\oe ur}
(les nombres, définitions et formules à restituer sans notes), puis les questions que le jury
peut poser sur ce chapitre, avec la réponse à tenir et l'endroit du mémoire où s'appuyer. Les
questions qui ne relèvent d'aucun chapitre (forme, déontologie, dispositif) sont regroupées à
la fin.

\textbf{Les pages} sont celles du fichier déposé à l'ENSAE, \texttt{KADDOURI\_Kelian\_MS25.pdf},
$199$~pages. La version de l'Institut, $195$~pages, a exactement les mêmes pages jusqu'à la
conclusion (p.~$149$). Elle n'a pas le chapitre du stage, donc ses annexes arrivent quatre pages
plus tôt : annexe~A p.~$152$ au lieu de $156$, annexe~C p.~$169$ au lieu de $173$, annexe~G
p.~$187$ au lieu de $191$.

\textbf{Trois règles pour l'oral}, tirées du \S~6.7 des Recommandations Jury : la présentation
dure \textbf{25 minutes au plus} et le dépassement est pénalisant ; le jury ne cherche pas à
déstabiliser mais à mesurer la maturité ; et \emph{si vous ne savez pas, dites-le}.
\end{cle}

\begin{cle}
\textbf{Quatre réflexes qui évitent les fautes graves.} (i)~Dire \og au secteur\fg{} avant tout
montant en milliers de millions. (ii)~Dire \og besoin de capital ORSA au titre de DORA\fg{},
jamais \og le SCR DORA de telle entité\fg{}. (iii)~Ne jamais additionner les quatre canaux, ni
l'interaction entre canaux avec celle entre piliers. (iv)~Ne jamais comparer un quantile de
sévérité d'un sinistre ($663$~M\euro) à un capital ($6\,049$ ou $8\,123$~M\euro) : ce ne sont
pas des objets de même échelle.
\end{cle}

% ---------------------------------------------------------------------
\section*{Les pages liminaires : résumé, abstract et notes de synthèse}
\addcontentsline{toc}{section}{Les pages liminaires : résumé, abstract et notes de synthèse}
\reperes{Résumé p.~3, abstract p.~4, glossaire p.~10, note de synthèse p.~12, executive summary
p.~16.}

\enbref{Ces pages sont lues en premier, et parfois seules : chaque nombre qu'elles portent doit
sortir sans hésitation, et dans la même formulation que le résumé.}

\begin{parcoeur}
\begin{fichetab}
UE 2022/2554 & le règlement DORA, applicable depuis le 17 janvier 2025 \\
$1\,024$ sommets & les $2^{10}$ orientations extrêmes des dix paires de piliers \\
$[6\,858\,;8\,697]$ M\euro & le capital borné sur l'ensemble admissible ; la largeur,
  $1\,839$, est le coût de l'ignorance de la direction, connu d'avance \\
$6\,049\to20\,188$ M\euro & conforme vers non conforme, quatre canaux, au secteur :
  facteur $3{,}34$, écart $14\,139$ \\
$9\,138$ et $5\,001$ & la somme des canaux isolés, puis leur interaction ($35\,\%$ de
  l'écart) \\
$10\,377$ et $11\,413$ & la fréquence seule, puis les trois autres canaux réunis \\
$131{,}5$ à $183{,}3$ M\euro & l'entité notionnelle, majorant d'ordre de grandeur ($169$ au
  point) \\
facteur $22$, $41$ & la variation du niveau avec la seule source de sévérité, puis avec la
  fréquence triplée en plus \\
$66\,\%$ et $28{,}9\,\%$ & le resserrement de la bande si le sens de six dépendances était
  établi, puis par la seule dépendance gouvernance-incidents \\
$2{,}78$ & le coût annuel de détention rapporté au prix de marché du transfert (taux révisé) \\
$1\,811$ années & l'historique nécessaire pour backtester le quantile à $99{,}5\,\%$ \\
\end{fichetab}
\end{parcoeur}

\begin{question}{\og Résumez votre mémoire en une minute.\fg{}}
\rep{Quatre phrases, dans l'ordre du résumé. Je chiffre ce que coûte en capital la
non-conformité à DORA, avec une cascade dirigée entre les cinq piliers plutôt qu'une copule. La
direction de cette cascade n'est pas identifiable sur données publiques alors que sa
co-occurrence l'est : je borne donc le capital sur les $1\,024$ configurations admissibles, entre
$6\,858$ et $8\,697$~millions, et la largeur est le coût de l'ignorance. Sur les quatre canaux que
la conformité déplace, le besoin de capital du secteur passe de $6\,049$ à $20\,188$~millions,
un facteur $3{,}34$, et c'est la fréquence des incidents qui commande. Enfin la limite devient
un livrable : je dis quelle donnée réclamer et ce qu'elle rapporterait.}
\appui{résumé p.~3.}
\end{question}

\begin{question}{\og La note parle d'un facteur $22$ sur le niveau, le résumé d'un facteur $41$. Lequel ?\fg{}}
\rep{Les deux : ce ne sont pas les mêmes variations. $22$ en changeant la seule source de
sévérité, OpRisk contre PRC ; $41$ en triplant en plus la fréquence. Les deux disent la même
chose : le niveau est fixé par des entrées dont le périmètre de collecte n'est pas transposable,
alors que les rapports au socle varient d'au plus $33\,\%$.}
\appui{chapitre~9, \S~9.4 p.~107 ; note de synthèse, rubrique Limites.}
\end{question}

\begin{question}{\og Deux canaux reposent sur des bornes posées, écrivez-vous. Lesquels, et qu'est-ce que cela change ?\fg{}}
\rep{La propagation ($g$ de $0{,}45$ à $0{,}90$) et l'accumulation par un prestataire commun
(de $0$ à $0{,}68$). La fréquence et la détection sont calibrables. Pour les deux canaux posés,
l'effet de fermeture est une borne et non un gain de remédiation assuré, et leur part de Shapley
n'est pas un budget.}
\appui{chapitre~7, \S~7.3 p.~86 ; annexe~B, \S~B.4 p.~166.}
\end{question}

% ---------------------------------------------------------------------
\section*{Chapitre 1. Risque cyber, résilience numérique et besoin de capital}
\addcontentsline{toc}{section}{Chapitre 1. Risque cyber, résilience numérique et besoin de capital}
\reperes{p.~21 à 35. 1.1 le risque (p.~21), 1.2 l'assurabilité (p.~24), 1.3 le marché français
(p.~26), 1.4 DORA (p.~28), 1.5 l'enjeu prudentiel et la problématique (p.~31).}

\enbref{Le cyber n'est pas un risque de plus : sa charge est portée par quelques événements, ses
causes se propagent par les prestataires partagés, et le marché français n'a pas l'historique
pour en estimer la queue. Ces trois traits commandent la théorie des valeurs extrêmes, la cascade
dirigée et la calibration sur une base internationale.}

\begin{parcoeur}
\begin{fichetab}
Incident, attaque & l'incident TIC est tout événement qui affecte la sécurité ; l'attaque exige
  trois conditions cumulatives (intention, impact significatif, aboutissement). DORA retient
  l'incident : la panne CrowdStrike de 2024 est dans son champ \\
Quatre familles & rançongiciel, fuite de données, déni de service, défaillance d'un prestataire,
  la seule non individualisable, qui produit des sinistres corrélés \\
Quatre cas & WannaCry 2017 (échelle), NotPetya 2017 (P4 puis P1 puis P2), MOVEit 2023
  (accumulation), CrowdStrike 2024 (dépendance sans adversaire, P4 et P3) \\
LUCY 2026 & $20\,996$ polices, $1\,251$ sinistres, sept exercices ; assurés $+49\,\%$, primes
  $-3\,\%$ ; sinistres sur primes $12\,\%$ en 2023, $27\,\%$ en 2025, $169\,\%$ en 2020 ; charge
  indemnisée $83{,}2$ M\euro{} \\
Deux régimes & 2024, sévérité (nombre $\times0{,}73$, coût moyen $\times1{,}96$) ; 2025,
  fréquence (nombre $\times2{,}79$, coût moyen $\times0{,}55$) ; un seul sinistre au-delà de
  $10$ M\euro{} en 2025 \\
Cinq piliers & P1 gouvernance (art.~5 à 16), P2 incidents (17 à 23), P3 tests (24 à 27, TLPT
  de l'art.~26 tous les trois ans), P4 tiers (28 à 44, clauses de l'art.~30, troisième rang de
  sous-traitance), P5 partage (art.~45) \\
Notification & alerte sous 4~h après classification et au plus 24~h après détection, rapport
  intermédiaire sous 72~h, rapport final sous un mois \\
AMF 2025 & $47$ notifications, $31$ incidents majeurs, $27$ d'origine tierce ($87\,\%$) : une
  part d'origine, pas une probabilité de transfert \\
CESIN, France Ass. & attaque par un tiers $35\,\%$, $43\,\%$ dans les grandes entreprises ;
  sanction d'une autorité dans $3\,\%$ des cas ; cyber premier risque pour la huitième année
  ($4{,}1$ sur $5$), fréquence $-0{,}45$ et sévérité $+0{,}25$ en 2026 \\
\end{fichetab}
\end{parcoeur}

\begin{question}{\og Pourquoi DORA, et pas NIS~2, le RGPD ou le règlement sur la cyberrésilience ?\fg{}}
\rep{Parce que c'est le seul de ces textes qui découpe la résilience en domaines de contrôle dont
la défaillance se lit dans des données d'incidents, et le seul dont la non-conformité s'adresse à
l'entité qui porte le capital. Les autres sont du contexte : NIS~2 et le RGPD visent la sécurité
et la confidentialité, la directive sur les entités critiques la criticité sectorielle, le
règlement sur la cyberrésilience le produit mis sur le marché.}
\appui{\S~1.4, p.~29.}
\end{question}

\begin{question}{\og Vous êtes co-auteur de l'étude LUCY. N'est-ce pas gênant de la citer ?\fg{}}
\rep{Elle est citée comme citation externe non recalculable : aucune de ses valeurs n'entre dans
un calcul ni ne porte un capital. Ce qui est vérifiable l'est : un script contrôle la fidélité de
la recopie, sinistres sur primes, somme des classes de taille sur sept exercices, multiplicateurs.
Et trois de ses constats servent le mémoire sans le flatter : l'absence de traduction
prudentielle, la queue nationale vide, et la décomposition entre nombre et coût moyen.}
\appui{\S~1.3, p.~26.}
\end{question}

\begin{question}{\og Vos milliers de millions contre les $83$~millions indemnisés du marché français : lequel est faux ?\fg{}}
\rep{Aucun, ils ne mesurent pas la même chose. LUCY publie une charge \emph{indemnisée}, nette de
franchise et plafonnée par la capacité, sur le marché français. Le mémoire modélise la perte
\emph{brute} d'entités financières, à l'échelle du secteur mondial. Les rapprocher ferait lire
une différence de définition comme une contradiction.}
\appui{\S~1.3, encadré p.~26.}
\end{question}

\begin{question}{\og Le marché voit presque tripler le nombre de sinistres, les enquêtes voient la fréquence baisser. Qui croire ?\fg{}}
\rep{Les deux. Le nombre de sinistres du marché est un total, porté par une base d'assurés
élargie de $49\,\%$ ; le CESIN et la cartographie mesurent une exposition par entité. On peut
compter plus de sinistres en assurant des entités chacune moins exposées. C'est pour cela que le
mémoire parle de \emph{nombre} et non de fréquence pour les multiplicateurs de LUCY.}
\appui{\S~1.5, encadré p.~34.}
\end{question}

\begin{question}{\og Le vrai coût de la non-conformité, n'est-ce pas l'amende ?\fg{}}
\rep{Non, et la donnée le dit : une sanction d'autorité n'apparaît que dans $3\,\%$ des cas du
CESIN, loin derrière la perturbation de la production ($28\,\%$) et la perte d'image ($26\,\%$).
Le modèle n'ajoute d'ailleurs aucune pénalité : la non-conformité déplace quatre paramètres de la
loi de perte. Un modèle qui ferait de l'amende le coût principal se tromperait de mécanisme.}
\appui{\S~1.5, p.~33 ; chapitre~7, \S~7.3.}
\end{question}

\begin{question}{\og L'AMF trouve $87\,\%$ d'incidents d'origine tierce. Pourquoi ne pas calibrer P4 dessus ?\fg{}}
\rep{Parce que c'est une part d'\emph{origine}, pas une probabilité de transfert d'un domaine de
contrôle à un autre. Elle confirme l'accumulation du pilier des tiers sur une population
entière, un prestataire compromis touchant plusieurs sociétés, mais ne dit rien de la force d'un
lien de P4 vers P2. Et l'écart entre $47$ notifications et $31$ incidents confirmés annonce ce que
sera le registre : une collecte au-dessus d'un seuil de matérialité, pas un recensement.}
\appui{\S~1.4, p.~30.}
\end{question}

\begin{question}{\og Pourquoi la Formule Standard ne suffit-elle pas ?\fg{}}
\rep{Elle charge le risque opérationnel par un pourcentage de primes ou de provisions : deux
entités de bilan identique et de maturité opposée sur les cinq piliers paient exactement la même
charge. Aucun sous-module cyber n'existe. Une copule ne comble pas le vide, parce qu'elle est
symétrique : elle ne distingue pas une défaillance qui entraîne les autres d'une défaillance qui
les subit.}
\appui{\S~1.5, p.~34.}
\end{question}

\lepiege{Les tentatives WannaCry bloquées (plus de $64$~millions en Russie) ne sont pas des
sinistres : elles mesurent la persistance d'une exposition, et ne se comparent à aucune grandeur
du modèle.}

% ---------------------------------------------------------------------
\section*{Chapitre 2. Le cadre prudentiel de capital et son articulation avec DORA}
\addcontentsline{toc}{section}{Chapitre 2. Le cadre prudentiel et son articulation avec DORA}
\reperes{p.~36 et 37.}

\enbref{Le risque opérationnel a la même définition en banque et en assurance mais pas le même
traitement : Bâle a retiré les modèles internes, Solvabilité~II garde la liberté de méthode, et
c'est dans l'ORSA, non dans le forfait du pilier~1, que se loge le mémoire.}

\begin{parcoeur}
\begin{fichetab}
Risque opérationnel & perte due à des procédures, des personnes ou des systèmes inadaptés ou
  défaillants, ou à des événements extérieurs ; inclut le juridique, exclut le stratégique et la
  réputation \\
Formule Standard & $\mathrm{SCR}=\mathrm{BSCR}+\mathrm{Adj}+\mathrm{SCR}_{Op}$, avec
  $\mathrm{SCR}_{Op}=\min(0{,}3\,\mathrm{BSCR},\,\mathrm{Op})+0{,}25\,\mathrm{Exp}_{ul}$ et
  $\mathrm{Op}=\max(\mathrm{Op}_{\text{primes}},\mathrm{Op}_{\text{provisions}})$ ; article~204
  du règlement délégué 2015/35 ; terme de croissance au-delà de $20\,\%$ de hausse des primes \\
Bâle & modèles internes AMA retirés en 2022 au profit de l'approche standardisée ; l'approche
  par distribution des pertes reste un outil de modélisation \\
Solvabilité~II & Formule Standard, paramètres propres (USP), ou modèle interne partiel ou
  global approuvé \\
Piliers 1 et 2 & le forfait du pilier~1, plafonné à $30\,\%$ du BSCR, contre l'ORSA du
  pilier~2, prospectif, qui admet les scénarios de stress stochastiques \\
Additivité & $\mathrm{SCR}_{Op}$ n'entre pas dans la matrice de corrélation : un euro d'écart
  DORA vaut un euro de capital (\S~10.12) \\
\end{fichetab}
\end{parcoeur}

\begin{question}{\og Écrivez la charge de risque opérationnel de la Formule Standard.\fg{}}
\rep{$\mathrm{SCR}=\mathrm{BSCR}+\mathrm{Adj}+\mathrm{SCR}_{Op}$, et
$\mathrm{SCR}_{Op}=\min(0{,}3\,\mathrm{BSCR},\mathrm{Op})+0{,}25\,\mathrm{Exp}_{ul}$, où
$\mathrm{Op}$ est le maximum d'une branche primes et d'une branche provisions, à coefficients
forfaitaires. Elle ne dépend que du volume : sa réponse à la conformité est nulle par
construction.}
\appui{chapitre~2, p.~36.}
\end{question}

\begin{question}{\og Pourquoi l'ORSA, et pas un modèle interne de pilier~1 ?\fg{}}
\rep{Parce qu'aucun module DORA n'existe dans la Formule Standard et qu'un modèle interne
approuvé exigerait une validation et une donnée d'entité que personne n'a. L'ORSA demande une
évaluation prospective de l'adéquation entre le profil de risque réel et les besoins de
solvabilité, et il admet les scénarios stochastiques. La grandeur calculée est donc un besoin de
capital de pilier~2, pas un SCR.}
\appui{chapitre~2, p.~37.}
\end{question}

\begin{question}{\og Bâle a abandonné les modèles internes. Pourquoi y revenir ?\fg{}}
\rep{Bâle a retiré l'AMA pour le calcul réglementaire, sans invalider l'approche par distribution
des pertes comme outil de modélisation. Le mémoire en reprend les deux marges, fréquence et
sévérité, sur la base même qui alimentait ces modèles, et remplace leur agrégation par la
cascade. L'assurance, elle, garde sous Solvabilité~II la liberté de méthode que la banque a
perdue.}
\appui{chapitre~2, p.~36 ; chapitre~4, \S~4.7 p.~51.}
\end{question}

\begin{question}{\og Pourquoi un euro de DORA vaudrait-il plus qu'un euro de risque diversifiable ?\fg{}}
\rep{Parce que la charge opérationnelle s'ajoute au BSCR hors matrice de corrélation : un écart
est reconnu en entier. Agrégé comme un module corrélé à $\rho$, il ne le serait qu'à hauteur de
$\tau\approx\rho$ ; à $\rho=0{,}25$, l'écart de traitement va de $3{,}4$ à $3{,}9$ sur les quatre
entités réelles. Cette conclusion ne dépend d'aucun paramètre de la cascade, seulement de la
place que le règlement donne au risque opérationnel.}
\appui{chapitre~10, \S~10.12 p.~126.}
\end{question}

\lepiege{Ne pas dire que la Formule Standard sous-estime le cyber : elle ne le voit pas. Et le
repère de $3\,\%$ des provisions employé au chapitre~10 n'est qu'un ordre de grandeur, qui
s'écarte de $-8\,\%$ et $+56\,\%$ des modules publiés.}

% ---------------------------------------------------------------------
\section*{Chapitre 3. État de l'art et positionnement}
\addcontentsline{toc}{section}{Chapitre 3. État de l'art et positionnement}
\reperes{p.~38 à 42. 3.1 la signature (p.~38), 3.2 le positionnement (p.~39), 3.3 la revue par
thème (p.~40), 3.4 où se situe le mémoire (p.~42).}

\enbref{Trois vagues de littérature (assurabilité et valeurs extrêmes, contagion entre entités,
Hawkes dans le temps) ; le mémoire en garde les marges et ajoute un mécanisme qu'aucune ne porte,
une dépendance dirigée et sensible à l'ordre entre les domaines de contrôle d'une entité,
traduite en capital.}

\begin{parcoeur}
\begin{fichetab}
Signature & sur les $26$ sous-ensembles d'au moins deux piliers, $22$ changent de criticité
  selon l'ordre ; écart jusqu'à $6$ points sur $10$ ; la chaîne P1 vers P2 vaut $8$, P2 vers P1
  vaut $2$ ($\mathrm{TRANS}_{12}=0{,}80$, $\mathrm{TRANS}_{21}=0{,}20$) \\
Les concurrents & AMDEC (ni direction ni propagation), copule (symétrique), réseau et
  épidémiologie (entre entités, exige la structure du réseau), Hawkes (dans le temps, entre types
  d'attaque, exige un horodatage fin), réseau bayésien (acyclique, tables conditionnelles au même
  mur d'identification) \\
Emprunts & au crédit, le branchement multitype et les modèles multi-états de migration ; à
  Leontief, la normalisation ; à Hillairet et Lopez, les états de conformité ; à Manski (1990,
  1995) et Tamer (2010), l'identification partielle \\
Précédents & Eling et Wirfs (2019), sur la même base SAS OpRisk, pour la sévérité ; Ranguin
  (2024), mémoire pour la CCR sur des sources proches, sans capital à un quantile réglementaire \\
Préprint climatique & entièrement synthétique : on lui emprunte une architecture, aucun chiffre \\
Non revendiqué & la non-transitivité, testée dans trois formulations, aucune ne tient \\
\end{fichetab}
\end{parcoeur}

\begin{question}{\og Qu'est-ce qui est vraiment nouveau dans votre travail ?\fg{}}
\rep{Le mécanisme de dépendance : dirigé et sensible à l'ordre entre piliers réglementaires d'une
même entité, traduit en capital, avec sa frontière d'identification démontrée puis convertie en
spécification de reporting. Il n'y a ni loi de sévérité ni processus de fréquence nouveaux : sur
les marges, je calibre la littérature. Neuf sur le mécanisme, sobre sur le reste, déclaré sur
les limites.}
\appui{\S~3.4, encadré p.~42.}
\end{question}

\begin{question}{\og Pourquoi pas un réseau bayésien, qui porte aussi une direction ?\fg{}}
\rep{Il suppose un graphe acyclique, alors que la contagion entre piliers a des cycles, que la
forme de Leontief gère. Et ses tables de probabilités conditionnelles se heurtent exactement au
mur d'identification de $W$ : il déplacerait le problème sans le résoudre.}
\appui{\S~3.2, p.~40.}
\end{question}

\begin{question}{\og Quel est le travail le plus proche du vôtre ?\fg{}}
\rep{Pour la sévérité, Eling et Wirfs (2019), qui étudient les pertes cyber par valeurs extrêmes
sur la même base SAS OpRisk. Pour les sources, le mémoire de Ranguin (2024) : son objet est la
sinistralité d'un portefeuille assuré, sans mesure de capital à un quantile réglementaire ni
dépendance entre domaines de contrôle.}
\appui{\S~3.3, p.~40.}
\end{question}

\begin{question}{\og Qu'empruntez-vous aux travaux de Caroline Hillairet ?\fg{}}
\rep{La logique d'états de conformité de Hillairet et Lopez, appliquée entre les piliers d'une
entité et non entre assurés. Et l'écartement du processus de Hawkes est testé contre ses propres
noyaux, ceux de Bessy-Roland, Boumezoued et Hillairet et de Boumezoued, Cherkaoui et Hillairet,
pas contre un noyau de paille.}
\appui{\S~3.2 et 3.3 ; chapitre~11, \S~11.3.}
\end{question}

\begin{question}{\og Ce préprint climatique qui confirme vos résultats, que vaut-il ?\fg{}}
\rep{Il est synthétique et le dit : aucun de ses chiffres n'est un repère empirique ni n'est
repris. J'emprunte une architecture, la séparation des échelles de temps, le protocole à marges
appariées, l'ablation. Là où nos classements divergent, la cause est dans son texte : sa sévérité
est bornée, donc sans brique de queue, alors que la queue est la brique la plus lourde ici.}
\appui{\S~3.3, encadré p.~41 ; chapitre~10, \S~10.7.}
\end{question}

% ---------------------------------------------------------------------
\section*{Chapitre 4. Données et limites}
\addcontentsline{toc}{section}{Chapitre 4. Données et limites}
\reperes{p.~44 à 63. 4.1 cartographie (p.~44), 4.2 trois biais (p.~45), 4.3 actualisation
(p.~46), 4.4 trois fréquences (p.~48), 4.5 deux biais chiffrés (p.~49), 4.6 Hackmageddon (p.~50),
4.7 OpRisk (p.~51), 4.8 profilage des variables (p.~53), 4.9 portée (p.~62), 4.10 pourquoi pas
plus de sources (p.~63).}

\enbref{Aucune base ne dit tout : chaque source porte ce qu'elle peut porter, ses biais sont
déclarés puis mesurés quand c'est possible, et la queue commune est une décision adossée à des
tailles d'effet, pas une commodité.}

\begin{parcoeur}[À savoir par c\oe ur : les sources et les fréquences]
\begin{fichetab}
PRC & $74\,783$ incidents dans la chronologie, $15\,053$ retenus sur 2019--2025 ; sévérité
  \emph{dérivée} par Jacobs, $\ln L=7{,}68+0{,}76\ln X$ ; $\hat\xi_{\text{PRC}}=1{,}03$
  $[0{,}97\,;1{,}09]$ \\
SAS OpRisk & $41\,975$ pertes ; $582$ cyber $\times$ finance sur 2000--2026 pour la GPD ;
  $1\,476$ TIC en finance sur 2005--2022 pour la fréquence et la descente ; le décile supérieur
  porte $92{,}3\,\%$ de la perte, le centile $65{,}1\,\%$ \\
Hackmageddon & $1\,041$ incidents de janvier à juin 2026, $840$ à vecteur identifié ($81\,\%$) ;
  hameçonnage $38{,}8\,\%$, exploitation $33{,}8\,\%$, chaîne d'approvisionnement $15{,}8\,\%$,
  identifiants $6{,}3\,\%$ ; citation, non recalculable \\
Trois fréquences & $\lambda_{ref}=341$/an (notifications PRC des organisations financières,
  code BSF ; $2\,150$ sur la base entière) ; $\lambda=21{,}6$/an (incidents TIC matériels du
  secteur, seule entrée du moteur) ; $\lambda_{\text{entité}}=0{,}092$/an (lu à la taille de
  l'entité ; les seaux valent $0{,}10$ et $0{,}21$) ; facteur $3\,719$ entre les extrêmes \\
Montants & nominaux ; inflation implicite $2{,}19\,\%$ par an (rapport $1{,}747$ en 2004), non
  appliquée à la calibration gelée \\
\end{fichetab}
\end{parcoeur}

\begin{parcoeur}[À savoir par c\oe ur : les biais et le profilage]
\begin{fichetab}
Trois biais déclarés & sélection (OpRisk, grandes entités), déclaration (PRC, les fuites plus
  que les rançongiciels sans exfiltration), survie (la ruine sort de la base : queue allégée) \\
Deux biais mesurés & élasticité sévérité-taille corrigée de la troncature $b=0{,}087$
  $[0{,}026\,;0{,}148]$, soit une sévérité multipliée par $0{,}85$ à l'entité ; chain-ladder :
  $97$ événements sur 2023--2025 deviennent $255$ à l'ultime \\
Compensation & $0{,}861\times1{,}147=0{,}987$ sur la perte moyenne ; pas sur le quantile \\
Complétude & un manquant PRC vaut la chaîne \texttt{UNKN} : un contrôle naïf donne $100\,\%$ ;
  le volume touché n'est renseigné qu'à $41{,}77\,\%$, et les lacunes se cumulent \\
Doublons & $1\,643$ lignes, $10{,}9\,\%$, borne basse ; le mémoire n'additionne jamais ces
  volumes \\
Queue commune & $V$ de Cramér $<0{,}20$ (max $0{,}176$), $\varepsilon^2<0{,}08$ (max
  $0{,}0785$), Spearman $<0{,}16$ (actif total $0{,}068$, dernier des trois) \\
Seuil & $20{,}03$ M\euro{}, percentile $84{,}4$ de la donnée courante ($q_{85}=22{,}03$) ;
  exactement $91$ excès ; maximum $=3{,}7\times q_{99}$ contre $52{,}7$ sur la base entière \\
\end{fichetab}
\end{parcoeur}

\begin{question}{\og Pourquoi calibrer une entité française sur des pertes de grandes banques mondiales ?\fg{}}
\rep{Parce que la base française est vide là où se joue le quantile : un seul sinistre au-delà de
$10$~M\euro{} en 2025 dans LUCY, et la classe la plus haute vide deux exercices sur sept. Le prix
est un biais de taille, que je mesure au lieu de le supposer : $b=0{,}087$, la sévérité transposée
à l'entité n'est multipliée que par $0{,}85$. C'est aussi pourquoi le chiffre d'entité est publié
comme un majorant.}
\appui{\S~4.5, p.~49 ; chapitre~1, \S~1.3.}
\end{question}

\begin{question}{\og Votre fréquence vient d'une base américaine sans seuil de matérialité.\fg{}}
\rep{C'est le cas de $\lambda_{ref}=341$, et c'est pour cela qu'il n'est qu'une clé de
répartition par vecteur d'attaque : il ne fixe aucun niveau. Le moteur prend $\lambda=21{,}6$,
les incidents TIC \emph{matériels} lus sur OpRisk, au-dessus d'un seuil de collecte.}
\appui{\S~4.4, encadré p.~49 ; chapitre~5, \S~5.3.}
\end{question}

\begin{question}{\og Trois nombres s'appellent fréquence et vont de $0{,}092$ à $341$. Lequel est le bon ?\fg{}}
\rep{Les trois, à trois niveaux d'observation : des notifications, des pertes matérielles du
secteur, les mêmes pertes rapportées à une firme de la taille visée. Le facteur $3\,719$ entre les
extrêmes est recalculé à chaque exécution. Un seul alimente le moteur, $\lambda=21{,}6$.}
\appui{\S~4.4, encadré p.~49.}
\end{question}

\begin{question}{\og Comment avez-vous traité les biais de vos bases ?\fg{}}
\rep{Trois sont déclarés avant usage, sélection, déclaration et survie, et deux sont mesurés.
Les deux corrections mesurées vont en sens opposés : $0{,}861\times1{,}147=0{,}987$, un effet net
de $1{,}3\,\%$ sur la perte moyenne. Mais elles ne se compensent pas sur le capital : à queue
lourde, le quantile suit la sévérité, et la correction de taille porte sur la sévérité.}
\appui{\S~4.2, p.~45, et \S~4.5, p.~49.}
\end{question}

\begin{question}{\og Pourquoi une seule queue pour toutes les catégories ?\fg{}}
\rep{Parce que trois instruments indépendants disent que la catégorie informe peu la sévérité :
$V$ de Cramér sous $0{,}20$, $\varepsilon^2$ sous $0{,}08$, corrélations de rang sous $0{,}16$.
Segmenter par les sept catégories bâloises ramènerait les $91$ excès à une douzaine par modalité,
sur lesquels aucune loi de queue ne s'ajuste. La queue commune est donc une mesure.}
\appui{\S~4.8, p.~59 ; annexe~B, \S~B.1.}
\end{question}

\begin{question}{\og Vos montants ne sont pas actualisés.\fg{}}
\rep{Exact : la chaîne lit les montants nominaux, et l'inflation implicite de la base vaut
$2{,}19\,\%$ par an. Elle n'est pas appliquée à la calibration gelée, ce qui minore la queue, et
son effet est mesuré au backtest, où la moitié de la dérive de $4{,}00\,\%$ par an est monétaire.}
\appui{\S~4.3, p.~46 ; chapitre~5, \S~5.7.}
\end{question}

\begin{question}{\og Une de vos sources n'est pas reproductible.\fg{}}
\rep{Hackmageddon, oui, déclarée comme citation et non comme sortie du dispositif. C'est aussi
la source dont le modèle dépend le moins : elle fixe les parts par vecteur d'attaque, donc une
structure et jamais un niveau, et la catégorie d'un incident ne porte pas de segmentation. Son
statut de citation s'étend au multiplicateur de fréquence, ce qui est déclaré.}
\appui{\S~4.6, encadré p.~50.}
\end{question}

\begin{question}{\og Votre queue n'est-elle pas tirée par un sinistre hors norme ?\fg{}}
\rep{Non. Sur le périmètre calibré, le maximum ne vaut que $3{,}7$ fois le centile $99$, contre
$52{,}7$ sur la base entière : le plus gros sinistre du fichier ne relève pas du risque
informatique, et la plus grosse perte calibrée appartient à un continuum.}
\appui{\S~4.8, encadré p.~56 ; chapitre~5, \S~5.6.}
\end{question}

\begin{question}{\og Pourquoi ne pas ajouter d'autres bases ?\fg{}}
\rep{Chacune apporterait ses propres biais sans lever la limite structurelle, l'absence d'un
registre harmonisé. Et c'est l'indice de queue qui domine l'incertitude : mieux exploiter les
sources existantes, par bootstrap et validation, réduit davantage l'incertitude qu'une quatrième
base.}
\appui{\S~4.10, p.~63.}
\end{question}

\lepiege{Ne jamais présenter $\hat\xi_{\text{PRC}}=1{,}03$ comme un risque plus lourd : la
sévérité PRC est une puissance du nombre d'enregistrements, et une puissance multiplie l'indice
de queue par son exposant. Les deux indices ne portent pas sur le même objet.}

% ---------------------------------------------------------------------
\section*{Chapitre 5. Le socle mécaniste : fréquence et sévérité}
\addcontentsline{toc}{section}{Chapitre 5. Le socle mécaniste : fréquence et sévérité}
\reperes{p.~65 à 74. 5.1 les deux théorèmes (p.~65), 5.2 estimer $\xi$ (p.~66), 5.3 la fréquence
(p.~67), 5.4 la sévérité (p.~69), 5.5 Hill contre maximum de vraisemblance (p.~71), 5.6
validation (p.~71), 5.7 backtest (p.~73).}

\enbref{Deux lois reprises de la littérature et calibrées : une binomiale négative pour le nombre
d'incidents, une loi de Pareto généralisée au-delà de $20{,}03$~M\euro{} pour la sévérité. La
validation établit les formes, pas le niveau, et le backtest dit en années ce qu'il ne peut pas
tester.}

\begin{parcoeur}[À savoir par c\oe ur : la sévérité]
\begin{fichetab}
Deux théorèmes & Fisher-Tippett-Gnedenko : le maximum renormalisé tend vers une GEV ;
  Pickands-Balkema-de Haan : les excès au-delà d'un seuil tendent vers une GPD de même $\xi$ \\
Moments & $\xi>0$ queue lourde ; variance infinie dès $\xi\ge1/2$, espérance infinie dès
  $\xi\ge1$ \\
Calibration & $u=20{,}03$ M\euro{}, $91$ excès, $\hat\xi=0{,}5954$, $\hat\sigma=57{,}97$,
  $\zeta_u=0{,}151$ (gelé) \\
Quantile & $q_\alpha=u+\frac{\sigma}{\xi}\bigl[\bigl(\frac{1-\alpha}{\zeta_u}\bigr)^{-\xi}-1\bigr]$,
  soit $663$ M\euro{} à $99{,}5\,\%$ ; moyenne de queue $1\,752{,}4$ M\euro{} \\
Intervalles & $\xi$ : bootstrap $[0{,}304\,;0{,}831]$, delta-méthode $[0{,}320\,;0{,}870]$ ;
  quantile $[411\,;1\,037]$, facteur $2{,}5$ \\
Hill & $1{,}32$ à $k=86$, $+120{,}8\,\%$ ; les six valeurs observées tombent dans l'intervalle à
  $90\,\%$ de $2\,000$ échantillons simulés sous $\xi=0{,}595$ \\
Seuil & $\hat\xi$ stable à $0{,}5\,\%$ près jusqu'au percentile $85$ ; de $0{,}60$ à $0{,}38$
  au-dessus (dans l'IC) ; $0{,}98$ en dessous (hors IC) ; règle de stabilité : $19{,}87$ contre
  $20{,}03$ ; règle permissive : $12{,}88$ et un quantile $+16{,}0\,\%$ \\
Tests & famille GPD : Anderson-Darling $p=0{,}96$, Kolmogorov-Smirnov $p=0{,}89$ ; paramètres
  publiés imposés : $0{,}99$ et $0{,}89$ ; couverture réelle de l'IC à $90\,\%$ : $86\,\%$ \\
Jackknife & un retrait à la fois : étendue $111{,}7$ M\euro{} ($17{,}9\,\%$ de l'IC) ; trois plus
  grosses pertes retirées : sous la borne basse ; cinq : $657{,}2\to327{,}3$, $\xi=0{,}22$ \\
\end{fichetab}
\end{parcoeur}

\begin{parcoeur}[À savoir par c\oe ur : la fréquence et le backtest]
\begin{fichetab}
Clé de répartition & $\lambda_{ref}=341=\sum_v\lambda_v$, $\lambda_v=\lambda_{ref}\,\pi_v$ ;
  elle n'alimente pas le moteur \\
Binomiale négative & $r=\lambda/(\varphi-1)$, $p=r/(r+\lambda)$ ; $\varphi$ gardé constant d'un
  scénario à l'autre, $r$ recalculé \\
Surdispersion & $1{,}19$ sur les $14\,944$ cellules firme-année ; $2{,}24$ $[1{,}24\,;5{,}18]$ sur
  la série détendancée (dix-huit points) ; $9{,}20$ posé au moteur ; $a=0{,}122$
  $[0{,}054\,;0{,}223]$ contre $0{,}60$ posé \\
Poisson & rejeté à $p\approx10^{-30}$ ; passer à la binomiale négative relève le SCR de
  $+9\,\%$ seulement ; quantile du nombre annuel $74$ (Gamma) contre $82$ (lognormal) \\
Backtest & origine glissante 2004--2025, $539$ incidents, douze années notées ; couverture
  $91{,}7\,\%$ (binomiale négative) contre $75{,}0\,\%$ (Poisson) ; forme de queue non rejetée,
  $p=0{,}4087$ \\
Dérive & échelle de queue $+4{,}00\,\%$ par an, dont la moitié d'inflation monétaire ; au moins
  $+24{,}7\,\%$ sur le quantile ; facteur entre états $3{,}344\to3{,}124$ ; écart $-1{,}9\,\%$ \\
Puissance & doubler le taux de dépassement se détecte à $80\,\%$ en $1\,811$ années à
  $99{,}5\,\%$, $905$ à $99\,\%$, $78$ à $90\,\%$ \\
\end{fichetab}
\end{parcoeur}

\begin{question}{\og Pourquoi les dépassements de seuil plutôt que les maxima par blocs ?\fg{}}
\rep{Les maxima par blocs ne gardent qu'une observation par bloc : sur vingt-deux années, vingt-deux
points. Les dépassements de seuil exploitent les $91$ pertes au-delà de $20{,}03$~M\euro, et le
théorème de Pickands, Balkema et de Haan garantit la même famille limite, avec le même indice de
queue.}
\appui{\S~5.1, p.~65.}
\end{question}

\begin{question}{\og Écrivez le quantile de sévérité et donnez sa valeur.\fg{}}
\rep{$q_\alpha=u+\frac{\sigma}{\xi}\bigl[\bigl(\frac{1-\alpha}{\zeta_u}\bigr)^{-\xi}-1\bigr]$,
valable pour $\xi$ entre $0$ et $1$ ; à $99{,}5\,\%$ il vaut $663$~M\euro{}, et la moyenne de
queue $1\,752$. Puis la phrase qui évite l'erreur : c'est le montant d'\emph{un} sinistre, pas un
capital. Le capital est le quantile de la charge annuelle, et atteindre ce quantile annuel demande
de monter à environ $99{,}98\,\%$ par sinistre.}
\appui{\S~5.2, proposition et encadré p.~66.}
\end{question}

\begin{question}{\og $\xi=0{,}595$ : votre moyenne est-elle infinie ?\fg{}}
\rep{Non. Entre $1/2$ et $1$, la variance est infinie et l'espérance finie. La proposition qui
donne le quantile et la moyenne de queue pose explicitement $\xi\in(0,1)$. C'est sous la
calibration PRC, à $1{,}03$, que l'espérance cesse d'exister.}
\appui{\S~5.1, p.~65 ; \S~5.2, p.~66.}
\end{question}

\begin{question}{\og Votre estimateur de Hill vaut $1{,}32$ : votre queue n'est-elle pas bien plus lourde ?\fg{}}
\rep{Si elle valait $1{,}32$, l'espérance n'existerait pas. Mais $1{,}32$ est exactement ce que Hill
produit en échantillon fini sous $\xi=0{,}595$ : sur $2\,000$ échantillons simulés sous le modèle
publié, les six valeurs observées tombent dans l'intervalle à $90\,\%$. Et Hill croît avec $k$
quand $\hat\xi$ décroît avec le seuil : si la queue valait $1{,}32$, les deux dérives seraient
inversées.}
\appui{\S~5.5, p.~71.}
\end{question}

\begin{question}{\og L'indice de queue est-il stable quand on bouge le seuil ?\fg{}}
\rep{Non, et le mémoire ne le prétend pas. Il est stable au voisinage du seuil ; en remontant du
percentile $85$ au $95$, il descend de $0{,}60$ à $0{,}38$, mais reste dans l'intervalle déjà
publié, donc sans incertitude nouvelle. En descendant, il remonte à $0{,}98$ et sort de
l'intervalle : c'est le biais de seuil attendu, et la raison de ne pas descendre plus bas.}
\appui{\S~5.6, p.~71.}
\end{question}

\begin{question}{\og Pourquoi ce seuil-là ? Il vous arrange.\fg{}}
\rep{C'est la question à laquelle je réponds le plus volontiers, parce que la mesure va contre
l'insinuation. La règle de stabilité, appliquée à dix candidats, retombe sur $19{,}87$~M\euro{}
contre $20{,}03$ publiés, avec le même quantile, et c'est aussi le seuil le mieux ajusté du
balayage. La règle la plus permissive descendrait à $12{,}88$ et donnerait un quantile
\emph{supérieur} de $16\,\%$ : le seuil publié ne maximise pas le chiffre du mémoire.}
\appui{\S~5.6, p.~72.}
\end{question}

\begin{question}{\og $582$ pertes, $91$ excès : n'est-ce pas trop peu, et tout tient-il à trois sinistres ?\fg{}}
\rep{C'est peu, et c'est pourquoi la queue est l'incertitude dominante, chiffrée en intervalle sur
$\xi$ et propagée jusqu'au capital. Aucun retrait isolé ne sort le quantile de son intervalle ;
en revanche, retirer les trois plus grosses pertes le fait passer sous la borne basse, et cinq le
divisent presque par deux. C'est une condition de lecture, pas une fragilité : à $99{,}5\,\%$ le
quantile \emph{est} porté par les plus grands sinistres, et les retirer retire l'objet mesuré.}
\appui{\S~5.6, p.~72.}
\end{question}

\begin{question}{\og Votre surdispersion de $9{,}2$ est-elle mesurée ?\fg{}}
\rep{Non, elle est posée. L'indice de dispersion vaut $1+\lambda/r$, donc il n'est pas invariant
d'échelle : $1{,}19$ sur les cellules firme-année, $2{,}24$ sur la série annuelle détendancée,
$9{,}20$ au moteur par prudence. Une échelle de lecture, une agrégation, puis une marge assumée.}
\appui{\S~5.3, p.~68.}
\end{question}

\begin{question}{\og Poisson ou binomiale négative, qu'est-ce que cela change au capital ?\fg{}}
\rep{Plus $9\,\%$ seulement : la moyenne ne bouge pas, la queue du nombre bouge, mais à
$\xi\approx0{,}6$ le quantile annuel est porté par la perte unique dominante. La binomiale
négative reste la bonne loi : elle couvre $91{,}7\,\%$ des années au backtest contre $75{,}0\,\%$
pour Poisson, et le log-score, qui pénalise la largeur, désigne le même gagnant.}
\appui{\S~5.3, p.~68 ; \S~5.7, p.~73.}
\end{question}

\begin{question}{\og Votre intervalle à $90\,\%$ couvre-t-il vraiment $90\,\%$ ?\fg{}}
\rep{Non, $86\,\%$ : il est trop étroit, donc l'incertitude sur $\xi$ est sous-estimée. Il
faudrait élargir la demi-largeur d'environ $10\,\%$. L'écart est petit mais son sens compte : il
va du même côté que les deux autres défauts d'estimation, celui de la sous-estimation.}
\appui{\S~5.6, p.~71.}
\end{question}

\begin{question}{\og On ne peut pas valider un quantile à $99{,}5\,\%$.\fg{}}
\rep{Exact, et le mémoire le chiffre au lieu de l'éluder : détecter un taux de dépassement double
du nominal avec une puissance de $80\,\%$ demanderait $1\,811$~années. Le backtest valide le
centre et le corps de la loi, pas le quantile réglementaire, et le présenter autrement serait
faux.}
\appui{\S~5.7, encadré p.~74.}
\end{question}

\begin{question}{\og Votre backtest rejette le niveau. Pourquoi ne pas corriger ?\fg{}}
\rep{Il le rejette pour un motif nommé et mesuré : l'échelle de sévérité dérive de $4{,}00\,\%$
par an dans la queue, dont la moitié d'inflation monétaire, soit au moins $+24{,}7\,\%$ sur le
quantile. La calibration étant gelée, l'écart est déclaré et chiffré plutôt que corrigé. Et il ne
gonfle pas la thèse : sous la sévérité dérivée, le facteur entre états baisse de $3{,}344$ à
$3{,}124$. La formule juste est l'invariance d'échelle, pas l'invariance à la dérive.}
\appui{\S~5.7, p.~73 et encadré p.~74.}
\end{question}

% ---------------------------------------------------------------------
\section*{Chapitre 6. La cascade dirigée}
\addcontentsline{toc}{section}{Chapitre 6. La cascade dirigée}
\reperes{p.~75 à 84. 6.1 Vasicek à un facteur (p.~75), 6.2 généralisation~A (p.~76), 6.3
généralisation~B (p.~77), 6.4 Leontief (p.~77), 6.5 branchement (p.~78), 6.6 le seuil $K_j$
(p.~80), 6.7 calibration (p.~80), 6.8 formalisation (p.~81), 6.9 vérifications (p.~82).}

\enbref{La dépendance n'est pas ajoutée par une copule, elle est produite par la propagation d'un
incident le long d'une matrice dirigée ; la normalisation de Leontief garantit que la cascade
s'éteint, et trois propriétés qui portent l'argument sont vérifiées numériquement.}

\begin{parcoeur}[À savoir par c\oe ur : de Vasicek à la cascade]
\begin{fichetab}
Vasicek & $X_i=\sqrt\rho\,Y+\sqrt{1-\rho}\,\varepsilon_i$, défaut si $X_i\ge K_i$ ; deux
  défauts pour le cyber : dépendance symétrique (le jugement d'expert est dirigé à $81\,\%$) et
  probabilité de défaut non mesurable \\
Généralisation A & un facteur par pilier ; c'est le Vasicek multifacteur de Li (2000), pas une
  contribution ; $\rho_{ij}=\rho_j+u_{ij}$ (compromis hiérarchique) \\
$\rho_{ij}$ libre & jamais lisible sur un taux de conformité (la marge vaut $0{,}35$ quel que soit
  $\rho$) ; $813$ périodes pour séparer un écart de $0{,}10$ même à facteur connu ; $474$ au
  niveau du pilier ; $81$ sur une population entière ; $374$ entités sous le regroupement publié \\
Convention & $W_{jk}$ va de la source $j$ vers la cible $k$ : la ligne émet, la colonne reçoit ;
  P1 a la plus grande somme de ligne ($2{,}6$) et une somme de colonne trois fois plus petite que
  celle de P2 \\
Leontief & $W=g\,\mathrm{TRANS}/\max_j\sum_k\mathrm{TRANS}_{jk}$, donc
  $\rho(W)\le\max_j\sum_kW_{jk}\le g<1$ ; $\rho(\mathrm{TRANS})=1{,}461$ et
  $(I-\mathrm{TRANS})^{-1}$ a $21$ entrées négatives sur $25$ \\
Deux rayons & $\rho(W)=0{,}562\,g$, soit $0{,}506$ à $g=0{,}9$ ; $R_0=\rho(M)=0{,}054$, avec
  $M_{jk}=\Phi(\text{base}_j+W_{jk})-\Phi(\text{base}_j)$ et une incidence de fond de $4{,}5\,\%$ ;
  seuils critiques $g=1{,}78$ et $8{,}1$, hors domaine \\
Progéniture & $0{,}109$ (P1), $0{,}069$ (P4), $0{,}058$ (P2), $0{,}054$ (P3), $0{,}026$ (P5) :
  l'ordre de ROOT, par cohérence interne (même expert), pas par validation \\
\end{fichetab}
\end{parcoeur}

\begin{parcoeur}[À savoir par c\oe ur : les propriétés]
\begin{fichetab}
Descendance & $\mathbb E|S_a|\le[(I-W)^{-1}\mathbf 1]_a$ ; l'écart, de $10{,}6\,\%$ (P5) à
  $19{,}9\,\%$ (P1), compte les collisions \\
Loi exacte & $\Prob(S_a=A)=R(A)\,B(A)$, $B(A)=\prod_{j\in A,k\notin A}(1-W_{jk})$ \\
Moteur & une marche auto-évitante, même matrice moyenne que la cascade de la définition ; passer au
  branchement déplace le SCR de $-4\,\%$ \\
Monotonie & $W\le W'$ entraîne $S_a(W)$ dominé par $S_a(W')$ ; mais le capital n'est pas monotone
  en $A_{jk}$, d'où des sommets sans garantie \\
Ordre & la criticité est invariante par permutation si et seulement si $\mathrm{TRANS}$ est
  symétrique ; $22$ sous-ensembles sur $26$ changent de criticité \\
Co-occurrence & $2r\,S_{jk}+O(\lVert W\rVert^2)$ au premier ordre sous amorce uniforme ; la loi
  complète distingue $W$ de $W^\top$ : $1{,}931$ contre $1{,}746$ pilier par sinistre,
  $62{,}31\,\%$ contre $50{,}77\,\%$ de multi-piliers \\
Transposer & capital $3{,}6\,\%$, perte moyenne $7{,}9\,\%$, priorité de P1 à P4 ; jusqu'à
  $1\,173$ M\euro{} sur l'ensemble admissible \\
Non-transitivité & réfutée : zéro violation sur $420$ triplets, valeur propre minimale
  $+0{,}405$, zéro cycle \\
Seuil $K_j$ & $K_j=F^{-1}(1-p_j)$, $p_j=\Prob(S\ge u_j)$, $u_j$ seuil de matérialité (art.~18) ;
  valeurs extrêmes : biais $-0{,}8\,\%$ contre $-63{,}9\,\%$ pour la lognormale \\
Calibration & probit, un logit gonflerait les magnitudes d'un facteur $2{,}1$ ; sur données
  simulées, corrélation $0{,}92$ et pente $1{,}08$ ; sur données réelles, rien \\
\end{fichetab}
\end{parcoeur}

\begin{question}{\og Votre cascade peut-elle exploser ?\fg{}}
\rep{Non, et c'est garanti plutôt qu'espéré. La normalisation impose que chaque ligne de $W$
somme à au plus $g<1$, donc $\rho(W)\le g<1$ et la série $I+W+W^2+\cdots$ converge. Le piège est
documenté : la matrice d'expert brute a un rayon spectral de $1{,}461$, et
$(I-\mathrm{TRANS})^{-1}$ contient $21$ entrées négatives sur $25$, un choc positif produisant un
stress négatif ailleurs.}
\appui{\S~6.4, p.~77 et 78.}
\end{question}

\begin{question}{\og Démontrez la stabilité au tableau.\fg{}}
\rep{$W$ est à coefficients positifs et chaque ligne somme à au plus $g$, puisque $s_{\max}$ est la
plus grande somme de ligne de $\mathrm{TRANS}$. Le rayon spectral est majoré par toute norme
d'opérateur, en particulier la norme infinie, maximum des sommes de lignes :
$\rho(W)\le\lVert W\rVert_\infty\le g<1$. La série de Neumann converge, et $(I-W)^{-1}$ existe, à
coefficients positifs.}
\appui{annexe~A, p.~156.}
\end{question}

\begin{question}{\og Vous confondez deux rayons spectraux.\fg{}}
\rep{Non, ils sont distingués. $\rho(W)$ gouverne le système latent simultané et doit être
inférieur à un pour que $(I-W)^{-1}$ existe ; $R_0=\rho(M)$ gouverne la cascade d'incidents, qui
s'éteint si et seulement si $R_0<1$, comme en épidémiologie. $0{,}506$ et $0{,}054$ : les deux sont
stables, mais pas du même ordre, et la forme simultanée approche son seuil plus de quatre fois
plus vite ($1{,}78$ contre $8{,}1$).}
\appui{\S~6.5, encadré p.~79.}
\end{question}

\begin{question}{\og Pourquoi un Vasicek, qui vient du risque de crédit ?\fg{}}
\rep{Parce que la structure est la même : une variable latente franchit un seuil, et un facteur
commun crée la dépendance. Le mémoire dit ce qu'il emprunte : le Vasicek multifacteur à
corrélation générale est celui de Li (2000), employé tel quel, et déclaré comme n'étant pas une
contribution. La contribution commence à la contagion dirigée.}
\appui{\S~6.2, encadré p.~77.}
\end{question}

\begin{question}{\og Pourquoi pas une sensibilité par entité et par pilier ?\fg{}}
\rep{Parce qu'elle ne se lit jamais sur un taux de conformité, dont la marge vaut $0{,}35$ quelle
que soit la sensibilité, mais seulement sur la co-variation dans le temps. Même en connaissant le
facteur, il faudrait $813$ périodes pour séparer un écart de $0{,}10$ ; élargir le panel sature,
$81$ périodes restant nécessaires sur une population entière. D'où le compromis hiérarchique.}
\appui{\S~6.2, p.~76.}
\end{question}

\begin{question}{\og Votre moteur ne simule pas la cascade de votre définition.\fg{}}
\rep{C'est exact, et c'est écrit. Le moteur simule une marche qui ne repasse pas par un pilier
déjà touché ; la définition pose un processus de branchement. Les deux ont la même matrice moyenne
et le même nombre moyen de descendants directs, et passer de l'un à l'autre déplace le capital de
$-4\,\%$. Le choix ne porte pas le résultat.}
\appui{\S~6.8, p.~82.}
\end{question}

\begin{question}{\og La progéniture retrouve exactement ROOT : c'est une validation ?\fg{}}
\rep{Non, une cohérence interne : ROOT et TRANS viennent du même expert. Sa portée réelle est une
réduction de paramètres, ROOT se déduisant de TRANS au lieu d'être posé, soit cinq paramètres
d'expert en moins.}
\appui{\S~6.5, encadré p.~79.}
\end{question}

\begin{question}{\og Qu'est-ce que la dépendance à l'ordre, et pourquoi une copule ne peut-elle pas la reproduire ?\fg{}}
\rep{La vraisemblance d'une chaîne ordonnée,
$\mathrm{ROOT}_{i_1}\prod_k\mathrm{TRANS}_{i_ki_{k+1}}$, est invariante par permutation si et
seulement si $\mathrm{TRANS}$ est symétrique. Elle ne l'est pas : $22$ des $26$ sous-ensembles
changent de criticité selon l'ordre. Toute copule et tout score additif sont invariants par
permutation, donc aucun ne reproduit cet objet, et l'hypothèse requise est minimale.}
\appui{\S~6.8 et 6.9, p.~82.}
\end{question}

\begin{question}{\og Pourquoi avoir abandonné la non-transitivité, qui donnait son titre au mémoire ?\fg{}}
\rep{Parce qu'elle a été testée dans trois formulations et qu'aucune ne tient : zéro violation de
transitivité sur $420$ triplets, une matrice de corrélation parfaitement valide (valeur propre
minimale $+0{,}405$), zéro cycle de dominance. La réfuter valait mieux que la maintenir. Un travail
qui réfute sa propre hypothèse de départ est plus crédible, pas moins.}
\appui{\S~6.9, p.~83.}
\end{question}

\begin{question}{\og Que coûte de se tromper de sens ?\fg{}}
\rep{Transposer la matrice déplace le capital de $3{,}6\,\%$ et la perte moyenne de $7{,}9\,\%$,
et surtout fait passer la priorité de remédiation de P1 à P4. Sur l'ensemble admissible, l'écart
entre une matrice et sa transposée atteint $1\,173$~M\euro. Indiscernables pour les statistiques
disponibles, matériellement différentes.}
\appui{\S~6.9, p.~83.}
\end{question}

\begin{question}{\og Pourquoi un probit plutôt qu'un logit ?\fg{}}
\rep{Parce que le modèle générateur a un bruit gaussien et un seuil : l'estimateur est
littéralement un probit, et un logit gonflerait les magnitudes d'un facteur $2{,}1$. Sur données
simulées, le probit retrouve $W$ (corrélation $0{,}92$, pente $1{,}08$) : ce qui échoue sur les
données réelles, c'est la donnée, pas la méthode.}
\appui{\S~6.7, p.~80.}
\end{question}

\lepiege{L'argument \og la co-occurrence ne voit que $S$\fg{} ne vaut qu'au premier ordre et sous
amorce uniforme. La loi complète des piliers touchés distingue $W$ de $W^\top$ (fourche contre
collisionneur). Ne jamais présenter la frontière comme une impossibilité de principe.}

% ---------------------------------------------------------------------
\section*{Chapitre 7. La conformité multi-états}
\addcontentsline{toc}{section}{Chapitre 7. La conformité multi-états}
\reperes{p.~85 à 91. 7.1 deux échelles de temps (p.~85), 7.2 Vasicek multi-états (p.~86), 7.3
les quatre canaux (p.~86), 7.4 les valeurs de $g$ (p.~87), 7.5 l'état par pilier (p.~89), 7.6
Markov (p.~90), 7.7 le couplage (p.~91).}

\enbref{L'état de conformité n'est ni binaire ni figé : un Vasicek à trois états le fixe à
l'échelle rapide, une chaîne de Markov le fait évoluer à l'échelle lente, et il pilote le capital
par quatre canaux, sans aucune pénalité ajoutée.}

\begin{parcoeur}
\begin{fichetab}
Latente & $C^\star=\gamma\,\Theta+\sqrt{1-\gamma^2}\,\varepsilon$, $\gamma=0{,}68$, donc une
  corrélation $\gamma^2=0{,}462$ entre deux piliers \\
Seuils & $K_{\text{bas}}=\Phi^{-1}(p_{\text{NC}})$,
  $K_{\text{haut}}=\Phi^{-1}(p_{\text{NC}}+p_{\text{PC}})$ ; état actuel posé NC $35\,\%$, PC
  $35\,\%$, C $30\,\%$ \\
Quatre canaux & fréquence $21{,}56\to53{,}62$ incidents par an ; détection, $p_u$ multiplié par
  $0{,}85$ ou $1{,}20$ ($0{,}128\to0{,}181$) ; propagation $g$ de $0{,}45$ à $0{,}90$ ;
  accumulation P4 de $0$ à $0{,}68$ \\
Multiplicateur & $\lambda_{\text{NC}}=\lambda_{\text{C}}\sum_v\pi_v m_v=21{,}56\times2{,}4875$ ;
  centres $2{,}40$, $2{,}30$, $2{,}25$, $5{,}50$, $1{,}45$ ; aux bornes, de $1{,}933$ à $3{,}042$,
  soit $41{,}7$ à $65{,}6$ incidents ; les identifiants font $6{,}3\,\%$ des incidents et
  $13{,}9\,\%$ du multiplicateur \\
Valeurs de $g$ & $0{,}45$, $0{,}68$, $0{,}90$, posées ; lemme de croissance sur vingt points à
  nombres aléatoires communs ; entité $116{,}4$, $143{,}2$, $169{,}0$ M\euro{} ($+45{,}2\,\%$) ;
  correspondance très étroite $+3{,}9\,\%$ ; forfait $450$ M\euro{} plat \\
Élasticité & $0{,}41$ sur $g\in[0{,}20\,;0{,}95]$ ; un écart de plus de $10\,\%$ dès que $g$
  passe de $0{,}90$ à $0{,}75$ \\
Dépendance entre états & surcroît $\delta$ admissible jusqu'à $0{,}38$ ; $\Prob(\text{tous NC})$
  de $0{,}068$ à $0{,}088$ ; capital espéré $+2$ M\euro{} ($0{,}02\,\%$) ; forme additive
  $R^2=0{,}9945$, contributions $3\,524$ (P1), $2\,863$ (P4), $2\,022$ (P2), $1\,577$ (P3), $897$
  (P5) \\
Markov & NC vers PC vers C, C absorbant ; séjours Erlang-$k$ (mode strictement positif) ;
  $\overline{\mathrm{SCR}}(t)=\sum_e\Prob(e,t)\,\mathrm{SCR}_e$, une moyenne de quantiles (écart
  de $9{,}3\,\%$ au quantile du mélange) \\
Couplage & un seul $\Theta$ : $\mu(\Theta)=\mu_0e^{\beta_{\text{rem}}\Theta}$ et
  $\mathrm{SCR}_e(\Theta)=\mathrm{SCR}_e\,e^{-\beta_{\text{casc}}\Theta}$ \\
\end{fichetab}
\end{parcoeur}

\begin{question}{\og Qu'est-ce qui vous dit que la conformité réduit la propagation ?\fg{}}
\rep{Rien dans les données, et c'est écrit : c'est la seule hypothèse qualitative qui reste, et
le règlement lui-même prend parti puisqu'il impose ces exigences. Mais la thèse n'en dépend pas :
le capital est croissant en $g$, donc toute correspondance qui respecte l'ordre reproduit
l'écart. L'élasticité d'arc vaut $0{,}41$, et un écart de plus de $10\,\%$ n'exige qu'un
écartement de $0{,}15$ sur $g$ : l'exigence porte sur un ordre de grandeur, pas sur une mesure.}
\appui{\S~7.4, encadré p.~88.}
\end{question}

\begin{question}{\og Vos trois valeurs de $g$ sont arbitraires.\fg{}}
\rep{Posées, oui, faute d'entité publiant son état de conformité. Ce qui s'établit sans elles
tient en deux temps : le capital est croissant en $g$, vérifié sur vingt points sans aucun pas
décroissant ; donc l'existence et le sens de l'écart sont invariants, et seule l'amplitude est un
scénario. Même la correspondance très étroite, $0{,}85$ à $0{,}90$, garde $+3{,}9\,\%$ quand le
forfait reste plat.}
\appui{\S~7.4, p.~87 et 88.}
\end{question}

\begin{question}{\og D'où viennent vos multiplicateurs de fréquence ?\fg{}}
\rep{Ils se composent vecteur par vecteur : la part de chaque vecteur dans Hackmageddon, fois le
centre d'une fourchette tirée d'un effet publié, formation contre l'hameçonnage, authentification
forte, et ainsi de suite. Le total vaut $2{,}4875$, soit $53{,}62$ incidents par an. Le
multiplicateur est semi-ancré : sa provenance est tracée, sa valeur posée, et aux bornes des
fourchettes la fréquence non conforme va de $41{,}7$ à $65{,}6$.}
\appui{\S~7.3, p.~87.}
\end{question}

\begin{question}{\og Être conforme aux piliers 1 et 2 aide pour le 3 : vous l'ignorez ?\fg{}}
\rep{La dépendance existe déjà : un facteur commun de charge $0{,}68$, soit une corrélation de
$0{,}462$ entre deux piliers. Ce qui manque est sa structure, échangeable. Ajouter un surcroît sur
les seules paires concernées bouge nettement la loi des configurations, mais le capital espéré de
$2$~M\euro{} seulement, et c'est structurel : le capital des $32$ configurations est presque
additif ($R^2=0{,}9945$), et l'espérance d'une fonction additive ne dépend que des marges.
Réserve : l'argument vaut pour l'espérance, pas pour un quantile.}
\appui{\S~7.5, p.~89.}
\end{question}

\begin{question}{\og D'où sort la corrélation de $0{,}462$ ?\fg{}}
\rep{Elle n'est pas posée à part : c'est le carré de la charge du facteur commun,
$\gamma^2=0{,}68^2$.}
\appui{\S~7.2 et 7.5.}
\end{question}

\begin{question}{\og Pourquoi des séjours de type Erlang plutôt qu'exponentiels ?\fg{}}
\rep{Une exponentielle a sa densité maximale en zéro : elle autoriserait une remédiation quasi
instantanée. Un projet DORA a une durée minimale, que capte le mode strictement positif d'une
Erlang. Et ce n'est pas la même Erlang que la loi mélangeante de la fréquence, dont la forme
estimée, $0{,}63$, n'est pas entière.}
\appui{\S~7.6, p.~90.}
\end{question}

\begin{question}{\og Votre trajectoire de capital est-elle un quantile ?\fg{}}
\rep{Non, une moyenne de quantiles pondérée par la loi de l'état à chaque horizon. Le quantile du
mélange serait autre chose, la queue d'un mélange se prélevant dans ses composantes les plus
lourdes : l'écart vaut $9{,}3\,\%$ sur les $32$ configurations. La convention est déclarée et
tenue dans tout le mémoire.}
\appui{\S~7.6, p.~91.}
\end{question}

\begin{question}{\og Pourquoi deux échelles de temps ?\fg{}}
\rep{Parce que la propagation à l'intérieur d'un sinistre et l'évolution de la conformité d'une
année sur l'autre ne vivent pas sur la même horloge : les confondre ferait expliquer un
enchaînement de quelques heures par une tendance longue. C'est l'architecture éprouvée de la
migration de rating, et le préprint climatique sépare ses indices pour la même raison.}
\appui{\S~7.1, p.~85 ; \S~7.7, p.~91.}
\end{question}

% ---------------------------------------------------------------------
\section*{Chapitre 8. L'identifiabilité de la contagion dirigée}
\addcontentsline{toc}{section}{Chapitre 8. L'identifiabilité de la contagion dirigée}
\reperes{p.~92 à 102. 8.2 trois échecs (p.~92), 8.3 réversibilité (p.~94), 8.4 MOVEit (p.~96),
8.5 corpus (p.~97), 8.6 corpus étendu (p.~98), 8.7 biais de narration (p.~99), 8.8 séquences
(p.~101), 8.9 bayésien (p.~101), 8.10 frontière (p.~102).}

\enbref{La direction de la contagion n'est pas identifiable sur les données publiques agrégées,
quand sa co-occurrence l'est : trois sources échouent pour trois raisons distinctes, la
réversibilité en donne la raison de principe, et un corpus de rapports d'enquête corrobore la
structure sans calibrer les coefficients.}

\begin{parcoeur}[À savoir par c\oe ur : les trois échecs]
\begin{fichetab}
Il faudrait & des co-occurrences dans une même entité, une antériorité fiable, une taxonomie par
  pilier \\
PRC & $74\,783$ incidents, $25\,009$ organisations ; catégorie indéterminée $52{,}7\,\%$ ;
  survenance exploitable dans $49{,}6\,\%$ ; $71$ transitions ($3{,}6$ par cellule), $34$ en
  finance : pas de taxonomie, pas de matière \\
OpRisk & $4\,329$ transitions annuelles (environ $103$ par cellule) ;
  $T(M)=\tfrac12\lVert M-M^\top\rVert_1=119$ contre $128\pm27$ au placebo, $z=-0{,}33$ ; aucune
  des $18$ paires significative : le pas annuel efface l'ordre \\
VERIS & $10\,591$ incidents ($980$ en finance) ; \texttt{control\_failure} renseigné $13$ fois
  ($0{,}12\,\%$), chronologie datée sur deux étapes $126$ fois ($1{,}19\,\%$), sur trois $10$
  fois ; les deux ensemble : $1$ incident, aucun en finance \\
Temps inversé & $M_{\text{inv}}=M^\top$ par construction : corrélation $-1$ quel que soit le
  contenu, le test est dégénéré \\
Réversibilité & courant $J_{jk}=\pi_jP_{jk}-\pi_kP_{kj}$ ; la forme de Dirichlet ne voit que la
  partie réversible ; production d'entropie $\sigma\ge0$, nulle si et seulement si réversible \\
MOVEit & $188$ organisations financières, fenêtres de $180$ jours ; différence de différences
  $+1{,}17$ $[+1{,}07\,;+1{,}27]$ ; placebo un an plus tôt $+1{,}11$ ; effet net $+0{,}06$ \\
\end{fichetab}
\end{parcoeur}

\begin{parcoeur}[À savoir par c\oe ur : le corpus et la frontière]
\begin{fichetab}
Sept rapports & CSRB, FCA et PRA, GAO, SEC, OCC ; $17$ transitions ; $z=+3{,}93$
  ($p=0{,}0004$) ; P1 : $9$ sorties, $0$ entrée ; P2 : $0$ sortie, $7$ entrées \\
Dix rapports & plus CrowdStrike 2024, Log4Shell, Raphaels Bank 2015 ; $22$ transitions,
  $z=+5{,}12$ ; sources primaires seules $+4{,}08$ ; perspective entité seule $+3{,}62$ \\
Matrice ordonnée & $p_{jk}\sim\mathrm{Beta}(1+m_{jk},1+N_j-m_{jk})$ ; $12$ liens sur $20$
  crédiblement sous $0{,}5$, aucun au-dessus ; P3 vers P2 $0{,}77$ $[0{,}48\,;0{,}95]$, P1 vers P4
  $0{,}68$ $[0{,}40\,;0{,}89]$ ; les quatre sorties de P2 et trois des quatre entrées de P1 \\
Avec l'expert & corrélation non significative, $+0{,}08$ (matrice) et $+0{,}31$
  (antisymétrique, $p=0{,}19$) ; dénominateur corrigé : niveau $0{,}27\to0{,}07$, part de
  direction $60\,\%$ inchangée \\
Biais de narration & $11$ des $22$ arêtes partent de P1 ; $p$ exacte $1{,}5\cdot10^{-5}$ ;
  quatre arêtes retournées au pire : $0{,}084$ ; les onze retirées : $0{,}0039$ ; les sept
  rapports retirés : $0{,}50$ \\
En capital & P1 muette : entité $169{,}0\to127{,}1$ M\euro{} ($-24{,}8\,\%$), soit $0{,}81$ fois
  la bande de $51{,}7$ ; priorité de P1 à P4 dès $a=0{,}50$ ; écart entre états
  $+45{,}2\,\%\to+29{,}1\,\%$ \\
Bayésien & $A=S(2\theta-1)$, $\theta\sim\mathrm{Beta}(\alpha_0+n_{lo},\beta_0+n_{hi})$ ;
  crédible $81\,\%$ plus resserré que les bornes ; deux paires sur six tranchées \\
Trois niveaux & structure revendiquée ; direction paire par paire non revendiquée ; amplitude non
  identifiable (une bande de $656$ M\euro{} même les dix sens connus) \\
\end{fichetab}
\end{parcoeur}

\begin{question}{\og Pourquoi la direction n'est-elle pas identifiable ?\fg{}}
\rep{Parce que trois sources échouent, chacune pour sa raison : la chronologie PRC n'a ni la
taxonomie ni la matière, OpRisk a le volume ($103$ observations par cellule) mais un pas annuel qui
efface l'ordre, et VERIS a les bons champs sans les remplir, une seule fois sur $10\,591$. Ce n'est
donc pas un manque de volume, et la théorie des processus réversibles explique pourquoi une
statistique invariante par renversement du temps ne peut pas voir un courant.}
\appui{\S~8.2, p.~92 à 94 ; \S~8.3.}
\end{question}

\begin{question}{\og C'est donc une impossibilité de principe ?\fg{}}
\rep{Non, et le mémoire le dit expressément. Au premier ordre et sous amorce uniforme, la
co-occurrence ne dépend que de $S$ ; au-delà, la loi complète des piliers touchés par incident
distingue $W$ de $W^\top$, par exemple une fourche d'un collisionneur. La frontière est celle des
statistiques que les sources fournissent. D'où la quatrième exigence du registre : consigner tous
les piliers qui ont cédé.}
\appui{\S~8.3, encadré p.~95 ; chapitre~6, proposition 6.7.}
\end{question}

\begin{question}{\og Le test du temps inversé ne suffisait-il pas ?\fg{}}
\rep{Il est dégénéré sur des comptages : renverser la chronologie donne la transposée par
construction, donc une corrélation de $-1$ quel que soit le contenu. Il ne prouve rien. Seul le
placebo par permutation des années à l'intérieur de chaque firme est valide.}
\appui{\S~8.2, encadré p.~93.}
\end{question}

\begin{question}{\og Votre expérience naturelle MOVEit, qu'a-t-elle donné ?\fg{}}
\rep{Rien, et c'est un résultat. La différence de différences brute était grande et
significative, $+1{,}17$ jour-brèche avec un intervalle excluant zéro. Mais le même dispositif calé
sur une fausse date un an plus tôt donnait $+1{,}11$ : l'effet venait de la sélection du groupe de
contrôle. Une fois le placebo retranché, il reste $+0{,}06$. Et la cible serait restée non
attribuable, la taxonomie étant un vecteur d'attaque.}
\appui{\S~8.4, p.~96.}
\end{question}

\begin{question}{\og Vous parlez de martingale : êtes-vous en mesure risque-neutre ?\fg{}}
\rep{Non. Le mot est employé au sens des chaînes de Markov : la forme de Dirichlet est la
variation quadratique de la martingale associée à une fonction test, sous la mesure physique.
Aucune mesure risque-neutre, aucun actif répliqué, aucun test de martingalité au sens des
générateurs de scénarios : le mémoire ne valorise rien.}
\appui{\S~8.3, p.~95.}
\end{question}

\begin{question}{\og Que prouvent vos rapports d'enquête ?\fg{}}
\rep{La structure : P1 en amont, P2 en aval, $z=+3{,}93$ puis $+5{,}12$, et le signal survit aux
sources primaires seules et à la perspective entité seule. Pas les coefficients : la corrélation
avec la matrice d'expert n'est pas significative. Le corpus établit très bien où la contagion ne
va pas, et pas du tout où elle va fort.}
\appui{\S~8.5 et 8.6, p.~97 à 99.}
\end{question}

\begin{question}{\og Votre corpus sur-attribue à la gouvernance, et vous concluez qu'elle domine.\fg{}}
\rep{C'est l'objection la plus juste, et elle est dans le mémoire avant d'être dans votre bouche :
une enquête officielle est mandatée pour trouver une cause racine. Mais le biais est mesuré. Il
faudrait avoir retourné quatre des onze arêtes de P1, dans la pire configuration, pour perdre la
significativité, et même en les retirant toutes la direction reste significative, portée par P3
et P4 vers P2. En capital, rendre P1 muette ne coûte que $0{,}81$ fois la bande déjà publiée. Ce
qui bouge, c'est la priorité, qui passe à P4 dès une dégradation de moitié.}
\appui{\S~8.7, p.~99 et 100.}
\end{question}

\begin{question}{\og Un seul codeur a lu vos rapports.\fg{}}
\rep{Oui, c'est une limite déclarée. Le second codage en aveugle est préparé, et le script
d'accord ne fabrique rien sans le second jeu. En attendant, le biais de narration est borné par
une borne adverse : accorder au critique toute son objection coûte moins que l'ignorance
directionnelle déjà publiée. Un second codage ajouterait une mesure d'intersubjectivité, pas une
correction de niveau.}
\appui{\S~8.7, p.~100 ; annexe~F, \S~F.3.}
\end{question}

\begin{question}{\og Pourquoi pas un bayésien, qui donnerait une loi plutôt qu'une bande ?\fg{}}
\rep{Le mémoire le fait : le crédible à $90\,\%$ est $81\,\%$ plus resserré que les bornes. Mais
seules deux paires documentées sur six y sont tranchées, et le crédible dépend d'un prior quand les
bornes n'en dépendent pas. Les deux sont publiés : les bornes sont opposables, le crédible est
informatif.}
\appui{\S~8.9, p.~101.}
\end{question}

\begin{question}{\og Et les triplets ordonnés, qui portent plus d'information que les paires ?\fg{}}
\rep{Explorés et refermés, pour trois raisons : la comparaison naturelle a un support vide, P2
n'émettant jamais et P1 n'étant jamais atteint ; le test de dépendance au prédécesseur n'a aucune
puissance ; l'acyclicité s'obtiendrait deux fois sur trois par hasard. Il faudrait $158$ à $589$
chemins passant par un même pilier, contre neuf aujourd'hui.}
\appui{\S~8.8, p.~101.}
\end{question}

\begin{question}{\og Vos praticiens vous contredisent.\fg{}}
\rep{Deux sur trois contestent le sens unique de deux arcs, gouvernance source pure et incidents
comme puits. La plupart de leurs exemples décrivent une rétroaction corrective, qui relève de
l'évolution des états de conformité, pas de la propagation d'une perte ; et l'ensemble admissible
ne fixant aucune direction, le capital publié ne bouge pas. Surtout, aucun ne distingue un ordre
d'une cause commune : ils corroborent la frontière, pas la matrice.}
\appui{\S~8.5, encadré p.~98 ; annexe~C, \S~C.7 p.~179.}
\end{question}

\begin{question}{\og Qu'est-ce que votre frontière en trois niveaux ?\fg{}}
\rep{La structure, P1 amont et P2 aval, est revendiquée avec son confondant déclaré. La direction
paire par paire ne l'est pas, deux paires sur six seulement étant tranchées : les bornes restent.
L'amplitude n'est pas identifiable par cette voie : même les dix sens connus, une bande de
$656$~M\euro{} subsisterait. On revendique un ordre partiel, pas une matrice.}
\appui{\S~8.10, p.~102.}
\end{question}

% ---------------------------------------------------------------------
\section*{Chapitre 9. Identification partielle et bornes de capital}
\addcontentsline{toc}{section}{Chapitre 9. Identification partielle et bornes de capital}
\reperes{p.~104 à 108. 9.1 le modèle en couches (p.~104), 9.2 ce que le capital devient
(p.~105), 9.3 poser ou borner (p.~106), 9.4 le niveau est illustratif (p.~107), 9.5 l'échelle de
repli (p.~107).}

\enbref{Puisque la direction n'est pas identifiée, on ne la pose pas : on énumère les $1\,024$
matrices extrêmes compatibles avec la co-occurrence et l'on publie le capital en bornes, le socle
sans contagion restant défendable sans aucun paramètre non identifié.}

\begin{parcoeur}
\begin{fichetab}
Socle & $W=0$ : $5\,275$ M\euro{}, le plancher défendable \\
Ensemble admissible & $\mathcal W(t)=\{S+A : A^\top=-A,\ |A_{jk}|\le t\,S_{jk},\
  \rho(S+A)<1\}$ ; la contrainte $|A|\le S$ vient de la seule positivité ; $t=1$ est l'ignorance
  totale \\
Bornes & $2^{10}=1\,024$ sommets ; $[6\,858\,;8\,697]$ M\euro{}, facteur $1{,}27$, soit $+30$ à
  $+65\,\%$ au-dessus du socle ; le point d'expert, $8\,110$, est un point de l'ensemble \\
Sommets & l'atteinte des extrêmes n'est pas garantie (monter $A_{jk}$ baisse $W_{kj}$) ; aucun
  des $3\,000$ points intérieurs ne sort des bornes, en VaR ni en perte moyenne \\
Perte moyenne & bornée dans $[849\,;1\,100]$ M\euro{}, moins bruitée que la VaR \\
Valeur de $t$ & passer de $t=1$ à $t=0{,}5$ fait tomber la largeur de $1\,839$ à $695$
  ($-62\,\%$) \\
Poser ou borner & poser $W$ : un point, $8\,110$, qu'un panel biaisé déplace de $777$ M\euro{}
  ($10\,\%$) sans signal ; borner : la bande ne bouge pas, $79\,\%$ du capital est déjà fixé \\
Dominance & P1 en tête dans $45\,\%$ des configurations, P4 dans $36\,\%$ ($81\,\%$ à eux deux),
  P3 dans moins de $10\,\%$, P5 jamais \\
Niveau et rapport & le niveau varie d'un facteur $41$ (socle PRC $239$, fréquence triplée
  $11\,709$) ; les rapports d'au plus $33\,\%$ ; le groupe de tête est le même \\
Repli et cran nul & un rang~1 capte $74\,\%$ de la direction et resserre de $25\,\%$ ; à $u=0$,
  P1 et P4 deviennent indiscernables ($1\,051$ contre $1\,052$ M\euro), séparés de $353$ avec la
  direction du classeur \\
\end{fichetab}
\end{parcoeur}

\begin{question}{\og Alors, quel est votre SCR ?\fg{}}
\rep{Le niveau n'est pas une mesure, et le mémoire le démontre plutôt que de s'en excuser : il
bouge d'un facteur $41$ en ne changeant que la source de sévérité et la fréquence. Ce que le modèle
détermine, c'est un intervalle et son prix : entre $6\,858$ et $8\,697$~M\euro{} sur les $1\,024$
sommets, et la largeur, $1\,839$, est le coût de l'ignorance de la direction, connu d'avance. Ce
qui résiste au changement de source, c'est le rapport au socle et le groupe de tête des piliers.}
\appui{\S~9.2, p.~105 ; \S~9.4, p.~107.}
\end{question}

\begin{question}{\og Pourquoi borner, au lieu d'estimer ?\fg{}}
\rep{Parce que la direction n'est pas \emph{mal estimée}, elle n'est \emph{pas identifiée}, et
les deux appellent des réponses opposées : un paramètre mal estimé appelle plus de données, un
paramètre non identifié appelle une autre approche. La borne a en outre une propriété qu'un
intervalle de confiance n'a pas : sa largeur se paie linéairement en $t$ et son maximum est connu
d'avance.}
\appui{\S~9.1 et 9.2 ; chapitre~8.}
\end{question}

\begin{question}{\og Qu'est-ce qui garantit que les sommets donnent les extrêmes ?\fg{}}
\rep{Rien a priori, et le mémoire le dit : monter $A_{jk}$ élève $W_{jk}$ mais abaisse $W_{kj}$,
donc la réponse n'est pas monotone. On procède par réfutation : $3\,000$ points intérieurs tirés
uniformément, aucun ne sort des bornes, ni en VaR ni en perte moyenne. C'est une hypothèse vérifiée
numériquement, pas un théorème, et c'est écrit ainsi.}
\appui{\S~9.2, encadré et figure 9.1, p.~105.}
\end{question}

\begin{question}{\og Que vaut le point d'expert, $8\,110$ ?\fg{}}
\rep{C'est un point de l'ensemble, jamais sa mesure. Si le panel qui le fixe est biaisé, il se
déplace de $777$~M\euro{} et la priorité change de pilier sans que rien ne le signale au lecteur.
La bande, elle, contient tous ces dérapages et ne bouge pas d'un euro.}
\appui{\S~9.3, p.~106.}
\end{question}

\begin{question}{\og À quoi servirait alors une élicitation ?\fg{}}
\rep{À départager P1 et P4, qui concentrent $81\,\%$ des configurations, et à resserrer la bande.
Elle ne peut pas faire remonter P2, P3 ou P5 : l'ensemble admissible les exclut déjà de la tête,
P5 n'y arrivant jamais premier.}
\appui{\S~9.3, p.~106.}
\end{question}

\begin{question}{\og Pourquoi regarder la perte moyenne et pas seulement la VaR ?\fg{}}
\rep{À queue lourde, la VaR obéit au principe de la perte unique dominante et reste peu sensible
à la structure de cascade : un tableau de sensibilité fondé sur elle serait plat et faussement
rassurant. La perte moyenne, bornée dans $[849\,;1\,100]$, est bien moins bruitée, et c'est sur la
décision, l'ordre de remédiation, que la contribution se lit.}
\appui{\S~9.2, encadré p.~106.}
\end{question}

\begin{question}{\og Si la direction pèse si peu sur le niveau, à quoi sert-elle ?\fg{}}
\rep{À trancher la décision. Au cran nul, où aucune source ne se distingue, P1 et P4 deviennent
indiscernables, $1\,051$ contre $1\,052$~M\euro{} de gain de remédiation ; avec la direction du
classeur, $353$~M\euro{} les séparent. Sans direction, on ne sait plus lequel remédier d'abord.}
\appui{\S~9.5, encadré p.~107.}
\end{question}

\begin{question}{\og Le niveau est-il au moins une indication ?\fg{}}
\rep{C'est une illustration, pas une mesure. Le niveau bouge d'un facteur $41$ quand on change
des \emph{entrées} dont le périmètre n'est pas transposable, alors que les rapports au socle
varient d'au plus $33\,\%$, trente fois moins, et que le groupe de tête reste le même. Le rapport
le moins stable, la largeur de bande rapportée au socle, varie de $32{,}9\,\%$ : la stabilité du
relatif n'est pas parfaite, elle est d'un autre ordre.}
\appui{\S~9.4, p.~107.}
\end{question}

% ---------------------------------------------------------------------
\section*{Chapitre 10. Résultats}
\addcontentsline{toc}{section}{Chapitre 10. Résultats}
\reperes{p.~109 à 128. 10.1 SCR par état et chiffre de tête (p.~109), 10.2 décomposition par
pilier (p.~113), 10.3 Shapley et Euler (p.~114), 10.4 trajectoire (p.~116), 10.5 priorisation
(p.~117), 10.6 robustesse (p.~118), 10.7 sensibilité à la propagation (p.~119), 10.8 copule
(p.~120), 10.9 regard réglementaire (p.~122), 10.10 ordre de grandeur (p.~122), 10.11 entités
réelles (p.~124), 10.12 agrégation (p.~126), 10.13 détenir ou transférer (p.~127).}

\enbref{Le combien : sur les quatre canaux, le besoin de capital du secteur passe de $6\,049$ à
$20\,188$~M\euro{} ; le chapitre l'attribue, le déroule dans le temps, le confronte à une copule, à
la Formule Standard et à quatre entités réelles, puis le traduit en décision.}

\begin{parcoeur}[À savoir par c\oe ur : le chiffre de tête]
\begin{fichetab}
Quatre lectures & score seul, $g$ figé au niveau NC : $7\,987\to16\,847$ ($\times2{,}11$) ;
  fréquence et propagation : $6\,664\to15\,074$ ($\times2{,}26$) ; plus détection :
  $5\,932\to16\,595$ ($\times2{,}80$) ; quatre canaux : $6\,049\to20\,188$ ($\times3{,}34$) \\
Pourquoi la quatrième & la seule où l'état conforme l'est sur tous les canaux du modèle ; les
  trois autres sont des remédiations partielles, publiées comme telles \\
Détection & résolue : Shapley $2\,928\pm343$, fermeture $4\,562\pm819$, signe stable sur seize
  graines \\
Interaction & canaux isolés $9\,138$, interaction $5\,001$ ($35\,\%$) ; produit des quatre facteurs
  isolés $3{,}47$ contre $3{,}34$ publié : une composition presque multiplicative \\
Fréquence & seule, $10\,377$ ($\times1{,}72$) ; les trois autres réunis, $11\,413$ ($\times1{,}89$) ;
  aucun autre canal seul n'atteint $8\,000$ \\
Réconciliation & conforme autour de $5\,900$ (lecture~C, $5\,932$) contre socle $5\,275$ à $W=0$ :
  deux points d'un même modèle \\
Sources & le même état conforme vaut $8\,301$ sous OpRisk et $2\,589$ sous PRC (lecture~A) \\
Surcoût (lecture B) & NC contre C, OpRisk : médiane $7\,401$, IC à $90\,\%$ $[2\,830\,;38\,793]$,
  $100\,\%$ des tirages positifs \\
Mesure & $\mathrm{VaR}_{99{,}5\,\%}(L)$ \emph{non centrée} : la perte attendue n'étant couverte
  par aucune provision, elle n'est pas déduite \\
\end{fichetab}
\end{parcoeur}

\begin{parcoeur}[À savoir par c\oe ur : attribution, décision, comparaison]
\begin{fichetab}
Entre piliers & interaction de signe non résolu en VaR ($-137$ à $+1\,395$ sur six graines),
  positive en perte moyenne ($+532$ à $+659$) ; à ne jamais ajouter à celle entre canaux \\
Shapley & $\varphi_j=\sum_S\frac{|S|!(5-|S|-1)!}{5!}[v(S\cup\{j\})-v(S)]$ ; P1 $34{,}0\,\%$, P4
  $25{,}7\,\%$, P2 $17{,}1\,\%$, P3 $13{,}6\,\%$, P5 $9{,}6\,\%$ ; somme $10\,927$, marginaux
  $9\,531$ ; amorces ROOT $30$, $27$, $18$, $15$, $9\,\%$ \\
Euler & $C_j=\mathbb E[L_j\mid L\ge\mathrm{VaR}]$ ; presque plat sous PRC (P2 et P4 $23\,\%$, P1
  $20$, P3 $18$, P5 $17$) ; indicatif sous OpRisk (variance infinie) \\
Queue & $\mathrm{TVaR}/\mathrm{VaR}\approx2{,}1$ contre l'asymptote $1/(1-\xi)=2{,}5$ \\
Trajectoire & lecture~B : $10\,752\to7\,041$ en cinq ans, vers un conforme d'environ $6\,925$ \\
Ordre optimal & sur les $120$ ordres, P1, P4, P2, P3, P5, l'ordre ROOT, distinct du glouton ;
  pire contre optimal $21\,\%$ \\
Ensemble admissible & P1 premier dans $43{,}6\,\%$, P4 dans $37{,}7\,\%$, écart de regret maximal
  $1{,}5\,\%$ ; P5 jamais premier \\
Élasticités & $\xi$ $3{,}79$ ; $\sigma$ $0{,}97$ ; $p_u$ $0{,}82$ ; $\lambda$ $0{,}45$ ; $\varphi$
  $-0{,}44$ ; $g$ $0{,}37$ ; seuil seul $0{,}03$ ; ablation : propagation $-20\,\%$, queue
  $-76\,\%$ \\
Copule & gaussienne à marges égales : $+0{,}1\,\%$ (et $+4{,}8\,\%$ en moyenne sur quatre
  graines) ; Student $+17\,\%$ ; surcoût $6\,760$ contre $10\,927$ ($-38\,\%$) ; interaction
  moyenne $0$ contre $+560$ ; part de l'effet : Formule Standard $0\,\%$, copule $62\,\%$ \\
\end{fichetab}
\end{parcoeur}

\begin{question}{\og Vous dites ne pas connaître la direction, et vous publiez un facteur $3{,}34$. Laquelle est votre réponse ?\fg{}}
\rep{Les deux, et elles ne mesurent pas la même chose. Le facteur $3{,}34$ mesure l'effet des
quatre canaux à direction fixée, celle du classeur ; la bande mesure ce que coûte l'ignorance de
la direction. Le chapitre réconcilie d'ailleurs son conforme d'environ $5\,900$ avec le socle de
$5\,275$ du chapitre~9 : deux points du même modèle, pas deux mesures du même objet.}
\appui{\S~10.1, encadré p.~109 ; chapitre~9.}
\end{question}

\begin{question}{\og Pourquoi publier la lecture à quatre canaux, celle qui donne le plus gros chiffre ?\fg{}}
\rep{Parce que c'est la seule où l'état conforme est conforme sur tous les canaux du modèle. Dans
les autres, l'entité dite conforme propage encore comme une entité défaillante, ou garde sa
détection ou son accumulation au niveau non conforme : c'est une remédiation partielle, et
l'appeler conforme sous-estimerait l'écart. Les trois autres lectures restent publiées, comme
remédiations partielles.}
\appui{\S~10.1, encadré p.~110.}
\end{question}

\begin{question}{\og Vos canaux isolés somment à $9\,138$, l'écart vaut $14\,139$ : où sont les $5\,001$ ?\fg{}}
\rep{Dans l'interaction, et elle vient d'une composition presque multiplicative : le produit des
quatre facteurs isolés vaut $3{,}47$ pour un facteur publié de $3{,}34$. Une décomposition
additive d'effets qui se multiplient produit mécaniquement une interaction positive. D'où la
conséquence de gestion : chiffrer une remédiation par l'effet de \emph{fermeture}, depuis l'état
non conforme, jamais par la somme des effets isolés.}
\appui{\S~10.1, encadré p.~110 ; annexe~B, \S~B.4.}
\end{question}

\begin{question}{\og Y a-t-il un double comptage entre vos canaux ?\fg{}}
\rep{Non : chaque canal agit sur un objet distinct du modèle, le nombre d'incidents, la détection,
la propagation, l'accumulation par un prestataire. Et la colonne des effets de fermeture, qui
semble compter deux fois, obéit à une identité : sa somme vaut $\sum_S|S|\,m(S)$, un croisé d'ordre
$k$ y étant compté $k$ fois, $1\times9\,138+2\times5\,345-3\times691+4\times346=19\,141$. La vraie
propriété est l'encadrement $9\,138\le14\,139\le19\,141$.}
\appui{annexe~B, encadré p.~166.}
\end{question}

\begin{question}{\og Votre modèle est-il super-additif ?\fg{}}
\rep{Il faut distinguer trois énoncés. Au sein d'un sinistre, les coûts des piliers touchés
s'additionnent : c'est une hypothèse de construction. En fonction des piliers non conformes, le
capital est presque additif, $R^2=0{,}9945$ : c'est une mesure. En fonction des quatre canaux, il
est super-additif, $+35\,\%$ : c'est aussi une mesure. Répondre \og oui\fg{} tout court serait faux.}
\appui{chapitre~11, \S~11.7 p.~140.}
\end{question}

\begin{question}{\og Shapley et Euler ne donnent pas le même classement.\fg{}}
\rep{C'est le résultat. Shapley regarde la \emph{source} de la non-conformité et concentre le
surcoût sur P1, $34\,\%$ ; Euler regarde le pilier \emph{touché} et répartit le capital presque
uniformément, P2 et P4 en légère tête. La remédiation se décide côté source ; le capital se loge
là où les cascades atterrissent.}
\appui{\S~10.3, encadré p.~116.}
\end{question}

\begin{question}{\og Votre SCR est-il centré ?\fg{}}
\rep{Non : partout le quantile $\mathrm{VaR}_{99{,}5\,\%}(L)$ non centré, au niveau de
Solvabilité~II. La perte attendue d'une non-conformité n'est couverte par aucune provision, donc
elle n'en est pas déduite. Si l'on demande l'ordre de grandeur : centré, le chiffre de tête
tomberait à environ $5\,430\to16\,320$, un facteur de l'ordre de $3{,}0$ (calcul hors mémoire).}
\appui{\S~10.3, p.~115.}
\end{question}

\begin{question}{\og Pourquoi une VaR et pas une TVaR ?\fg{}}
\rep{Parce que Solvabilité~II définit le capital comme une VaR à $99{,}5\,\%$ à un an, et qu'un
travail qui se compare au dispositif prudentiel doit parler sa mesure. La TVaR intervient comme
diagnostic : noyau d'allocation d'Euler, et forme de queue, $\mathrm{TVaR}/\mathrm{VaR}\approx2{,}1$
contre $2{,}5$ en théorie. L'argument décisif contre elle est d'existence : sous la calibration
PRC, $\hat\xi=1{,}03>1$, l'espérance est infinie et la mesure n'existe pas.}
\appui{\S~10.3, p.~115.}
\end{question}

\begin{question}{\og Quel canal une direction des risques doit-elle maîtriser en premier ?\fg{}}
\rep{La fréquence, et le test de résistance inversé le montre. Seule, elle porte le capital à
$10\,377$, quand aucun autre canal seul n'atteint $8\,000$. Au-delà de $11\,413$, ce que les trois
autres canaux atteignent réunis, aucune cible n'est atteignable sans elle. La lecture inverse
désigne le canal dont la maîtrise interdit les scénarios sévères.}
\appui{\S~10.5, p.~118 ; annexe~B, \S~B.4 p.~167.}
\end{question}

\begin{question}{\og Qu'est-ce qu'un dirigeant fait de votre travail lundi matin ?\fg{}}
\rep{Trois choses. Il sait quel pilier remédier en premier, et que l'ordre optimal n'est pas
l'ordre glouton, parce que les interactions imposent de raisonner sur tout l'horizon : $21\,\%$ de
capital intégré entre le pire ordre et l'optimal. Il sait comparer remédier, détenir et transférer
sur un seul axe en euros par an. Et il sait quelle donnée réclamer à son registre d'incidents.}
\appui{\S~10.5, p.~117 ; \S~10.13 ; chapitre~12.}
\end{question}

\begin{question}{\og Votre recommandation dépend-elle de paramètres que personne ne sait calibrer ?\fg{}}
\rep{Non : l'ordre de remédiation optimal ne dépend que des SCR par configuration, donc il est
invariant aux taux de transition de la chaîne de Markov, pourtant les paramètres les moins
calibrables du modèle. Ces taux fixent le calendrier, jamais la séquence.}
\appui{\S~10.6, encadré p.~118.}
\end{question}

\begin{question}{\og P1 est-il vraiment le premier pilier à remédier ?\fg{}}
\rep{À matrice fixée, oui : Shapley $34\,\%$, ordre optimal P1 puis P4. Sur l'ensemble admissible,
P1 et P4 sont presque à égalité, $43{,}6\,\%$ contre $37{,}7\,\%$ des configurations, avec un
écart de regret de $1{,}5\,\%$. Sous le biais de narration, la priorité bascule vers P4 à mi-chemin.
Ce qui est robuste, c'est que P1 et P4 devancent les trois autres, et que P5 n'est jamais premier.}
\appui{\S~10.6, encadré p.~119.}
\end{question}

\begin{question}{\og Votre modèle est-il sensible aux probabilités de propagation, qui ne sont pas observées ?\fg{}}
\rep{C'est le levier le moins sensible : élasticité $0{,}37$ contre $3{,}79$ pour l'indice de
queue, et retirer la propagation fait tomber la VaR de $20\,\%$ quand retirer la queue lourde la
fait tomber de $76\,\%$. La direction, elle, se borne au lieu de se stresser : $1\,839$ à ignorance
totale, soit $34{,}9\,\%$ du socle. La vigilance porte d'abord sur les valeurs extrêmes.}
\appui{\S~10.7, p.~120.}
\end{question}

\begin{question}{\og Pourquoi une cascade et pas une copule ?\fg{}}
\rep{La copule a été testée en benchmark adverse. À marges égales, la gaussienne retombe sur le
capital de la cascade, $+0{,}1\,\%$ : le niveau central n'identifie pas la dépendance. Mais elle ne
sait pas faire le reste : le surcoût sort à $6\,760$ contre $10\,927$, l'interaction moyenne est
nulle par construction contre $+560$, et la dépendance ne durcit pas quand la conformité se
dégrade. La dépendance n'est pas le mécanisme, c'en est l'ombre portée.}
\appui{\S~10.8, p.~120 à 122.}
\end{question}

\begin{parcoeur}[À savoir par c\oe ur : l'entité, l'agrégation, le transfert]
\begin{fichetab}
Descente & $\lambda(A)=\lambda_c(A/A_c)^{b_\lambda}$, $b_\lambda=0{,}074$ ;
  $m(A)=(A/A_m)^b$, $b=0{,}087$, $m=0{,}854$ à la taille de l'entité notionnelle \\
Échelles & secteur $\lambda=21{,}56$ : $8\,123$ M\euro{} (moyenne $983$) ; seau d'au moins dix
  événements $0{,}210$ : $488$ ; seau toutes firmes $0{,}100$ : $229$ ; $\lambda$ lu à la taille
  $0{,}092$ : $198$ ; les deux canaux composés : $169$ (moyenne $3{,}6$), bande
  $[131{,}5\,;183{,}3]$ \\
Non-linéarité & diviser $\lambda$ par $235$ divise la moyenne par $270$ et le capital par $48$ :
  le quantile est porté par le plus grand sinistre \\
Deux incertitudes & à l'entité, la bande de direction mesure $51{,}7$ M\euro{} ; l'IC de $\xi$
  donne $140{,}9$ à $191{,}6$, soit $51{,}5$ : du même ordre \\
Quatre SFCR 2024 & non-vie A : actif $2\,319$, SCR $425$, besoin $112$ ($26{,}4\,\%$) ; non-vie
  B : $5\,305$, $812$, $136$ ($16{,}7\,\%$) ; vie C : $37\,845$, $1\,587$, $196$ ($12{,}3\,\%$) ;
  vie D : $309\,800$, $14\,800$, $294$ ($2{,}0\,\%$) \\
Hors domaine & un facteur $134$ sur l'actif ne divise la sévérité que par $1{,}53$ et la
  fréquence par $1{,}44$ ; au critère de $10\,\%$ du SCR, la bascule tombe entre $37\,845$ et
  $309\,800$ ; l'entité notionnelle ($19\,231$) est dedans : $169$ est un majorant ; le test du
  levier (médian $10{,}0\,\%$) donne $9\,\%$ et ne l'accable pas \\
Forfait de $3\,\%$ & $54{,}1$ contre $59$ publié ($-8\,\%$), $62{,}1$ contre $39{,}7$ ($+56\,\%$) :
  un ordre de grandeur, ni la charge ni une borne \\
Agrégation & $\tau(\rho,d)=\frac{\sqrt{1+2\rho d+d^2}-1}{d}=\rho+\frac d2(1-\rho^2)+O(d^2)$ ;
  à $\rho=0{,}25$, rapport additif sur corrélé de $3{,}4$ à $3{,}9$ ; l'écart pèse de $0{,}58$ à
  $8{,}21\,\%$ du SCR publié \\
Transfert & coût du capital $6\,\%$, $4{,}75\,\%$ après la directive (UE) 2025/2 ; sinistres sur
  primes $48{,}5\,\%$ ; multiple $\mathrm{SCR}/\mathbb E[L]$ de $47$ à l'entité contre $8{,}3$ au
  secteur ; détention $8{,}0$ M\euro/an ($10{,}1$ à $6\,\%$), $2{,}2$ à $2{,}8$ fois la perte
  attendue ; portée optimale $169$, coût $8{,}0\to2{,}9$ ($-64\,\%$) ; seuil d'inefficacité
  $78\,\%$ ; une entité conforme : $-81\,\%$ \\
\end{fichetab}
\end{parcoeur}

\begin{question}{\og Ce montant sectoriel, que vaut-il pour mon entité ?\fg{}}
\rep{À la taille d'une entité, fréquence et sévérité lues à sa taille, le même modèle donne
$169$~M\euro, dans une bande de $131{,}5$ à $183{,}3$. Mais le mémoire déclare lui-même l'entité
notionnelle hors du domaine de validité : la sévérité décroît beaucoup moins vite que la taille, un
facteur $134$ sur l'actif ne la divisant que par $1{,}53$. Ce $169$ est donc un majorant d'ordre
de grandeur. Ce qui transfère à une entité, c'est le rapport au socle et le groupe de tête des
piliers.}
\appui{\S~10.10, p.~122 à 124 ; \S~10.11, encadré p.~126.}
\end{question}

\begin{question}{\og Pourquoi votre entité notionnelle est-elle hors domaine ?\fg{}}
\rep{Sur quatre assureurs réels, la charge passe de $2{,}0\,\%$ du SCR publié pour le plus grand à
$26{,}4\,\%$ pour le plus petit. En posant qu'un seul risque de non-conformité ne peut peser plus de
$10\,\%$ du SCR total, la bascule tombe entre $37\,845$ et $309\,800$~M\euro{} d'actif, et l'entité
notionnelle, $19\,231$, est dedans. L'élasticité corrige l'échelle de la sévérité, jamais sa forme,
qui reste celle de grandes institutions. Le test du levier, lui, ne l'accable pas.}
\appui{\S~10.11, p.~125 et 126.}
\end{question}

\begin{question}{\og Pourquoi anonymiser des entités dont les rapports sont publics ?\fg{}}
\rep{Par prudence et non par confidentialité : le modèle n'utilise que le bilan, et l'état de
conformité y est \emph{supposé}. Nommer une entité reviendrait à lui attribuer une non-conformité
que rien n'observe. La table de correspondance existe et vit dans un script : le travail reste
vérifiable par qui a le dépôt.}
\appui{\S~10.11, encadré p.~124.}
\end{question}

\begin{question}{\og Pourquoi refusez-vous d'appeler cela un SCR ?\fg{}}
\rep{Parce qu'aucun module DORA n'existe dans la Formule Standard. La grandeur calculée est un
besoin de capital ORSA, donc de pilier~2, et le rapport au SCR publié est une mise à l'échelle,
pas une part : diviser un besoin de pilier~2 par le SCR réglementaire n'en fait pas une composante.
Écrire \og le SCR DORA de telle entité\fg{} ferait croire à un objet qui n'existe pas.}
\appui{\S~10.11, encadré p.~124 ; chapitre~2.}
\end{question}

\begin{question}{\og La Formule Standard n'est-elle pas au moins une borne ?\fg{}}
\rep{Non. Sa réponse à la conformité est nulle par construction : $0\,\%$ de l'effet, contre
$62\,\%$ pour une copule et $100\,\%$ pour la cascade. Et le repère de $3\,\%$ des provisions
s'écarte de $-8\,\%$ et $+56\,\%$ des modules publiés : il ne se cite que comme ordre de grandeur. La
coïncidence de $488$ contre $450$ au seuil de dix événements ne survit pas à la correction
d'échelle.}
\appui{\S~10.9, p.~122 ; \S~10.11, p.~125.}
\end{question}

\begin{question}{\og Que retient le capital total de votre écart ?\fg{}}
\rep{Tout, sous la convention de Solvabilité~II, puisque la charge opérationnelle s'ajoute au BSCR
sans passer par la matrice de corrélation. Agrégé comme un module corrélé, il ne serait retenu qu'à
hauteur de $\tau$, qui tend vers $\rho$ pour un petit écart : à $\rho=0{,}25$, l'écart de
traitement vaut $3{,}4$ à $3{,}9$. Un euro économisé sur DORA vaut donc trois à quatre fois plus,
en capital, qu'un euro sur un risque qui se diversifie.}
\appui{\S~10.12, p.~126 et encadré p.~127.}
\end{question}

\begin{question}{\og Détenir ou transférer ?\fg{}}
\rep{À l'échelle d'entité, détenir coûte $8{,}0$~M\euro{} par an au taux révisé, $2{,}2$ à $2{,}8$
fois la perte attendue : immobiliser coûte plus que le risque. Un traité en excès de perte de portée
égale au SCR, $169$, fait tomber le coût total à $2{,}9$, soit $-64\,\%$. C'est une borne
supérieure, entamée par la capacité du marché, les exclusions et le risque de contrepartie ; mais il
faudrait que $78\,\%$ de la couverture soit inefficace pour renverser la conclusion. Et une entité
conforme voit ce coût baisser de $81\,\%$.}
\appui{\S~10.13, p.~127 et 128.}
\end{question}

\begin{question}{\og Votre bande d'identification est-elle informative, ou trop large pour rien apprendre ?\fg{}}
\rep{Elle est informative, et la comparaison le montre : à l'échelle d'entité, elle mesure
$51{,}7$~M\euro, quand faire varier le seul indice de queue sur son intervalle à $90\,\%$ déplace
le capital de $51{,}5$. Ignorer la direction coûte autant qu'ignorer la queue, ce qui fait de la
spécification de reporting un levier du même ordre qu'une meilleure estimation de la queue.}
\appui{\S~10.10, encadré p.~123.}
\end{question}

\lepiege{Le multiple de capital de $47$ à l'entité contre $8{,}3$ au secteur n'est pas une
incohérence : c'est la mutualisation lue à l'envers, le capital d'une entité isolée étant dominé
par un sinistre unique, celui du secteur par l'accumulation.}

% ---------------------------------------------------------------------
\section*{Chapitre 11. Robustesse et incertitude du chiffre reporté}
\addcontentsline{toc}{section}{Chapitre 11. Robustesse et incertitude du chiffre reporté}
\reperes{p.~130 à 142. 11.1 la table de robustesse (p.~131), 11.2 limites du domaine et de ce
travail (p.~132), 11.3 trois étages (p.~132), 11.4 le moteur contrôlé (p.~134), 11.5 VaR prédictive
et robuste (p.~136), 11.6 le défaut de calibration (p.~139), 11.7 défaillances simultanées et
additivité des coûts (p.~140).}

\enbref{La robustesse en une table, la séparation entre ce que la donnée interdit à quiconque et ce
que ce travail a choisi, les trois étages d'incertitude sur une même échelle, un moteur contrôlé
par trois chemins, le choix du chiffre reporté et le registre des écarts chiffrés plutôt que
corrigés.}

\begin{parcoeur}[À savoir par c\oe ur : incertitude et moteur]
\begin{fichetab}
Verdicts & robuste : l'effet DORA, la priorité P1 et P4, le rôle de l'asymétrie, l'absence de
  Hawkes, les lois du socle hors échantillon ; voulu : le capital réagit à la conformité
  ($\times3{,}34$) ; assumé : la distribution large, le niveau de sévérité hors échantillon,
  l'additivité des coûts \\
Deux natures & limites du domaine (VERIS $0{,}12\,\%$, $1{,}19\,\%$, un incident ; coût
  multi-piliers inobservé ; sévérité en euros seulement dans les bases bancaires) contre limites
  de ce travail (élicitation sans réponse, gel, source citée, deux écarts publiés) \\
Même échelle & charge annuelle : paramètre $[3\,042\,;27\,152]$, $\times8{,}93$ ; identification
  $[6\,858\,;8\,697]$, $\times1{,}27$ ; modèle, dépendance $[5\,322\,;9\,806]$, $\times1{,}84$ ;
  modèle, famille de queue $[4\,330\,;39\,901]$, $\times9{,}21$ \\
Bande reportée & paramètre plus identification ; le modèle en sort, publié comme axe de
  sensibilité, parce qu'un choix de famille ne se moyenne pas \\
Famille & lognormale $4\,330$ contre GPD à $\xi=0{,}90$ $24\,016$ ($\times5{,}5$) ; cascade
  $8\,318$, entre gaussienne et Gumbel ; plafond de $40$ M\euro{} : $495$ \\
Hawkes & une cascade sans auto-excitation donne $n=0{,}480$ au jour contre $0{,}551$ mesuré ;
  pic à $3{,}3$~h ; $0{,}551$ tombe vers zéro si les co-occurrences du jour passent en exogène \\
Moteur & Monte-Carlo $6\,049$ et $20\,188$ ; Fourier $5\,945$ et $19\,890$ ($-1{,}7$ et
  $-1{,}5\,\%$, soit $0{,}75$ et $0{,}59$ erreur type) ; perte unique au second ordre $-9{,}9$ et
  $-7{,}7\,\%$, au premier $-17{,}5$ et $-26{,}0\,\%$ \\
\end{fichetab}
\end{parcoeur}

\begin{parcoeur}[À savoir par c\oe ur : chiffre reporté et écarts]
\begin{fichetab}
Six postures & point $657$ ; prédictive $645\pm7$ ($648$ ailleurs, le même nombre) ; robuste à
  $\beta=90\,\%$ $939\pm13$, à $95\,\%$ $1\,037\pm9$, à $99\,\%$ $1\,247\pm24$ ; pire cas de
  modèle ($\xi=0{,}90$) $1\,331$ \\
Robuste à $95\,\%$ & c'est la borne haute de l'IC à $90\,\%$ déjà publié, $1\,037$ : la même
  quantité sous deux noms ; la prédictive ne s'écarte du point qu'en profondeur ($+4\,\%$ à
  $99{,}9\,\%$) \\
Posture retenue & le point avec sa bande, pour trois raisons : réglementaire, de mesure, de
  précision \\
$p_u$ gelé & $0{,}1509=88/583$ contre $91/582=0{,}1564$ ; quantile $662{,}78\to678{,}60$
  ($+2{,}4\,\%$, $15{,}8$ M\euro) \\
Deux colonnes & défauts involontaires, tous sous-estimants : $p_u$, couverture $86\,\%$, dérive
  (au moins $+24{,}7\,\%$) ; choix posés, tous prudents : $a=0{,}60$ contre $0{,}22$,
  $\varphi=9{,}20$ contre $2{,}24$ ; périmètre de sens opposés : attritionnel absent, sévérité non
  plafonnée \\
Multi-piliers & $62{,}31\,\%$ des sinistres à l'état non conforme contre $31{,}15\,\%$ à l'état
  conforme ; $1{,}931$ pilier en moyenne contre $1{,}380$ \\
Additivité des coûts & perte $(\sum_jX_j)\,k^{\theta-1}$ ; pour $\theta\in[0{,}70\,;1{,}30]$,
  écart de $9\,706$ à $20\,487$ ($76\,\%$ de l'écart, quinze fois son bruit de $\pm734$) ;
  élasticités $0{,}95$ (NC) contre $0{,}39$ (C), facteur $2{,}42$ ; sens indéterminé \\
Plafond et saturation & le plafond amortit, l'interaction change de signe à $\theta=1$ ; à
  $\kappa=0{,}5\,\%$, $96{,}2$ M\euro{} pour tous les $\theta$ ; bande conjointe $177{,}8$ sous la
  somme des deux, $203{,}7$ \\
\end{fichetab}
\end{parcoeur}

\begin{question}{\og Vos limites ne sont-elles pas des insuffisances de votre travail ?\fg{}}
\rep{Il faut en séparer deux natures. Les unes tiennent à ce qu'aucune source ne permet de calibrer,
à personne : la direction n'est enregistrée nulle part, VERIS en porte le champ et le remplit dans
$0{,}12\,\%$ des cas. Les autres sont des choix de ce travail : le gel, la source citée,
l'élicitation restée sans réponse, deux écarts publiés plutôt que corrigés. Une limite du domaine se
corrige par le registre, une limite de ce travail en y revenant.}
\appui{\S~11.2, p.~132.}
\end{question}

\begin{question}{\og Des trois étages d'incertitude, lequel domine ?\fg{}}
\rep{Sur une même grandeur, la charge annuelle, le paramètre et le modèle sont du même ordre,
facteurs $8{,}93$ et $9{,}21$, et tous deux écrasent l'identification, $1{,}27$. Le facteur $2{,}5$
souvent cité porte sur le quantile d'un sinistre : le comparer tel quel à l'étage de modèle
comparerait deux échelles. La bande reportée cumule paramètre et identification ; l'étage de modèle
est publié comme axe de sensibilité.}
\appui{\S~11.3, p.~133.}
\end{question}

\begin{question}{\og Pourquoi sortir l'incertitude de modèle de la bande ?\fg{}}
\rep{Parce que les deux premiers étages sont statistiques, portables par un intervalle, alors que
le troisième est épistémique : un choix de famille ne se moyenne pas. Sur l'axe de la queue, il
pousse $\xi$ à une valeur posée hors de l'intervalle bootstrap, et le cumuler mettrait un scénario
et une borne de confiance dans le même objet. Il reste publié deux fois, comme pire cas et comme
bande de dépendance.}
\appui{\S~11.3, p.~133 et 134.}
\end{question}

\begin{question}{\og Et si vous aviez pris une autre méthode ?\fg{}}
\rep{Chaque brique a des alternatives, fréquence, sévérité, dépendance, mesure, agrégation, et je
ne prétends pas détenir la méthode. Les concurrents les plus sérieux sont testés et l'effet du
choix est borné. Le niveau bouge dans la bande ; l'ordre d'amorce des piliers, P1, P4, P2, P3, P5,
ne bouge pas d'une famille à l'autre.}
\appui{\S~11.3, encadré p.~134.}
\end{question}

\begin{question}{\og Pourquoi pas un processus de Hawkes ?\fg{}}
\rep{Il a été testé, sur des données datées, contre les noyaux de la littérature : l'excitation est
intra-journalière, pic à $3{,}3$~heures, et le ratio de branchement de $0{,}551$ tombe vers zéro dès
que les co-occurrences du même jour passent en exogène. Surtout, une cascade sans aucune
auto-excitation donne déjà $0{,}480$ au jour : les deux ne sont pas distinguables à la résolution
disponible. La cascade est retenue par parcimonie, puisqu'elle porte en plus la structure entre
piliers. Équivalence observationnelle, pas réfutation.}
\appui{\S~11.1, p.~131 ; \S~11.3, p.~134.}
\end{question}

\begin{question}{\og Comment savez-vous que votre moteur de simulation ne se trompe pas ?\fg{}}
\rep{Il est contrôlé par des chemins qui ne partagent ni code ni aléa : l'inversion de Fourier de la
loi composée s'en écarte de $-1{,}7$ et $-1{,}5\,\%$, soit moins d'une erreur type, donc un écart non
résolu ; l'approximation par perte unique reste à moins de $10\,\%$ au second ordre. Limite du
contrôle : les chemins partagent les lois et les tables de cascade, ils vérifient l'agrégation,
jamais la vérité du modèle.}
\appui{\S~11.4, p.~135 et 136.}
\end{question}

\begin{question}{\og Quel nombre reportez-vous : le point, le prédictif ou le robuste ?\fg{}}
\rep{Le point avec sa bande, pour trois raisons. Réglementaire : une borne haute sur un ensemble
d'ambiguïté est un supremum sur une famille de lois, pas un quantile de la loi de perte. De mesure :
la prédictive vaut $645$ contre $657$, deux fois le bruit. De précision : les postures robustes
sont les moins reproductibles. Et la robuste à $95\,\%$ est déjà publiée, c'est la borne haute de
l'intervalle, $1\,037$.}
\appui{\S~11.5, p.~137 et 138.}
\end{question}

\begin{question}{\og \og Le quantile d'un objet incertain est mauvais\fg{} : que répondez-vous ?\fg{}}
\rep{Que la remarque est juste, mais pas par le mécanisme attendu. Intégrer l'incertitude de
paramètre avant le quantile ne déplace pas le point à $99{,}5\,\%$ ($645$ contre $657$) ; l'écart
n'apparaît qu'en profondeur, $+4\,\%$ à $99{,}9\,\%$. La fragilité est dans la largeur de la bande,
pas dans le point. D'où la règle : jamais le point nu, toujours avec sa bande.}
\appui{\S~11.5, p.~136.}
\end{question}

\begin{question}{\og Vous savez que votre taux de dépassement est incohérent avec votre seuil, et vous ne corrigez pas ?\fg{}}
\rep{L'écart vaut $2{,}4\,\%$ sur le quantile, un quarantième de son intervalle à $90\,\%$. Rejouer
la chaîne réinjecterait un bruit de Monte-Carlo du même ordre : on remplacerait un écart connu et
signé par un déplacement non attribuable. L'écart est déclaré, chiffré et signé : anti-conservateur
pour la solvabilité, minorant pour la thèse. Seul le premier engage.}
\appui{\S~11.6, p.~139.}
\end{question}

\begin{question}{\og Vous avez gelé votre calibration en cours de route. N'est-ce pas un aveu ?\fg{}}
\rep{C'est une décision de conduite de projet, datée. Au-delà d'un certain volume de résultats
publiés, un modèle qui bouge encore invalide en silence ce qui est écrit. Le gel sépare la
correction d'erreur, restée autorisée et pratiquée, de la recalibration, qui ne l'était plus, et
son prix est publié : les écarts connus sont chiffrés plutôt que corrigés.}
\appui{\S~11.6 ; chapitre~14, \S~14.1 p.~151.}
\end{question}

\begin{question}{\og Vos erreurs vont-elles toutes dans le même sens ?\fg{}}
\rep{Oui, et c'est pourquoi elles sont rangées en deux colonnes. Les trois défauts involontaires
d'estimation sous-estiment tous le capital ; les deux choix délibérés, la charge systémique et la
surdispersion, le majorent. La prudence du capital publié vient donc de choix déclarés, pas de ses
écarts d'estimation. Seule l'additivité des coûts a un sens indéterminé.}
\appui{\S~11.6, tableau p.~139.}
\end{question}

\begin{question}{\og Les défaillances simultanées de plusieurs piliers ne sont pas traitées.\fg{}}
\rep{La prémisse est fausse : c'est la sortie même du modèle, et à l'état non conforme elle est
majoritaire, $62{,}31\,\%$ des sinistres touchant plus d'un pilier contre $31{,}15\,\%$ à l'état
conforme. La loi est exacte, calculée par énumération. Ce qui reste une hypothèse, c'est la façon
dont les coûts de ces piliers s'agrègent.}
\appui{\S~11.7, p.~140.}
\end{question}

\begin{question}{\og Quelle est l'hypothèse la plus lourde de votre modèle ?\fg{}}
\rep{L'additivité des coûts au sein d'un sinistre multi-piliers. Relâchée par un exposant, elle
déplace l'écart entre états de $9\,706$ à $20\,487$~M\euro, $76\,\%$ de l'écart publié, et elle
mord plus à l'état non conforme, élasticité $0{,}95$ contre $0{,}39$. Son sens ne se déclare pas :
la mutualisation de la remédiation et la saturation de la capacité tirent en sens contraire. Ce qui
tient : le signe et l'ordre de grandeur de l'écart, sur toute la plage.}
\appui{\S~11.7, p.~141.}
\end{question}

% ---------------------------------------------------------------------
\section*{Chapitre 12. La donnée manquante comme livrable, hiérarchisé}
\addcontentsline{toc}{section}{Chapitre 12. La donnée manquante comme livrable}
\reperes{p.~143 à 145. 12.1 le registre (p.~143), 12.2 ce que l'information vaut (p.~144), 12.3
quelle dépendance documenter (p.~145).}

\enbref{La limite d'identification se retourne en spécification de reporting, chiffrée par la
valeur d'information et hiérarchisée par dépendance : c'est le seul endroit du mémoire où une
limite produit un livrable.}

\begin{parcoeur}
\begin{fichetab}
Le registre & horodater la survenance au mois ou à la semaine ; catégoriser par domaine de
  contrôle ; consolider une cause commune en un événement à $N$ victimes ; consigner tous les
  piliers qui ont cédé ($62{,}31\,\%$ de multi-piliers sous $W$, $50{,}77\,\%$ sous $W^\top$) \\
Instruments & le registre d'information du règlement d'exécution 2024/2956 sert la troisième
  exigence ; la notification des incidents majeurs porte les trois autres, par le format de
  l'horodatage et la taxonomie \\
VERIS & le champ du contrôle défaillant et la chronologie en quatre étapes existent : rendre
  obligatoires deux champs normalisés, pour un coût de conception nul \\
Valeur & $L(\mathcal V)=\sup_{\mathcal V}\mathrm{SCR}-\inf_{\mathcal V}\mathrm{SCR}$ ;
  $V(\mathcal I)=1-L(\mathcal W_{\mathcal I})/L(\mathcal W)\in[0,1]$ ; monotone et sans
  dimension \\
$66\,\%$ & si l'on retenait le sens que sept rapports documentent sur six dépendances, alors que
  deux seulement sont tranchées : la valeur d'une information, pas un acquis \\
$28{,}9\,\%$ & la seule dépendance gouvernance-incidents, que le corpus ne couvre pas ; zéro pour
  tests-tiers \\
Amplitude & même les dix sens connus, une bande subsiste ($656$ M\euro) : l'identification
  partielle est structurelle \\
\end{fichetab}
\end{parcoeur}

\begin{question}{\og Que répondez-vous à un superviseur ?\fg{}}
\rep{Une recommandation chiffrée et opposable : rendre obligatoires, dans la notification des
incidents majeurs sous DORA, deux champs \emph{déjà normalisés} dans VERIS, le domaine de contrôle
défaillant et la chronologie. Le coût de conception est nul. Et un superviseur qui ne pourrait
exiger qu'une chose sait laquelle : la dépendance entre gouvernance et gestion des incidents, qui
lèverait $28{,}9\,\%$ de l'ambiguïté de capital à elle seule.}
\appui{chapitre~12, p.~143 à 145.}
\end{question}

\begin{question}{\og $66\,\%$ de resserrement : c'est un résultat ?\fg{}}
\rep{Non, c'est la valeur d'une information de cette étendue. Les sept premiers rapports
documentent le sens de six dépendances, mais le mémoire n'en revendique que deux paire par paire.
Retenu sur les six, ce sens ferait tomber la largeur de $66\,\%$ : c'est l'ordre de grandeur de ce
qu'un registre bien spécifié apporterait, pas un acquis.}
\appui{\S~12.2, p.~144.}
\end{question}

\begin{question}{\og DORA ne prévoit-il pas déjà ce registre ?\fg{}}
\rep{En partie. Le registre d'information porte les accords contractuels avec les prestataires,
donc la dépendance \emph{déclarée} : il sert la consolidation des causes communes. Il ne porte pas
l'ordre des défaillances. La notification des incidents majeurs, elle, existe : il s'agit d'en
fixer le format d'horodatage et la taxonomie, pas de créer une collecte nouvelle.}
\appui{\S~12.1, p.~143.}
\end{question}

\begin{question}{\og Si l'on connaissait tous les sens, auriez-vous un point ?\fg{}}
\rep{Non : connaître le sens d'un lien ne donne pas sa force. Même les dix directions établies, une
bande d'amplitude subsisterait. L'identification partielle n'est pas un état provisoire que la
donnée finirait par lever, et c'est une raison de plus de ne pas promettre à un superviseur plus
qu'un registre ne peut tenir.}
\appui{\S~12.2, encadré p.~144 ; chapitre~8, \S~8.10.}
\end{question}

\begin{question}{\og Votre registre est-il réaliste ?\fg{}}
\rep{Les praticiens interrogés le jugent réaliste, sous réserve de règles de qualification
précises, le classement par domaine pouvant être subjectif, et avec un point de vigilance sur la
confidentialité. L'une décrit même un usage déjà établi : un tableau de bord des incidents par
prestataire, sur douze mois, dans une analyse de risque au titre de DORA.}
\appui{annexe~C, \S~C.7, p.~179 à 182.}
\end{question}

% ---------------------------------------------------------------------
\section*{Chapitre 13. Conclusion}
\addcontentsline{toc}{section}{Chapitre 13. Conclusion}
\reperes{p.~146 à 149. 13.1 ce qui est établi, 13.2 ce qui est démontré non mesurable, 13.3 pour
le régulateur, 13.4 pour l'entité, 13.5 limites, 13.6 ouvertures.}

\enbref{Ce qui est établi (le chiffre, le mécanisme, un acquis de méthode), ce qui est démontré non
mesurable (la direction, convertie en bornes), ce qui en découle pour le régulateur et pour
l'entité, puis les limites et trois ouvertures.}

\begin{parcoeur}
\begin{fichetab}
Fréquence & seule, $10\,377$, que la forme fermée de la perte unique retrouve à moins de $1\,\%$ \\
Méthode & le dispositif de vérification a trouvé trois valeurs fausses qu'aucune relecture
  n'avait vues ; il ne vérifie ni qu'un nombre est au bon endroit, ni ce que la phrase en fait \\
Détenir, transférer & rapport du coût de détention au prix du transfert : $2{,}78$ au taux
  révisé, $3{,}51$ au taux historique ; ses deux termes sont proportionnels au même SCR ; l'effet
  de forme est borné à $+36\,\%$ sur une plage de SCR d'un facteur quatre \\
Limites & élicitation sans réponse ; niveaux de propagation posés ; additivité des coûts
  ($0{,}95$ contre $0{,}39$) ; l'ordre entre P1 et P4 dépend de la direction \\
Ouvertures & l'élicitation ; le pilier des tiers sur données clients ; une criticité au sens de
  l'audit, le risque résiduel après contrôles \\
\end{fichetab}
\end{parcoeur}

\begin{question}{\og Qu'apporte votre travail à la science actuarielle ?\fg{}}
\rep{Un mécanisme de dépendance dirigée dont trois propriétés sont démontrées, dont celle-ci : une
matrice et sa transposée, indiscernables pour les statistiques disponibles, donnent un capital et
une décision différents. Une frontière d'identification établie sur données réelles. Et une méthode
importée de l'économétrie, l'identification partielle, appliquée à une exigence de capital.}
\appui{\S~13.1 et 13.2, p.~146 et 147.}
\end{question}

\begin{question}{\og Quelle est la faiblesse de votre travail ?\fg{}}
\rep{À dire sans détour, parce que le jury le sait et que l'esquive coûte plus que l'aveu : le
niveau absolu du capital n'est pas une mesure, et l'ordre entre les deux piliers de tête dépend
d'une direction que je ne sais pas identifier. Ce que je sais dire, c'est que ces deux-là dominent
les trois autres, et combien l'ignorance coûte.}
\appui{\S~13.5, p.~148.}
\end{question}

\begin{question}{\og Pourquoi votre rapport entre détention et transfert serait-il plus robuste que le niveau ?\fg{}}
\rep{Pour une raison qui est une identité : ses deux termes sont proportionnels au même SCR, qui se
simplifie, donc un changement d'échelle pur le laisse inchangé. Il vaut $2{,}78$ au taux révisé et
$3{,}51$ au taux historique. Reste un effet de forme, la descente d'échelle déplaçant fréquence et
sévérité séparément, borné à $+36\,\%$ : plus robuste que le niveau, sans être invariant.}
\appui{\S~13.4, p.~148.}
\end{question}

% ---------------------------------------------------------------------
\section*{Chapitre 14. Enseignements du stage (version ENSAE)}
\addcontentsline{toc}{section}{Chapitre 14. Enseignements du stage}
\reperes{p.~151 à 153, dans la seule version ENSAE. 14.1 le cadre, 14.2 les enseignements de
l'ENSAE, 14.3 l'encadrement et les tiers, 14.4 ce que j'en retiens, 14.5 ce que je ferais
autrement.}

\enbref{Le recul sur la mission, à la première personne et sans aucun nombre : quatre contraintes de
cadre, les enseignements mobilisés, la part de l'encadrement et des tiers, et ce que je ferais
autrement.}

\begin{parcoeur}
\begin{fichetab}
Quatre contraintes & deux livrables, le mémoire et l'analyse LUCY co-écrite avec Hugo Rapior ; une
  donnée sous licence, où la traçabilité remplace la reproductibilité complète ; des entités
  anonymisées ; le gel de la calibration \\
Enseignements & valeurs extrêmes, modèles à latente gaussienne, processus stochastiques (Markov,
  branchement), statistique inférentielle, Monte-Carlo, jeux coopératifs (Shapley), assurance
  non-vie ; les deux qui ont le plus pesé : l'économétrie de l'identification, et $\xi$ comme
  condition d'existence des moments \\
Encadrement & Hugo Rapior, point hebdomadaire en diapositives ; Caroline Hillairet, sur le
  traitement de l'incertitude (quantile prédictif, sensibilité à la propagation) et sur Hawkes \\
Tiers & un second codage sollicité, non reçu ; trois praticiens ont répondu, non nommés \\
Autrement & le dispositif de vérification dès le début ; le gel de la calibration plus tôt \\
\end{fichetab}
\end{parcoeur}

\begin{question}{\og Qu'est-ce que l'ENSAE vous a apporté qui vous a servi ?\fg{}}
\rep{Surtout deux choses. L'économétrie de l'identification : reconnaître qu'un paramètre n'est
pas identifié, plutôt que mal estimé, a changé la méthode. Et la lecture de l'indice de queue comme
condition d'existence des moments, qui permet d'écarter des variantes par un argument d'existence
plutôt que de degré. La programmation scientifique a compté autant : tenir une chaîne de scripts
reproductible pendant des mois est un problème d'ingénierie.}
\appui{\S~14.2, p.~152.}
\end{question}

\begin{question}{\og Quelle a été la part de votre maître de stage et de votre tutrice ?\fg{}}
\rep{Leurs demandes n'ont pas porté sur des calculs, et c'est ce qui en a fait la valeur. Hugo
Rapior m'a fait retirer des lectures que la mesure ne soutenait pas et préciser un vocabulaire qui
masquait une borne prise à l'envers. Caroline Hillairet a porté sur le traitement de l'incertitude
et a produit deux sections, le quantile prédictif et la sensibilité aux probabilités de
propagation ; l'écartement du processus de Hawkes a été testé contre ses propres noyaux.}
\appui{\S~14.3, p.~152.}
\end{question}

\begin{question}{\og Que feriez-vous autrement ?\fg{}}
\rep{Deux choses. Mettre le dispositif de vérification en place dès le début, et non à
mi-parcours, ce qui a obligé à rattacher après coup chaque section à ses scripts. Et geler la
calibration plus tôt, des semaines de recalibrations partielles ayant fait bouger des nombres déjà
écrits.}
\appui{\S~14.5, p.~153.}
\end{question}

\begin{question}{\og Qu'est-ce qui distingue une mission d'un exercice académique ?\fg{}}
\rep{En conseil, un résultat doit pouvoir être contrôlé par quelqu'un qui n'a pas fait le calcul.
Rattacher chaque nombre publié à la sortie du script qui le produit répond à cette exigence, et
cela a représenté une part substantielle du temps.}
\appui{\S~14.4, p.~153.}
\end{question}

% ---------------------------------------------------------------------
\section*{Les annexes A à G}
\addcontentsline{toc}{section}{Les annexes A à G}
\reperes{A démonstrations p.~156 ; B adaptations par pilier p.~160 ; C élicitation p.~173 ; D table
des paramètres p.~183 ; E notations p.~185 ; F pistes p.~187 ; G inventaire de l'IA p.~191. Version
Institut : quatre pages plus tôt.}

\subsection*{Annexe A. Les démonstrations à savoir refaire au tableau}

\begin{parcoeur}
\begin{fichetab}
Forme réduite & $\mathbb E\,\mathcal M_n(a,k)=(W^n)_{ak}$ ; la série de Neumann converge vers
  $(I-W)^{-1}$ si et seulement si $\rho(W)<1$ ; $\mathbf 1\{k\in S_a\}\le\sum_n\mathcal M_n(a,k)$
  donne la majoration de $\mathbb E|S_a|$ \\
Sous-criticité & $\rho(W)\le\lVert W\rVert_\infty=\max_j\sum_kW_{jk}\le g<1$ \\
Loi exacte & deux événements sur des arêtes disjointes, donc indépendants : $B(A)$ pour les
  fuites, $R(A)$ par récurrence sur les sous-ensembles \\
Monotonie & couplage par uniformes communes : $U_{jk}<W_{jk}\Rightarrow U_{jk}<W'_{jk}$, donc
  $S_a(W)\subseteq S_a(W')$ presque sûrement \\
Ordre & pour des chaînes de longueur deux, l'invariance équivaut à
  $\mathrm{TRANS}_{ij}=\mathrm{TRANS}_{ji}$ \\
Co-occurrence & contre-exemple de la fourche : $W_{12}=W_{13}=w$ ; sous $W$, la paire $\{2,3\}$
  co-occurre avec la probabilité $rw^2$, sous $W^\top$ jamais, pour la même partie symétrique \\
Shapley & poids $\frac{s!(n-s-1)!}{n!}$ : moyenne sur les $n!$ ordres ; la somme des
  contributions le long d'un ordre est télescopique, d'où $\sum_i\varphi_i=v(\text{tous})$ \\
\end{fichetab}
\end{parcoeur}

\begin{question}{\og Montrez que $W$ et $W^\top$ peuvent être distingués.\fg{}}
\rep{Trois piliers, $W_{12}=W_{13}=w$, tout le reste nul. Sous $W$, une amorce en~1 touche~2 et~3
ensemble avec la probabilité $w^2$ : la paire $\{2,3\}$ co-occurre. Sous $W^\top$, les arêtes vont
de~2 et de~3 vers~1, et aucun sinistre ne touche~2 et~3 ensemble. Même partie symétrique, lois
différentes : c'est la structure en V de l'identification des graphes causaux.}
\appui{annexe~A, p.~158.}
\end{question}

\begin{question}{\og Pourquoi les parts de Shapley somment-elles exactement à $100\,\%$ ?\fg{}}
\rep{C'est l'efficience. Le poids de Shapley fait de $\varphi_i$ la moyenne, sur les $n!$ ordres
d'arrivée, de la contribution marginale de $i$. À ordre fixé, la somme des contributions est
télescopique et vaut $v(\text{tous})-v(\emptyset)$ ; en moyennant, $\sum_i\varphi_i=v(\text{tous})$.
Limite : c'est une partition exacte du surcoût, pas une prévision de ce qu'on gagnerait en retirant
un seul joueur.}
\appui{annexe~A, p.~158 et 159.}
\end{question}

\subsection*{Annexe B. Adaptations par pilier}

\begin{parcoeur}
\begin{fichetab}
Queue par pilier & $\xi$ de $0{,}87$ à $1{,}37$ selon la catégorie bâloise ; $111$ observations
  pour la catégorie TIC, sous le seuil de $150$ ; $120$ assignations : capital $+55$ à $+96\,\%$ ;
  pilier de tête dominant dans $83\,\%$ \\
P4 en choc commun & intensité bornée dans $[0\,;0{,}68]$ ; $3{,}7$ piliers emportés d'un coup,
  $+19\,\%$ de capital ; au moins trois piliers touchés : $0{,}22$ en cascade, $0{,}90$ en choc
  commun ; co-occurrence symétrique, donc identifiée \\
Seuil et copule & $K_4=1{,}03$ pour $p_4=0{,}151$ ; $\lambda_U=2-2^{1/\theta}=0{,}53$ à
  $\theta=1{,}8$ ; cinq piliers avec P4 : $0{,}17$ en gaussien, $0{,}39$ en Gumbel \\
Temps d'arrêt & $\Prob(\tau\le T)=2\,\Phi(-z^\star/\sqrt T)$, par le principe de réflexion
  (Markov fort et symétrie, pas la martingalité) ; $0{,}14$ à $\rho=0{,}50$, $0{,}23$ à
  $\rho=0{,}75$ \\
Ancrages & ESAs : $29\,\%$ des $3\,383$ incidents majeurs de 2025 d'origine tierce ; BCE : plus
  de $30\,\%$ du budget des grandes banques à dix fournisseurs ; $19$ prestataires critiques
  désignés \\
Détection & VCDB ($n=278$) : $47\,\%$ le jour même, $8\,\%$ en semaines, $45\,\%$ en mois ; canal
  seul de $1\,775$ à $8\,052$ ; le multiplicateur $0{,}85$ et $1{,}20$ correspond à
  $\eta\approx0{,}36$ \\
Coin supérieur & sur $65$ paires emboîtées, relâcher un canal ne baisse jamais le capital, graine
  par graine ; faux trajectoire par trajectoire ($16{,}2\,\%$ et $30{,}7\,\%$) ; un coin peut être
  infaisable \\
Paires & fréquence et détection $1\,739\pm708$, fréquence et accumulation $1\,666\pm354$,
  fréquence et propagation $1\,182\pm444$ ($86\,\%$ de l'ordre~2) ; propagation et accumulation
  $-330\pm242$, substituts \\
Retour & portage $848$ M\euro/an à $6\,\%$ ($672$ à $4{,}75\,\%$) ; sinistralité évitée $3\,247$,
  $3{,}8$ fois plus ; retour de $6$ à $35$ ans, un majorant \\
\end{fichetab}
\end{parcoeur}

\begin{question}{\og Votre retour sur investissement de $6$ à $35$ ans est-il un plancher ?\fg{}}
\rep{C'est un majorant de la durée. Il ne monétise que le portage du capital libéré, $848$~M\euro{}
par an ; la sinistralité évitée, $3\,247$ par an, en vaut $3{,}8$ fois plus. Minorer le bénéfice
majore la durée : le compte réel est plus court, jamais plus long. Additionner les deux donnerait
$1$ à $7$ ans, mais mêlerait un flux de résultat et un coût du capital.}
\appui{annexe~B, \S~B.4, p.~170.}
\end{question}

\begin{question}{\og Un test de résistance sur vos quatre canaux ?\fg{}}
\rep{Il se réduit à une évaluation : l'état non conforme est le coin supérieur du pavé, relâcher un
canal ne baissant jamais le capital sur les $65$ paires emboîtées, graine par graine. Deux
réserves : la propriété vaut pour le quantile, pas trajectoire par trajectoire, et deux canaux
étant des bornes posées, le coin est un majorant sur un pavé déclaré, pas la description d'une
entité.}
\appui{annexe~B, \S~B.4, p.~166 et 167.}
\end{question}

\begin{question}{\og Pourquoi le pilier des tiers échappe-t-il à votre limite d'identification ?\fg{}}
\rep{Parce que son accumulation est un choc commun : P4 et un autre pilier tombent ensemble, sans
direction. C'est une co-occurrence symétrique, donc la partie identifiée du modèle, celle que la
donnée voit, MOVEit et la concentration de l'informatique en nuage.}
\appui{annexe~B, \S~B.2, encadré p.~161.}
\end{question}

\subsection*{Annexe C. Le protocole d'élicitation, diffusé et resté sans réponse}

\begin{parcoeur}
\begin{fichetab}
Cooke & quantiles $5$, $50$ et $95\,\%$ ; calibration sur des graines de réponse connue ;
  information ; poids égal au produit des deux, nul sous un seuil \\
Questionnaire & troisième version, seize questions dont dix graines (G1 à G10) et six liens (B1 à
  B6) ; valeurs vraies non publiées, pour ne pas brûler l'instrument ; graine $\xi=0{,}90$
  retirée (valeur posée, paramètre estimé) \\
Convention & $j$ vers $k$ ; une inversion de lettres ferait saisir $W^\top$, sans contrôle
  numérique possible \\
Diffusion & une commission de l'Institut et des praticiens du cabinet : aucune réponse, aucun
  panel \\
Pourquoi c'est tenable & à un ou deux experts, la pondération de Cooke dégénère ; un panel biaisé
  déplacerait le capital de $777$ M\euro{} sans signal, quand les bornes ne bougent pas \\
Relectures & trois réponses : zéro, deux et deux désaccords ; la cause commune avancée sans y être
  invité ; manquent les déclencheurs transversaux (gestion des changements, des
  vulnérabilités) \\
Coût & sans panel, P1 et P4 restent à égalité ($43{,}6$ contre $37{,}7\,\%$) ; reprise sans coût
  de développement avec une dizaine de répondants \\
\end{fichetab}
\end{parcoeur}

\begin{question}{\og Pourquoi n'avez-vous pas d'avis d'experts sur cette matrice ?\fg{}}
\rep{Le protocole a été préparé, puis diffusé à une commission de l'Institut et à des praticiens du
cabinet, et il est resté sans réponse : aucun panel n'a pu être constitué. Ce n'est donc pas un
choix de ne pas essayer. Le mémoire chiffre en outre ce qu'un essai raté aurait coûté : un panel
biaisé déplace le capital de $777$~M\euro{} sans qu'aucun signal ne l'indique, quand les bornes
restent invariantes. Une élicitation faiblement étalonnée échangerait une ignorance mesurée contre
une erreur invisible.}
\appui{annexe~C, \S~C.4 et C.5, p.~177 et 178.}
\end{question}

\begin{question}{\og Vous refusez le jugement d'expert, mais vous codez vous-même vos rapports. Contradiction ?\fg{}}
\rep{L'asymétrie est assumée et justifiée par trois différences. Le codage est auditable : le
protocole est écrit et les récits sont publics, n'importe qui peut recoder. Son biais principal est
mesuré par un test contre une loi nulle exacte. Et il produit un ordre partiel, que la suite borne,
là où une élicitation produirait des coefficients entrant directement dans le capital. Reste une
limite déclarée : le codeur unique.}
\appui{annexe~C, \S~C.5, p.~178.}
\end{question}

\begin{question}{\og Que vous ont dit les praticiens ?\fg{}}
\rep{Le premier ne contredit rien, et son exemple admet deux lectures opposées, ce qui illustre la
frontière. Le deuxième conteste deux phrases et avance de lui-même la cause commune : c'est la
corroboration la plus nette de la frontière. Le troisième répète les deux désaccords. Aucun ne
distingue un ordre d'une cause commune ; les retours sont déclarés sans être traités, les prendre
en compte supposant de recalibrer.}
\appui{annexe~C, \S~C.7, p.~179 à 182.}
\end{question}

\subsection*{Annexes D et E. Les tables}

\begin{parcoeur}
\begin{fichetab}
Trois natures & calibré, posé, semi-ancré ; la direction $u_{ij}$ est non identifiée \\
Forte sensibilité & seulement $\xi$, calibré, et les leviers voulus : $g$, le score de
  conformité, le multiplicateur de fréquence \\
Posés non leviers & copule ($\Delta$ de moins de $1{,}4\,\%$ sur $\theta\in[1,4]$), $\gamma$, seuil
  de crise $-2{,}5$, probabilités d'état, matérialité : influence faible à modérée \\
Notations & $g$ mesure une criticité, jamais un montant ; $\mathrm{VaR}$ et $\mathrm{TVaR}$ pour
  la charge annuelle $L$, $q$ et $\bar q$ pour un sinistre $X$ \\
\end{fichetab}
\end{parcoeur}

\begin{question}{\og Qu'est-ce qui est calibré, qu'est-ce qui est posé ?\fg{}}
\rep{Calibrés : la queue de sévérité, l'échelle, la fréquence, l'élasticité à la taille. Posés : la
matrice et le gain de propagation, la surdispersion du moteur et la charge systémique, tous deux
prudents, le plafond, les multiplicateurs de matérialité. Semi-ancrés : le multiplicateur de
fréquence, la charge $\gamma$, les probabilités d'état. Aucune conclusion ne repose sur un paramètre
à la fois non calibré et déterminant hors de son rôle attendu.}
\appui{annexe~D, p.~183 et 184.}
\end{question}

\subsection*{Annexe F. Pistes pour de futurs mémoires}

\begin{parcoeur}
\begin{fichetab}
Exécutables & exécuter l'élicitation préparée ; documenter la dépendance gouvernance-incidents
  ($28{,}9\,\%$) ; le second codage en aveugle ; modéliser la dérive d'échelle mesurée \\
Donnée d'entité & calibrer l'amplitude des liens, pas seulement leur sens ; passer de l'échelle du
  secteur à une entité réelle \\
Structure & remplacer l'exposant d'additivité par un modèle de coût à contrainte de capacité ;
  étendre la cascade entre entités par le prestataire partagé, dont le graphe est observé \\
\end{fichetab}
\end{parcoeur}

\subsection*{Annexe G. Inventaire de l'usage de l'IA générative}

\begin{parcoeur}
\begin{fichetab}
Usage & rédaction et reformulation du manuscrit ; écriture et relecture des scripts ; quatre
  agents de relecture automatique ; le dispositif de vérification lui-même \\
Vérification & $1\,885$ nombres vérifiables dans $18$ chapitres, confrontés à $117$ sorties de
  référence versionnées \\
Limite & huit erreurs sémantiques que le dispositif confirmait toutes (cinq de relation, un
  renvoi, un superlatif, un comptage) ; deux venaient d'un message écrit en dur dans un script ;
  règle : un superlatif ou un comptage se calcule \\
Non délégué & la problématique, le choix de la cascade, la décision de borner, les sources et
  leurs biais, la hiérarchie des limites, l'interprétation \\
\end{fichetab}
\end{parcoeur}

\begin{question}{\og Avez-vous utilisé l'intelligence artificielle générative ?\fg{}}
\rep{Oui, et c'est déclaré en fin de résumé et détaillé à l'annexe~G, comme le référentiel l'exige :
assistance à la rédaction, écriture et relecture des scripts, agents de relecture, dispositif de
vérification. Aucun résultat chiffré n'en sort : tout nombre publié vient de l'exécution d'un
script du dépôt, et un dispositif rattache chaque nombre à la sortie qui le produit. L'annexe dit
aussi ce que ce dispositif ne détecte pas.}
\appui{résumé p.~3 ; annexe~G, p.~191 et 192.}
\end{question}

\begin{question}{\og Comment garantissez-vous qu'aucun chiffre n'est faux ?\fg{}}
\rep{Je ne le garantis pas, et l'annexe~G dit pourquoi. Les $1\,885$ nombres vérifiables sont
confrontés aux sorties des scripts cités, et tous sont confirmés. Mais une relecture au filtre des
superlatifs et des comptages a trouvé huit erreurs sémantiques que le dispositif confirmait
toutes : le nombre sortait du bon script, seule la phrase était fausse. Un taux de confirmation ne
mesure pas la justesse du texte.}
\appui{annexe~G, \S~G.2 et G.3, p.~191 et 192.}
\end{question}

\begin{question}{\og Qu'est-ce que l'IA n'a pas fait ?\fg{}}
\rep{Tout ce qui engage : la problématique, le choix de la cascade plutôt que d'une copule, la
décision de ne pas poser de matrice arbitraire et de publier en bornes, le traitement de
l'élicitation restée sans réponse, la qualification des sources, la hiérarchie des limites,
l'interprétation. Et deux décisions vont contre le sens d'une rédaction assistée : réfuter
l'hypothèse qui donnait son titre au mémoire, et requalifier le chiffre central en borne
supérieure.}
\appui{annexe~G, \S~G.4, p.~192.}
\end{question}

% ---------------------------------------------------------------------
\section*{Questions transversales : la forme, les choix, la posture}
\addcontentsline{toc}{section}{Questions transversales : la forme, les choix, la posture}

\noindent Ces questions ne visent pas un chapitre mais les \emph{décisions} qui ont fait le mémoire,
et c'est là qu'un jury mesure la maturité. Chacune a une réponse en quelques phrases, et aucune ne
consiste à s'excuser.

\begin{question}{\og Votre mémoire fait $199$~pages, le référentiel recommande environ $70$~pages de corps.\fg{}}
\rep{Le référentiel écrit \og environ\fg{} et ne prévoit aucune sanction. Le volume vient des données et
de la robustesse : vingt pages de données, treize de robustesse, et chaque nombre publié rattaché au
script qui le produit. Ce qui pouvait sortir est sorti : les démonstrations standard, les pièces
justificatives et les compléments sont déposés sur Zenodo, sous DOI.}
\appui{notice des annexes, p.~155.}
\end{question}

\begin{question}{\og Pourquoi avoir déporté des annexes hors du document ?\fg{}}
\rep{Pour alléger la lecture sans perdre la traçabilité. Les démonstrations propres à la
contribution, celles de la cascade, restent dans le document ; sont sortis les théorèmes standard
des valeurs extrêmes, les pièces justificatives et les compléments de chapitres. Le tout est déposé
de façon pérenne, sous l'identifiant \texttt{10.5281/zenodo.22874524}, donc citable.}
\appui{notice, p.~155 ; annexe~A, p.~156.}
\end{question}

\begin{question}{\og Vous citez une source que personne ne peut recalculer. Est-ce sérieux ?\fg{}}
\rep{Elle porte la structure du risque, les parts par vecteur d'attaque, jamais son niveau. Son
statut est déclaré dans un encadré : consultée puis non conservée, donc non reproductible par un
script du dépôt. J'ai préféré assumer une citation datée plutôt que fabriquer une reproductibilité
que je n'ai pas. La retirer retirerait une décomposition, pas une incertitude.}
\appui{chapitre~4, \S~4.6, encadré p.~50.}
\end{question}

\begin{question}{\og Vous exposez beaucoup vos limites. Ce n'est pas vous desservir ?\fg{}}
\rep{Une limite déclarée et chiffrée est un actif ; la même limite trouvée par le jury est une
faute. Le mémoire a d'ailleurs réfuté sa propre thèse d'origine, celle qui lui donnait son titre, et
requalifié son chiffre d'entité en borne supérieure. Ce qui tient dans ce travail, ce n'est pas
l'assurance du niveau : tout ce qui est affirmé est démontré, mesuré, ou déclaré posé.}
\appui{chapitre~11, \S~11.2 ; chapitre~13, \S~13.5.}
\end{question}

\begin{question}{\og Qu'auriez-vous fait avec six mois de plus ?\fg{}}
\rep{Trois choses, écrites à l'annexe~F avec ce qui existe déjà pour chacune. Exécuter
l'élicitation, qui est prête et ne manque que de répondants. Faire coder le corpus en aveugle par un
second lecteur. Et, sur la structure, remplacer l'additivité des coûts par un modèle à contrainte de
capacité, puis étendre la cascade entre entités par le prestataire partagé, où le graphe est
\emph{observé} et où l'extension n'hériterait donc pas du problème d'identification.}
\appui{annexe~F, p.~187 à 190.}
\end{question}

\begin{question}{\og En une phrase, que faut-il retenir ?\fg{}}
\rep{Que la non-conformité à DORA a un coût en capital que la Formule Standard ne voit pas, un
facteur $3{,}34$ au secteur sur les quatre canaux, et que ce qu'on ne peut pas savoir, la direction
de la contagion, se borne et se chiffre au lieu de se poser.}
\appui{conclusion, p.~146 et 147.}
\end{question}

\laphrase{Je n'ai pas cherché à rendre le chiffre plus sûr qu'il n'est. J'ai cherché à rendre exact
ce que l'on peut en dire, et à chiffrer ce que l'on ne peut pas.}

\begin{cle}
\textbf{La règle qui vaut pour toutes ces réponses.} Aucune ne consiste à défendre un chiffre :
toutes disent ce que le chiffre vaut et ce qu'il ne vaut pas. C'est la posture que le \S~4.6 du
référentiel valorise explicitement, et c'est celle qui a produit ce mémoire.
\end{cle}

% =====================================================================
\newpage
\section*{Annexe~: le modèle en formules, du sinistre au capital}
\addcontentsline{toc}{section}{Annexe : le modèle en formules}

\begin{cle}
\textbf{À quoi sert cette page.} Si l'on demande d'écrire le modèle au tableau, c'est cela qu'il
faut savoir poser, dans cet ordre, sans notes. Chaque bloc porte son \textbf{statut}, et ce statut
est la moitié de la réponse~: \emph{estimé} sur données, \emph{posé} faute de source,
\emph{démontré}, ou \emph{hypothèse de construction non testée}. Confondre les quatre est la
seule faute qui ne se rattrape pas.
\end{cle}

\subsection*{1. La chaîne, en une ligne}

\[
\text{fréquence} \;\longrightarrow\; \text{cascade entre piliers} \;\longrightarrow\;
\text{coût du sinistre} \;\longrightarrow\; \text{charge annuelle} \;\longrightarrow\;
\text{capital}
\]

\noindent L'état de conformité n'ajoute aucune pénalité~: il \textbf{déplace quatre paramètres} de
cette chaîne, et rien d'autre.

\subsection*{2. La fréquence \quad\small\textcolor{navy}{statut~: estimée}\normalsize}

Le nombre annuel d'incidents matériels suit une binomiale négative, retenue contre la loi de
Poisson parce que la variance observée excède la moyenne~:
\[
N \sim \mathcal{BN}(r,p), \qquad \mathbb{E}[N]=\lambda, \qquad
\frac{\mathrm{Var}(N)}{\mathbb{E}[N]} \;=\; 1+\frac{\lambda}{r}.
\]
\lepiege{L'indice de dispersion \textbf{n'est pas invariant d'échelle}~: il dépend de $\lambda$,
donc il vaut autre chose sur des cellules firme-année, sur la série agrégée et au moteur. Trois
valeurs du projet sont trois \emph{échelles de lecture}, pas trois mesures contradictoires.}

\subsection*{3. La sévérité \quad\small\textcolor{navy}{statut~: estimée}\normalsize}

Au-delà du seuil $u$, les excès suivent une loi de Pareto généralisée, ce que le théorème de
Balkema et de Haan justifie asymptotiquement~:
\[
\mathbb{P}(X>u+y \mid X>u) \;=\; \Bigl(1+\xi\,\frac{y}{\sigma}\Bigr)^{-1/\xi},
\qquad y\ge 0 .
\]
D'où le quantile de \emph{sévérité d'un sinistre}, avec $\zeta_u=\mathbb{P}(X>u)$~:
\[
q_\alpha \;=\; u \;+\; \frac{\sigma}{\xi}
\left[\Bigl(\frac{1-\alpha}{\zeta_u}\Bigr)^{-\xi}-1\right].
\]
\cheznous{$u=20{,}03$~M\euro, $91$ excès, $\xi=0{,}5954$, $\sigma=57{,}97$.}
\lepiege{$\xi$ n'est pas qu'un paramètre de forme, c'est une \textbf{condition d'existence}~: la
moyenne existe si $\xi<1$, la variance si $\xi<1/2$. C'est cet argument, et non un argument de
degré, qui écarte deux variantes de modèle.}

\subsection*{4. La cascade \quad\small\textcolor{navy}{statut~: normalisation démontrée, amplitude posée}\normalsize}

On part d'une matrice d'influence brute $\mathrm{TRANS}$, asymétrique, et on la normalise à la
Leontief par l'\textbf{émission} maximale~:
\[
s_j=\sum_k \mathrm{TRANS}_{jk}, \qquad s_{\max}=\max_j s_j, \qquad
W \;=\; g\,\frac{\mathrm{TRANS}}{s_{\max}}, \qquad
e_j \;=\; g\,\frac{s_j}{s_{\max}} \;\le\; g .
\]
\noindent La convention est que $W_{jk}$ va de la \textbf{source} $j$ vers la \textbf{cible} $k$~:
la ligne émet, la colonne reçoit. La normalisation \emph{garantit} que le pilier le plus prolifique
engendre au plus $g$ descendants directs, donc $\rho(W)\le g<1$ et la sous-criticité est acquise
par construction, non supposée.

\cheznous{$s_{\max}=2{,}60$, $\rho(W)=0{,}506$, $R_0=0{,}054$, et le $g$ critique vaut $1{,}78$,
hors du domaine admissible.}
\lepiege{Diviser par la \emph{réception} maximale au lieu de l'émission donnait $e_{P_1}=1{,}02$
pour $g=0{,}9$~: la borne que la normalisation existe pour poser était franchie, et $g$ cessait de
désigner ce qu'il désigne.}

\subsection*{5. L'ensemble touché, et le coût d'un sinistre}

Le moteur tire une \textbf{marche auto-évitante}~: depuis le pilier courant $j$, la cascade
continue avec probabilité $e_j$, le successeur est tiré proportionnellement à
$\mathrm{TRANS}_{j\cdot}$, et retomber sur un pilier déjà touché l'éteint. On obtient l'ensemble
$S_a$ des piliers atteints depuis l'amorce $a$, et le coût du sinistre vaut
\[
C \;=\; \sum_{j\in S_a} c_j
\qquad\quad
\small\textcolor{accent}{\textbf{hypothèse de construction, non testée}}\normalsize
\]
\lepiege{C'est la plus lourde des hypothèses du mémoire. Relâchée par un exposant sur $|S_a|$, elle
déplace l'écart de capital de $9\,706$ à $20\,487$~M\euro. Ce qui survit à toute la plage~: le
signe et l'ordre. Ce qui ne se déclare pas~: le \emph{sens} de l'effet.}

\subsection*{6. La charge annuelle et le capital \quad\small\textcolor{navy}{statut~: définition}\normalsize}

\[
L \;=\; \sum_{i=1}^{N} C_i, \qquad\qquad
\mathrm{SCR} \;=\; \mathrm{VaR}_{99{,}5\,\%}(L)\quad\text{(non centrée)}.
\]
\lepiege{Le capital publié est le quantile \emph{entier}, sans retrancher $\mathbb E[L]$ : la
perte attendue d'une non-conformité n'est couverte par aucune provision, donc elle n'en est pas
déduite. Centré, le chiffre de tête tomberait à environ $5\,430\to16\,320$, un facteur de l'ordre
de $3{,}0$ (calcul hors mémoire).}
\lepiege{$\mathrm{VaR}$ et $\mathrm{TVaR}$ sont réservées à la charge \textbf{annuelle agrégée}
$L$~; $q_\alpha$ du bloc~3 est un quantile de \emph{sévérité d'un sinistre}. Les confondre fait
lire une différence d'échelle comme une contradiction, et c'est arrivé en séance. Le niveau
$99{,}5\,\%$ n'a de sens réglementaire que sur un horizon annuel~: une perte isolée n'en porte
aucun.}

\subsection*{7. Les états de conformité \quad\small\textcolor{navy}{statut~: structure posée, charge posée}\normalsize}

La conformité de chaque pilier dérive d'une latente à facteur commun, transposée du risque de
crédit~:
\[
C^\star_j \;=\; \gamma\,\Theta \;+\; \sqrt{1-\gamma^{2}}\;\varepsilon_j,
\qquad \Theta,\varepsilon_j \ \text{indépendantes} \ \mathcal{N}(0,1),
\]
\[
\mathrm{corr}(C^\star_j, C^\star_k) \;=\; \gamma^{2} \;=\; 0{,}4624
\quad\text{pour } \gamma=0{,}68 .
\]
\noindent Des seuils ordonnés $K_{\text{bas}}<K_{\text{haut}}$ découpent les états non conforme,
partiel et conforme.
\lepiege{La corrélation de $0{,}462$ \textbf{n'est pas posée}~: elle est le \emph{carré} de la
charge de facteur commun. Si l'on demande d'où elle sort, c'est la réponse.}

\subsection*{8. L'identification partielle \quad\small\textcolor{navy}{statut~: démontré}\normalsize}

Toute matrice se décompose de façon unique en une part symétrique et une part antisymétrique,
orthogonales au sens de Frobenius~:
\[
W = S + A, \qquad S=\tfrac12\,(W+W^{\!\top}), \qquad A=\tfrac12\,(W-W^{\!\top}).
\]
\noindent La donnée identifie $S$, la co-occurrence~; elle n'identifie pas $A$, la direction. On
ne pose donc pas $A$, on la \textbf{borne}, par un degré d'ignorance $t$~:
\[
\mathcal{W}(t) \;=\; \bigl\{\,W=S+A \ :\ |A_{jk}| \le t\,S_{jk}\,\bigr\},
\qquad
\bigl[\underline{\mathrm{SCR}},\ \overline{\mathrm{SCR}}\bigr]
= \bigl[\min_{\mathcal{W}(t)}, \ \max_{\mathcal{W}(t)}\bigr].
\]
\cheznous{Cinq piliers donnent dix paires, donc $2^{10}=1\,024$ sommets, énumérés
exhaustivement. Bornes $[6\,858\,;8\,697]$~M\euro, largeur $1\,839$.}
\laphrase{L'ignorance se paye linéairement en $t$, et son coût maximal est connu d'avance. C'est
plus favorable qu'un intervalle de confiance ordinaire, qui ne dit pas ce qu'il coûterait de
savoir.}

\subsection*{9. Les quatre canaux et leur interaction \quad\small\textcolor{navy}{statut~: mesuré}\normalsize}

La non-conformité déplace quatre paramètres, et l'on énumère les $2^4=16$ configurations. Pour un
sous-ensemble de canaux $S$, l'effet propre se lit par une décomposition de Möbius~:
\[
m(S) \;=\; \sum_{T\subseteq S} (-1)^{|S|-|T|}\,\Delta(T),
\qquad
\Delta(\text{tous}) \;=\; \sum_{S} m(S) .
\]
\cheznous{$\lambda$ de $21{,}6$ à $53{,}6$, $\zeta_u$ de $0{,}128$ à $0{,}181$, $g$ de $0{,}45$ à
$0{,}90$, accumulation tiers de $0$ à $0{,}68$. Capital de $6\,049$ à $20\,188$~M\euro, facteur
$3{,}34$, écart $14\,139$. Les canaux isolés somment à $9\,138$~: il manque $5\,001$~M\euro{}
d'interaction, soit $+35\,\%$.}
\lepiege{\textbf{Ne jamais additionner les quatre effets isolés} pour chiffrer une remédiation
partielle. Et ne pas confondre cette interaction, entre \emph{canaux}, avec celle entre
\emph{piliers}~: deux partitions du même écart, et la seconde a un signe non résolu.}

\subsection*{10. Le pont entre les deux échelles}

Question classique~: pourquoi un quantile de sévérité de quelques centaines de millions
cohabite-t-il avec un capital de plusieurs milliers de millions~? Parce qu'à $\lambda\simeq21{,}6$
sinistres par an, atteindre le quantile \emph{annuel} à $99{,}5\,\%$ demande un quantile \emph{par
sinistre} de l'ordre de $99{,}98\,\%$, bien plus haut dans la queue. L'approximation par perte
unique en donne la forme, la survie de la somme valant asymptotiquement $\mathbb E[K]$ fois celle
d'une perte isolée :
\[
\mathrm{VaR}_\alpha(L)\;\approx\;q_{\,1-(1-\alpha)/\mathbb E[K]}(X),
\qquad\text{d'où un quantile qui croît comme } (\lambda\,p_u)^{\xi}.
\]
\cheznous{Pour le canal de fréquence, cette forme prédit un facteur
$(53{,}62/21{,}56)^{0{,}5954}=1{,}720$ entre états, contre $1{,}715$ simulé ; pour la détection,
$1{,}228$ contre $1{,}239$. Au second ordre, elle retrouve le quantile publié à moins de $10\,\%$.}
\lepiege{Le quantile de sévérité ($663$~M\euro) n'est pas un capital, et le rapport entre les deux
objets n'est pas fixe : à l'échelle d'entité, le capital ($169$~M\euro) passe même sous le quantile
unitaire. Les comparer ne conclut jamais, dans un sens comme dans l'autre.}

% =====================================================================
\newpage
\section*{Annexe~: les nombres à connaître par c\oe ur}
\addcontentsline{toc}{section}{Annexe : les nombres à connaître}

% TABLE RESSERREE LE 23 SEPTEMBRE 2026. A l'interligne de 1,25 elle debordait sa page de
% 67,8 pt, un tabular ne sachant pas se couper : le bas de la table sortait du cadre sans
% que rien ne le signale a la lecture du source. Defaut anterieur, trouve a la compilation.
% PASSEE EN \footnotesize LE 30 SEPTEMBRE 2026 : en \small, avec les lignes ajoutees, l'encadre
% final tombait seul sur une page.
\begin{center}\footnotesize
\renewcommand{\arraystretch}{1.02}
\begin{tabular}{@{}l l l@{}}
\toprule
\textbf{Nombre} & \textbf{Ce que c'est} & \textbf{Chapitre} \\
\midrule
\multicolumn{3}{@{}l}{\textcolor{navy}{\textbf{Le chiffre de tête, à savoir sans hésiter}}} \\
$6\,049 \to 20\,188$~M\euro & besoin de capital, conforme vers non conforme, au secteur & 10 \\
$3{,}34$ & le facteur entre les deux états & 10 \\
$14\,139$~M\euro & l'écart, dont $9\,138$ portés par les canaux isolés & 10 \\
$5\,001$~M\euro\ ($35\,\%$) & la part portée par l'interaction des canaux & 10 \\
$10\,377$ / $11\,413$ & la fréquence seule / les trois autres canaux réunis & 10 \\
$[6\,858\,;\,8\,697]$~M\euro & \textbf{la réponse à \og quel est votre SCR\fg{}} & 9 \\
$1\,839$~M\euro & largeur de bande = coût de l'ignorance, connu d'avance & 9 \\
\midrule
$582$ / $91$ & pertes cyber $\times$ finance / excès au-delà du seuil & 4, 5 \\
$u=20{,}03$~M\euro & le seuil des excès (percentile $84{,}4$) & 5 \\
$\xi=0{,}5954$ / $1{,}03$ & indice de queue OpRisk / PRC & 5 \\
$663$ / $1\,752$~M\euro & quantile et moyenne de queue d'\emph{un} sinistre, pas un capital & 5 \\
$[411\,;1\,037]$, $\times2{,}5$ & intervalle bootstrap de ce quantile & 5 \\
$p_u=0{,}1509$ & taux gelé ($91/582=0{,}1564$ ; $+2{,}4\,\%$ sur le quantile) & 11 \\
$341$ / $21{,}6$ / $0{,}092$ & les trois fréquences (facteur $3\,719$) & 4 \\
$1{,}19$ / $2{,}24$ / $9{,}20$ & surdispersion : cellules, série, moteur (posé) & 5 \\
$p\approx10^{-30}$ & rejet de Poisson contre la binomiale négative & 5 \\
$0{,}087$ $[0{,}026\,;0{,}148]$ & élasticité sévérité-taille, corrigée du seuil & 4 \\
\midrule
$g=0{,}45$ / $0{,}68$ / $0{,}90$ & gains de propagation posés, conforme / partiel / non conforme & 7 \\
$\rho(\mathrm{TRANS})=1{,}461$ & pourquoi il faut renormaliser & 6 \\
$\rho(W)=0{,}506$, $R_0=0{,}054$ & les deux rayons spectraux & 6 \\
$62{,}31\,\%$ / $31{,}15\,\%$ & sinistres multi-piliers, non conforme / conforme & 11 \\
$R^2=0{,}9945$ & le capital, presque additif en piliers & 7 \\
$K_4=1{,}03$ ; $\lambda_U=0{,}53$ & seuil de P4 ; dépendance de queue de Gumbel & B \\
\midrule
$119$ contre $128\pm27$ & placebo sur OpRisk, $z=-0{,}33$ & 8 \\
$z=+3{,}93\to+5{,}12$ & placebo sur $7$ puis $10$ rapports & 8 \\
$0{,}12\,\%$ / $1{,}19\,\%$ / $1$ & VERIS : champ, chronologie, intersection & 8 \\
$12$ / $0$ sur $20$ & liens $p_{jk}$ crédiblement faibles / forts & 8 \\
$5\,275$~M\euro & socle sans contagion & 9 \\
$777$~M\euro & ce qu'un panel biaisé déplacerait sans signal & 9 \\
\midrule
$8\,123$ / $169$~M\euro & besoin du secteur / de l'entité, majorant & 10 \\
$47$ contre $8{,}3$ & multiple de capital, entité contre secteur & 10 \\
$2{,}0$ à $26{,}4\,\%$ & charge rapportée au SCR publié, quatre SFCR & 10 \\
$-64\,\%$ ; $\kappa^\star=78\,\%$ & gain du traité ; seuil d'inefficacité & 10 \\
$4{,}75\,\%$ & coût du capital, directive (UE) 2025/2 & 10 \\
$657$ / $645$ / $1\,037$ & point, prédictive, robuste à $95\,\%$ & 11 \\
$\times8{,}93$ / $\times1{,}27$ / $\times9{,}21$ & paramètre, identification, famille de queue & 11 \\
$91{,}7\,\%$ / $75{,}0\,\%$ & backtest : binomiale négative / Poisson & 5 \\
$4{,}00\,\%$ par an & dérive d'échelle de la queue & 5 \\
$1\,811$ années & pour tester $99{,}5\,\%$ à $80\,\%$ de puissance & 5 \\
$66\,\%$ / $28{,}9\,\%$ & valeur d'information : six sens / meilleure paire & 12 \\
\midrule
\multicolumn{3}{@{}l}{\textcolor{navy}{\textbf{Robustesse du classement, à sortir si l'on attaque le corpus}}} \\
$43{,}6\,\%$ / $37{,}7\,\%$ & P1 / P4 premiers sur l'ensemble admissible & 10 \\
$1{,}5\,\%$ & écart de regret maximal entre les deux & 10 \\
$a=0{,}25$ / $a=0{,}50$ & le classement tient / bascule vers P4 & 8, 10 \\
$-24{,}8\,\%$ ($0{,}81\times$) & chute du niveau d'entité si P1 devient muette & 8 \\
$34\,\%$ & part de Shapley de P1 dans le surcoût & 10 \\
\bottomrule
\end{tabular}
\end{center}

\vspace{4mm}
\begin{cle}
\textbf{La phrase qui tient tout.} Si l'on ne devait retenir qu'une chose de ce document~:
\emph{le niveau de capital bouge dans une bande large ; ce qui résiste, c'est le signe et l'ordre
de grandeur de l'écart entre états face à un forfait qui ne bouge pas, le rapport au socle, et le
groupe de tête des piliers, P1 et P4 devant les trois autres}. L'ordre entre P1 et P4, lui, dépend
d'une direction que la donnée ne fixe pas, et le dire fait partie de la réponse.
\end{cle}

\end{document}
